% Generated mechanically from frozen input inputs/p10/ja/restricted.tex.
% Do not edit this generated file; edit the bound locale source instead.

% OK, start here.
%

\setcounter{chapter}{87}
\chapter{形式空間の代数化}



\phantomsection
\label{restricted-section-phantom}


\section{序論}
\label{restricted-section-introduction}

\noindent
本章の主要な目標は，Artin の dilatation（膨張）に関する定理を証明することである。
定理 \ref{restricted-theorem-dilatations} を参照せよ。収縮に関する結果は Artin's Axioms，
Section \ref{artin-section-contractions} で論じる。どちらの結果にも形式代数空間に
関する内容が必要となるので，本章の中ほどでは前章に引き続いて形式代数空間を論じる。
Formal Spaces，Section \ref{formal-spaces-section-introduction} を参照せよ。本章の
第一部は代数的な準備に充て，主として rig-\'etale 代数の代数化を扱う。

\medskip\noindent
$A$ を Noether 環，$I \subset A$ をイデアルとする。本章の第一部
（Sections \ref{restricted-section-two-categories} -- \ref{restricted-section-approximation-principal}）では，
Noether 環 $A$ 上位相的有限型な $I$-進完備代数 $B$ の圏を論じる。$A$ に
$I$-進位相を入れたとき，ある制限冪級数環内の（閉）イデアル $J$ によって
$B = A\{x_1, \ldots, x_n\}/J$ と書けることを示す。また，素朴な余接複体
$\NL_{B/A}^\wedge$ には適切な概念があることを示す。$I$ のある冪が
$\NL_{B/A}^\wedge$ を零化するとき，$B$ を $(A, I)$ 上の rig-\'etale 代数と呼ぶ。
rig-smooth 代数も同様に定義される。$A$ が G-環ならば，Popescu の定理を用いて，
$(A, I)$ 上の任意の rig-smooth 代数 $B$ は有限型 $A$-代数の完備化であることを示す。
非形式的には，$B$ を「代数化」できるという。しかし第一部の主結果は，$A$ が
G-環であると仮定しなくても，$(A, I)$ 上の任意の rig-\'etale 代数 $B$ を
代数化できるというものである。補題 \ref{restricted-lemma-approximate} を参照せよ。この種の
代数化に関する文献への案内は，Remark \ref{restricted-remark-discussion} にある。第一部の
一般的な参考文献は \cite{EGA}，\cite{Abbes}，\cite{Fujiwara-Kato} である。

\medskip\noindent
本章の第二部（Sections \ref{restricted-section-finite-type-red} --
\ref{restricted-section-formal-modifications}）では，形式代数空間の射の諸類型を適切な一般性の
もとで論じる（主として局所 Noether 形式代数空間を扱う）。そのうち最も興味深いのは，
最後の節に現れる「形式修正」の概念である。我々の定義が \cite{ArtinII} における
Artin の定義と一致することを注意深く確認する。

\medskip\noindent
最後に，本章の第三部かつ最終部（Sections \ref{restricted-section-completion-and-morphisms} --
\ref{restricted-section-modifications}）では主定理を証明し，いくつかの応用を与える。実際には，
より強い結果である定理 \ref{restricted-theorem-dilatations-general} から Artin の定理を導く。
この定理をきわめて粗く述べると，$f : \mathfrak X \to \mathfrak X'$ が
rig-\'etale 射で，$\mathfrak X'$ が局所 Noether 代数空間の形式完備化であるならば，
$\mathfrak X$ もそうである，というものである。Artin の研究では，射 $f$ は固有かつ
rig-全射であると仮定されている。

\section{二つの圏}
\label{restricted-section-two-categories}

\noindent
$A$ を環，$I \subset A$ をイデアルとする。本節では ${}^\wedge$ は
$I$-進完備化を表すものとする。$A_n = A/I^n$ と置くと，$A$ の $I$-進完備化は
$A^\wedge = \lim A_n$ である。$\mathcal{C}$ を次の圏とする：
\begin{equation}
\label{restricted-equation-C}
\mathcal{C} =
\left\{
\begin{matrix}
\text{逆系 }\ldots \to B_3 \to B_2 \to B_1 \\
\text{ここで }B_n\text{ は有限型 }A_n\text{-代数であり，}\\
B_{n + 1} \to B_n\text{ は }A_{n + 1}\text{-代数写像であり，}\\
\text{しかも次を誘導する：}B_{n + 1}/I^nB_{n + 1} \cong B_n
\end{matrix}
\right\}.
\end{equation}
$\mathcal{C}$ の射は準同型の系によって与えられる。$\mathcal{C}'$ を次の圏とする：
\begin{equation}
\label{restricted-equation-C-prime}
\mathcal{C}' =
\left\{
\begin{matrix}
A\text{-代数 }B\text{ であって }I\text{-進完備であり，}\\
\text{しかも }B/IB\text{ は }A/I\text{ 上有限型である}
\end{matrix}
\right\}.
\end{equation}
$\mathcal{C}'$ の射は $A$-代数写像である。函手
\begin{equation}
\label{restricted-equation-from-complete-to-systems}
\mathcal{C}' \longrightarrow \mathcal{C},\quad
B \longmapsto (B/I^nB)
\end{equation}
が存在する。実際，$B/IB$ は $A/I$ 上有限型なので，環写像
$A_n = A/I^n \to B/I^nB$ は Algebra, Lemma
\ref{algebra-lemma-finite-type-mod-nilpotent} により有限型である。

\begin{lemma}
\label{restricted-lemma-topologically-finite-type}
$A$ を環，$I \subset A$ を有限生成イデアルとする。函手
$$
\mathcal{C} \longrightarrow \mathcal{C}',\quad
(B_n) \longmapsto B = \lim B_n
$$
は (\ref{restricted-equation-from-complete-to-systems}) の準逆である。完備化
$A[x_1, \ldots, x_r]^\wedge$ は $\mathcal{C}'$ に属し，$\mathcal{C}'$ の任意の対象は，
あるイデアル $J \subset A[x_1, \ldots, x_r]^\wedge$ によって
$$
B = A[x_1, \ldots, x_r]^\wedge / J
$$
と書ける。
\end{lemma}

\begin{proof}
$(B_n)$ を $\mathcal{C}$ の対象とする。Algebra, Lemma
\ref{algebra-lemma-limit-complete} により，$B = \lim B_n$ は $I$-進完備であり，
$B/I^nB = B_n$ である。したがって $B$ は $\mathcal{C}'$ の対象であり，商をとることで
対象 $(B_n)$ を復元できる。逆に，$B$ が $\mathcal{C}'$ の対象ならば，仮定により
$B = \lim B/I^nB$ である。ゆえに $B \mapsto (B/I^nB)$ は補題の函手の準逆である。

\medskip\noindent
$A[x_1, \ldots, x_r]^\wedge = \lim A_n[x_1, \ldots, x_r]$ なので，補題の最初の主張により，
これは $\mathcal{C}'$ の対象である。最後に，$B$ を $\mathcal{C}'$ の対象とする。
$b_1, \ldots, b_r \in B$ を，それらの $B/IB$ における像が $B/IB$ を
$A/I$-代数として生成するように選ぶ。$B$ は $I$-進完備なので，$A$-代数写像
$A[x_1, \ldots, x_r] \to B$，$x_i \mapsto b_i$ は $A$-代数写像
$A[x_1, \ldots, x_r]^\wedge \to B$ に延長される。証明を完了するには，この写像が
全射であることを示せばよい。これは Algebra, Lemma
\ref{algebra-lemma-completion-generalities} から従う。実際，写像
$A[x_1, \ldots, x_r] \to B$ は $I$ を法として全射であり，$B = B^\wedge$ である。
\end{proof}


\noindent
$A$ が Noether でない場合，$\mathcal{C}'$ の対象の商が $\mathcal{C}'$ の対象になるとは
限らないことを注意しておく。Examples, Lemma
\ref{examples-lemma-noncomplete-quotient} を参照せよ。次に，$A$ が Noether ならば
この現象は起こらないことを示す。

\begin{lemma}
\label{restricted-lemma-topologically-finite-type-Noetherian}
\begin{reference}
\cite[Proposition 7.5.5]{EGA1}
\end{reference}
$A$ を Noether 環，$I \subset A$ をイデアルとする。このとき，
\begin{enumerate}
\item 圏 $\mathcal{C}'$（\ref{restricted-equation-C-prime}）の任意の対象は Noether である。
\item $B \in \Ob(\mathcal{C}')$ かつ $J \subset B$ がイデアルならば，$B/J$ は
$\mathcal{C}'$ の対象である。
\item 有限型 $A$-代数 $C$ に対し，$I$-進完備化 $C^\wedge$ は $\mathcal{C}'$ に属する。
\item 特に，完備化 $A[x_1, \ldots, x_r]^\wedge$ は $\mathcal{C}'$ に属する。
\end{enumerate}
\end{lemma}

\begin{proof}
(4) は Algebra, Lemma \ref{algebra-lemma-completion-Noetherian-Noetherian} から従う。
実際，$A[x_1, \ldots, x_r]$ は Noether である（Algebra, Lemma
\ref{algebra-lemma-Noetherian-permanence}）。(1) を示すには，補題
\ref{restricted-lemma-topologically-finite-type} により，いま証明した多項式環の完備化の場合へ
帰着できる。(2) は，任意の有限 $B$-加群が $IB$-進完備であることを述べる
Algebra, Lemma \ref{algebra-lemma-completion-tensor} から従う。(3) も (4) と同様に従う。
\end{proof}

\begin{remark}[基底変換]
\label{restricted-remark-base-change}
$\varphi : A_1 \to A_2$ を環写像，$I_i \subset A_i$ をイデアルとし，ある $c \geq 1$ に
対して $\varphi(I_1^c) \subset I_2$ であるとする。すると，すべての $n \geq 1$ に対して
環写像 $A_{1, cn} = A_1/I_1^{cn} \to A_2/I_2^n = A_{2, n}$ が誘導される。
$\mathcal{C}_i$ を $(A_i, I_i)$ に対する圏 (\ref{restricted-equation-C}) とする。基底変換函手
\begin{equation}
\label{restricted-equation-base-change-systems}
\mathcal{C}_1 \longrightarrow \mathcal{C}_2,\quad
(B_n) \longmapsto (B_{cn} \otimes_{A_{1, cn}} A_{2, n})
\end{equation}
が存在する。$\mathcal{C}_i'$ を $(A_i, I_i)$ に対する圏
(\ref{restricted-equation-C-prime}) とする。$I_2$ が有限生成ならば，基底変換函手
\begin{equation}
\label{restricted-equation-base-change-complete}
\mathcal{C}_1' \longrightarrow \mathcal{C}_2',\quad
B \longmapsto (B \otimes_{A_1} A_2)^\wedge
\end{equation}
が存在する。実際，この場合，完備化は完備である（Algebra, Lemma
\ref{algebra-lemma-hathat-finitely-generated}）。$I_1$ と $I_2$ がともに有限生成ならば，
二つの基底変換函手は函手 (\ref{restricted-equation-from-complete-to-systems}) を介して一致する。
補題 \ref{restricted-lemma-topologically-finite-type} により，これらの函手は圏同値である。
\end{remark}

\begin{remark}[閉埋め込みによる基底変換]
\label{restricted-remark-take-bar}
$A$ を Noether 環，$I \subset A$ をイデアルとする。$\mathfrak a \subset A$ を
イデアルとし，$\bar A = A/\mathfrak a$ と書く。$\bar I \subset \bar A$ をイデアルとし，
ある $c, d \geq 1$ に対して $I^c \bar A \subset \bar I$ かつ
$\bar I^d \subset I\bar A$ であるとする。この場合，$(A, I)$ から
$(\bar A, \bar I)$ への基底変換函手 (\ref{restricted-equation-base-change-complete}) は
$B \mapsto \bar B = B/\mathfrak aB$ によって与えられる。すなわち，
\begin{equation}
\label{restricted-equation-base-change-to-closed}
\bar B = (B \otimes_A \bar A)^\wedge = (B/\mathfrak a B)^\wedge =
B/\mathfrak a B.
\end{equation}
最後の等号が成り立つのは，任意の有限 $B$-加群が Algebra, Lemma
\ref{algebra-lemma-completion-tensor} により $I$-進完備であり，さらに
$\mathfrak a$ によって零化される場合には Algebra, Lemma
\ref{algebra-lemma-change-ideal-completion} により $\bar I$-進完備でもあるからである。
\end{remark}

\section{素朴な余接複体}
\label{restricted-section-naive-cotangent-complex}

\noindent
$A$ を Noether 環，$I \subset A$ をイデアルとする。$B$ を $I$-進完備な
$A$-代数で，$A/I \to B/IB$ が有限型，すなわち
(\ref{restricted-equation-C-prime}) の対象であるものとする。補題
\ref{restricted-lemma-topologically-finite-type-Noetherian} により，ある有限生成イデアル
$J$ に対して
$$
B = A[x_1, \ldots, x_r]^\wedge / J
$$
と書ける。上のような表示を一つ選び，この状況における
{\it 素朴な余接複体}を次の式で定義する：
\begin{equation}
\label{restricted-equation-NL}
\NL_{B/A}^\wedge = (J/J^2 \longrightarrow \bigoplus B\text{d}x_i).
\end{equation}
各項は次数 $-1$ と $0$ に置かれ，写像は $g \in J$ の剰余類を微分
$\text{d}g = \sum (\partial g/\partial x_i) \text{d}x_i$ に送る。ここで偏微分は，
$g$ を冪級数とみなしてとる。次の補題は，$\NL_{B/A}^\wedge$ がホモトピーを除いて
良定義であることを示す。

\begin{lemma}
\label{restricted-lemma-NL-up-to-homotopy}
$A$ を Noether 環，$I \subset A$ をイデアルとする。$B$ を
(\ref{restricted-equation-C-prime}) の対象とする。素朴な余接複体
$\NL_{B/A}^\wedge$ は $K(B)$ において良定義である。
\end{lemma}

\begin{proof}
この補題の意味は，第二の表示 $B = A[y_1, \ldots, y_s]^\wedge / K$ が与えられたとき，
$B$-加群の複体
$$
(J/J^2 \to B\text{d}x_i)
\quad\text{および}\quad
(K/K^2 \to \bigoplus B\text{d}y_j)
$$
がホモトピー同値になるということである。これは Algebra, Lemma
\ref{algebra-lemma-NL-homotopy} の証明とまったく同様に示せる。

\medskip\noindent
ステップ 1．$B$ における $x_i$ の像へ写る
$g_i(y_1, \ldots, y_s) \in A[y_1, \ldots, y_s]^\wedge$ を選ぶと，与えられた
$B$ への全射と両立する（一意な）連続 $A$-代数準同型
$$
A[x_1, \ldots, x_r]^\wedge \to A[y_1, \ldots, y_s]^\wedge,\quad
x_i \mapsto g_i(y_1, \ldots, y_s)
$$
を得る。このような写像を表示の射と呼ぶ。この写像は $J$ を $K$ に写すので，
$B$-加群写像 $J/J^2 \to K/K^2$ を誘導する。$\text{d}x_i$ を
$\sum (\partial g_i/\partial y_j)\text{d}y_j$ に送ることにより，複体の写像
$$
(J/J^2 \to \bigoplus B\text{d}x_i)
\longrightarrow
(K/K^2 \to \bigoplus B\text{d}y_j)
$$
を得る。もちろん，二つの表示の役割を交換して同じ構成を行えば，逆向きの複体の写像を得る。

\medskip\noindent
% P10-E0383
ステップ 2．上の構成は表示の射の合成と両立する。したがって証明を完了するには，
$B$ における $x_i$ の像へ写る
$g_i(x_1, \ldots, x_r) \in A[x_1, \ldots, x_n]^\wedge$ が与えられたとき，誘導される複体の写像
$$
(J/J^2 \to \bigoplus B\text{d}x_i)
\longrightarrow
(J/J^2 \to \bigoplus B\text{d}x_i)
$$
が恒等写像にホモトピックであることを示せば十分である。そのためには，
$\text{d}x_i \mapsto g_i(x_1, \ldots, x_n) - x_i$ という規則で定められる写像
$h : \bigoplus B \text{d}x_i \to J/J^2$ を考え，計算すればよい。
\end{proof}

\begin{lemma}
\label{restricted-lemma-NL-is-completion}
$A$ を Noether 環，$I \subset A$ をイデアルとする。$A \to B$ を有限型環写像とし，
表示 $\alpha : A[x_1, \ldots, x_n] \to B$ を選ぶ。このとき複体として
$\NL_{B^\wedge/A}^\wedge = \lim \NL(\alpha) \otimes_B B^\wedge$ であり，
$D(B^\wedge)$ において
$\NL_{B^\wedge/A}^\wedge = \NL_{B/A} \otimes_B^\mathbf{L} B^\wedge$ である。
\end{lemma}

\begin{proof}
この主張には意味がある。実際，$B^\wedge$ は
(\ref{restricted-equation-C-prime}) の対象であることが，補題
\ref{restricted-lemma-topologically-finite-type-Noetherian} から分かる。
$J = \Ker(\alpha)$ と置く。$A[x_1, \ldots, x_n]$ 上の有限加群について
$I$-進完備化をとる函手は完全であり，函手
$M \mapsto M \otimes_{A[x_1, \ldots, x_n]} A[x_1, \ldots, x_n]^\wedge$
と一致する。Algebra, Lemmas \ref{algebra-lemma-completion-tensor} および
\ref{algebra-lemma-completion-flat} を参照せよ。さらに，環写像
$A[x_1, \ldots, x_n] \to A[x_1, \ldots, x_n]^\wedge$ と
$B \to B^\wedge$ は平坦である。したがって
$B^\wedge = A[x_1, \ldots, x_n]^\wedge / J^\wedge$ であり，
$$
(J/J^2) \otimes_B B^\wedge = (J/J^2)^\wedge = J^\wedge/(J^\wedge)^2
$$
Algebra, Section \ref{algebra-section-netherlander} にあるように
$\NL(\alpha) = (J/J^2 \to \bigoplus B\text{d}x_i)$ なので，複体
$\NL_{B^\wedge/A}^\wedge$ は $\NL(\alpha) \otimes_B B^\wedge$ に等しい。
最後の主張は，$\NL_{B/A}$ が $\NL(\alpha)$ とホモトピー同値であり，かつ環写像
$B \to B^\wedge$ が平坦であることから従う（したがって
$B \to B^\wedge$ に沿う導来基底変換は通常の基底変換にすぎない）。
\end{proof}

\begin{lemma}
\label{restricted-lemma-NL-is-limit}
$A$ を Noether 環，$I \subset A$ をイデアルとする。$B$ を
(\ref{restricted-equation-C-prime}) の対象とする。このとき，
\begin{enumerate}
\item $D(B)$ の pro-対象
$\{\NL_{B/A}^\wedge \otimes_B B/I^nB\}$ と $\{\NL_{B_n/A_n}\}$ は厳密に同型である
（説明は証明を参照せよ）。
\item $D(B)$ において $\NL_{B/A}^\wedge = R\lim \NL_{B_n/A_n}$ である。
\end{enumerate}
ここで $B_n$ と $A_n$ は Section \ref{restricted-section-two-categories} で定めたものである。
\end{lemma}

\begin{proof}
この主張の意味は次のとおりである。各 $n$ に対して，良定義な $B_n$-加群の複体
$\NL_{B_n/A_n}$ と遷移写像
$\NL_{B_{n + 1}/A_{n + 1}} \to \NL_{B_n/A_n}$ がある。
Algebra, Section \ref{algebra-section-netherlander} を参照せよ。したがって
$$
\ldots \to \NL_{B_3/A_3} \to \NL_{B_2/A_2} \to \NL_{B_1/A_1}
$$
を $B$-加群の複体の逆系，したがって特に $D(B)$ の逆系とみなせる。さらに，
$R\lim \NL_{B_n/A_n}$ はこの逆系のホモトピー極限である。Derived Categories,
Section \ref{derived-section-derived-limit} を参照せよ。

\medskip\noindent
表示 $B = A[x_1, \ldots, x_r]^\wedge / J$ を選ぶ。これにより表示
$$
B_n = B/I^nB = A_n[x_1, \ldots, x_r]/J_n
$$
が定まり，ここで
$$
J_n = JA_n[x_1, \ldots, x_r] =
J/(J \cap I^nA[x_1, \ldots, x_r]^\wedge)
$$
二項複体 $J_n/J_n^2 \longrightarrow \bigoplus B_n \text{d}x_i$ は
$\NL_{B_n/A_n}$ を表す。Algebra, Section
\ref{algebra-section-netherlander} を参照せよ。Noether 環
$A[x_1, \ldots, x_r]^\wedge$ における Artin--Rees
（Algebra, Lemma \ref{algebra-lemma-Artin-Rees}）と補題
\ref{restricted-lemma-topologically-finite-type-Noetherian} により，すべての $n \geq c$ に対して
標準的な全射
$$
J/I^nJ \to J_n \to J/I^{n - c}J \to J_{n - c},\quad n \geq c
$$
が存在するような $c \geq 0$ をとれる。少し考えれば，これらの写像が微分と両立し，
複体の写像
$$
\NL_{B/A}^\wedge \otimes_B B/I^nB \to
\NL_{B_n/A_n} \to
\NL_{B/A}^\wedge \otimes_B B/I^{n - c}B \to
\NL_{B_{n - c}/A_{n - c}}
$$
を得ることが分かる。これらは逆系
$\{\NL_{B/A}^\wedge \otimes_B B/I^nB\}$ と $\{\NL_{B_n/A_n}\}$ の遷移写像と
両立する。これで補題の (1) が証明された。

\medskip\noindent
(1) と，pro-同型な系は同じ $R\lim$ をもつことから，(2) を示すには
$\NL_{B/A}^\wedge$ が
$R\lim \NL_{B/A}^\wedge \otimes_B B/I^nB$ に等しいことを示せば十分である。
ところが，$\NL_{B/A}^\wedge$ は有限 $B$-加群からなる二項複体 $M^\bullet$ であり，
これらの加群は，例えば Algebra, Lemma \ref{algebra-lemma-completion-tensor} により
$I$-進完備である。ゆえに
$M^\bullet = \lim M^\bullet/I^nM^\bullet =
R\lim M^\bullet/I^n M^\bullet$ である。More on Algebra, Lemma
\ref{more-algebra-lemma-compute-Rlim-modules} および Remark
\ref{more-algebra-remark-how-unique} を参照せよ。
\end{proof}

\begin{lemma}
\label{restricted-lemma-NL-base-change}
$(A_1, I_1) \to (A_2, I_2)$ を Remark \ref{restricted-remark-base-change} のものとし，
$A_1$ と $A_2$ は Noether であるとする。$B_1$ を $(A_1, I_1)$ に対する
(\ref{restricted-equation-C-prime}) の対象とし，$B_2$ を $B_1$ の基底変換とする。このとき標準写像
$$
\NL_{B_1/A_1} \otimes_{B_2} B_1 \to \NL_{B_2/A_2}
$$
% P10-E0384
が存在し，これは $H^0$ 上で同型を，$H^{-1}$ 上で全射を誘導する。
\end{lemma}

\begin{proof}
表示 $B_1 = A_1[x_1, \ldots, x_r]^\wedge/J_1$ を選ぶ。
$A_2/I_2^n[x_1, \ldots, x_r] =
A_1/I_1^{cn}[x_1, \ldots, x_r] \otimes_{A_1/I_1^{cn}} A_2/I_2^n$
なので，
$$
A_2[x_1, \ldots, x_r]^\wedge =
(A_1[x_1, \ldots, x_r]^\wedge \otimes_{A_1} A_2)^\wedge
$$
である。ここでは両辺で $I_2$-進完備化を用いる
（ただし当然，$A_1[x_1, \ldots, x_r]^\wedge$ については $I_1$-進完備化を用いる）。
$J_2 = J_1 A_2[x_1, \ldots, x_r]^\wedge$ と置く。同様に議論すれば，
$B_2$ の $A_2$ 上の表示
\begin{align*}
B_2
& =
(B_1 \otimes_{A_1} A_2)^\wedge \\
& =
\lim \frac{A_1/I_1^{cn}[x_1, \ldots, x_r]}{J_1(A_1/I_1^{cn}[x_1, \ldots, x_r])}
\otimes_{A_1/I_1^{cn}} A_2/I_2^n \\
& =
\lim \frac{A_2/I_2^n[x_1, \ldots, x_r]}{J_2(A_2/I_2^n[x_1, \ldots, x_r])} \\
& =
A_2[x_1, \ldots, x_r]^\wedge/J_2
\end{align*}
を得る。その結果，次の可換図式を得る：
$$
\xymatrix{
\NL^\wedge_{B_1/A_1} : \ar[d] &
J_1/J_1^2 \ar[r]_-{\text{d}} \ar[d] & \bigoplus B_1\text{d}x_i \ar[d] \\
\NL^\wedge_{B_2/A_2} : &
J_2/J_2^2 \ar[r]^-{\text{d}} & \bigoplus B_2\text{d}x_i
}
$$
誘導される射 $J_1/J_1^2 \otimes_{B_1} B_2 \to J_2/J_2^2$ は全射である。
実際，$J_2$ は $J_1$ の像によって生成される。これにより補題に表示した射が定まる。
この射がホモトピーを除いて良定義であること
（すなわち，表示の選択にホモトピーを除いて依存しないこと）の証明は省略する。
誘導されるコホモロジー加群上の写像に関する主張は，上の議論から容易に従う
（詳細は省略する）。
\end{proof}

\begin{lemma}
\label{restricted-lemma-exact-sequence-NL}
$A$ を Noether 環，$I \subset A$ をイデアルとする。$B \to C$ を
(\ref{restricted-equation-C-prime}) の射とする。このとき正完全列
$$
\xymatrix{
C \otimes_B H^0(\NL_{B/A}^\wedge) \ar[r] &
H^0(\NL_{C/A}^\wedge) \ar[r] &
H^0(\NL_{C/B}^\wedge) \ar[r] & 0 \\
H^{-1}(\NL_{B/A}^\wedge \otimes_B C) \ar[r] &
H^{-1}(\NL_{C/A}^\wedge) \ar[r] &
H^{-1}(\NL_{C/B}^\wedge) \ar[llu]
}
$$
が存在する。説明は証明を参照せよ。
\end{lemma}

\begin{proof}
テンソル積 $\NL_{B/A}^\wedge \otimes_B C$ をとることには意味がある。実際，補題
\ref{restricted-lemma-NL-up-to-homotopy} により $\NL_{B/A}^\wedge$ はホモトピーを除いて
良定義である。また，$(B, IB)$ は $B$ が Noether 環である組であり
（補題 \ref{restricted-lemma-topologically-finite-type-Noetherian}），$C$ は対応する圏
(\ref{restricted-equation-C-prime}) に属する。したがって，$6$ 項完全列のすべての項が良定義である。
% P10-E0385

\medskip\noindent
表示 $B = A[x_1, \ldots, x_r]^\wedge/J$ と表示
$C = B[y_1, \ldots, y_s]^\wedge/J'$ を選ぶ。これらの表示を組み合わせると，表示
$$
C = A[x_1, \ldots, x_r, y_1, \ldots, y_s]^\wedge/K
$$
を得る。すると，次の可換図式を得ることを読者は確認できる：
% P10-E0386
$$
\xymatrix{
0 \ar[r] &
\bigoplus C \text{d}x_i \ar[r] &
\bigoplus C \text{d}x_i \oplus \bigoplus C \text{d}y_j \ar[r] &
\bigoplus C \text{d}y_j \ar[r] &
0 \\
&
J/J^2 \otimes_B C \ar[r] \ar[u] &
K/K^2 \ar[r] \ar[u] &
J'/(J')^2 \ar[r] \ar[u] &
0
}
$$
各行は完全である。左側の垂直射は，$\NL_{B/A}^\wedge$ を定義する射と
$\text{id}_C$ とのテンソル積であることに注意せよ。補題は蛇の補題
（Algebra, Lemma \ref{algebra-lemma-snake}）を適用すれば従う。
\end{proof}

\begin{lemma}
\label{restricted-lemma-transitive-lci-at-end}
補題 \ref{restricted-lemma-exact-sequence-NL} と同じ仮定のもとで，すべての $n$ に対して
$B/I^nB \to C/I^nC$ が局所完全交叉準同型であると仮定する。このとき
$H^{-1}(\NL_{B/A}^\wedge \otimes_B C) \to H^{-1}(\NL_{C/A}^\wedge)$ は単射である。
\end{lemma}

\begin{proof}
各 $n \geq 1$ に対し，$A_n = A/I^n$，$B_n = B/I^nB$，
$C_n = C/I^nC$ と置く。次を得る：
\begin{align*}
H^{-1}(\NL_{B/A}^\wedge \otimes_B C)
& =
\lim H^{-1}(\NL_{B/A}^\wedge \otimes_B C_n) \\
& =
\lim H^{-1}(\NL_{B/A}^\wedge \otimes_B B_n \otimes_{B_n} C_n) \\
& =
\lim H^{-1}(\NL_{B_n/A_n} \otimes_{B_n} C_n).
\end{align*}
最初の等号は More on Algebra, Lemma
\ref{more-algebra-lemma-consequence-Artin-Rees} と，
$H^{-1}(\NL_{B/A}^\wedge \otimes_B C)$ が有限 $C$-加群であり，したがって例えば
Algebra, Lemma \ref{algebra-lemma-completion-tensor} により $I$-進完備であることから従う。
第二の等号は明らかである。第三の等号は補題 \ref{restricted-lemma-NL-is-limit} から従う。
写像 $H^{-1}(\NL_{B_n/A_n} \otimes_{B_n} C_n) \to
H^{-1}(\NL_{C_n/A_n})$ は More on Algebra, Lemma
\ref{more-algebra-lemma-transitive-lci-at-end} により単射である。上と同様に
$H^{-1}(\NL_{C/A}^\wedge) = \lim H^{-1}(\NL_{C_n/A_n})$ でもあるから，証明は完了する。
\end{proof}

\section{Rig-smooth 代数}
\label{restricted-section-rig-smooth}

\noindent
次の定義の動機については，More on Algebra, Remark
\ref{more-algebra-remark-smoothness-ext-1-zero} を参照されたい。

\begin{definition}
\label{restricted-definition-rig-smooth-homomorphism}
$A$ を Noether 環，$I \subset A$ をイデアルとする。$B$ を
(\ref{restricted-equation-C-prime}) の対象とする。ある整数 $c \geq 0$ が存在して，すべての
$B$-加群 $N$ に対し $I^c$ が $\Ext^1_B(\NL_{B/A}^\wedge, N)$ を零化するとき，
$B$ は {\it $(A, I)$ 上 rig-smooth} であるという。
\end{definition}

\noindent
この条件の意味を具体的に記述しよう。

\begin{lemma}
\label{restricted-lemma-equivalent-with-artin-smooth}
$A$ を Noether 環，$I \subset A$ をイデアルとする。$B$ を
(\ref{restricted-equation-C-prime}) の対象とする。
$B = A[x_1, \ldots, x_r]^\wedge/J$ と書き
（補題 \ref{restricted-lemma-topologically-finite-type-Noetherian}），その素朴な余接複体を
$\NL_{B/A}^\wedge = (J/J^2 \to \bigoplus B\text{d}x_i)$
（\ref{restricted-equation-NL}）とする。次は同値である：
\begin{enumerate}
\item $B$ は $(A, I)$ 上 rig-smooth である。
\item $D(B)$ の対象 $\NL_{B/A}^\wedge$ は，イデアル $IB$ に関して
More on Algebra, Lemma \ref{more-algebra-lemma-ext-1-annihilated} の同値な条件
(1) -- (6) を満たす。
\item ある $c \geq 0$ が存在して，すべての $a \in I^c$ に対し写像
$h : \bigoplus B\text{d}x_i \to J/J^2$ が存在し，
$a : J/J^2 \to J/J^2$ は $h \circ \text{d}$ に等しい。
\item 元 $b_1, \ldots, b_s \in B$ が存在して
$V(b_1, \ldots, b_s) \subset V(IB)$ となり，しかも各
$l = 1, \ldots, s$ に対して，$m \geq 0$，$f_1, \ldots, f_m \in J$，および
$|T| = m$ を満たす部分集合 $T \subset \{1, \ldots, n\}$ が存在して，次が成り立つ：
\begin{enumerate}
\item $\det_{i \in T, j \leq m}(\partial f_j/ \partial x_i)$ は $B$ において
$b_l$ を割り切る。
\item $b_l J \subset (f_1, \ldots, f_m) + J^2$。
\end{enumerate}
\end{enumerate}
\end{lemma}


\begin{proof}
(1)，(2)，(3) の同値性は More on Algebra, Lemma
\ref{more-algebra-lemma-ext-1-annihilated} から直ちに従う。

\medskip\noindent
$b_1, \ldots, b_s$ が (4) のとおりであると仮定する。$B$ は Noether なので，
包含 $V(b_1, \ldots, b_s) \subset V(IB)$ から，ある $c \geq 0$ に対し
$I^cB \subset (b_1, \ldots, b_s)$ であることが従う
（例えば Algebra, Lemma
\ref{algebra-lemma-Noetherian-power-ideal-kills-module} による）。
$1 \leq l \leq s$，$m \geq 0$，$f_1, \ldots, f_m \in J$，および
$|T| = m$ を満たす $T \subset \{1, \ldots, n\}$ を，(4)(a)，(b) を満たすように選ぶ。
$b_l$ を可逆にすると，
$$
\NL_{B/A}^\wedge \otimes_B B_{b_l} =
\left(
\bigoplus\nolimits_{j \leq m} B_{b_l} f_j
\longrightarrow
\bigoplus\nolimits_{i = 1, \ldots, n} B_{b_l} \text{d}x_i
\right)
$$
であり，さらにこの射は直和因子
$\bigoplus_{i \in T} B_{b_l} \text{d}x_i$ の包含と同型である。したがって
$H^{-1}(\NL_{B/A}^\wedge)$ は $b_l$-冪捩れであり，
$H^0(\NL_{B/A}^\wedge)$ は $b_l$ を可逆にすると有限自由になる。
包含 $I^cB \subset (b_1, \ldots, b_s)$ と合わせると，
$H^{-1}(\NL_{B/A}^\wedge)$ は $IB$-冪捩れである。したがって
More on Algebra, Lemma \ref{more-algebra-lemma-ext-1-annihilated} の条件 (4) が
成り立つ。このようにして (4) から (2) が従う。

\medskip\noindent
同値な条件 (1)，(2)，(3) が成り立つと仮定する。(4) を証明するが，読者には
ぜひ自ら確認していただきたい。複体 $\NL_{B/A}^\wedge$ は
$D^b_{\textit{Coh}}(\Spec(B))$ の対象を定め，その Zariski 開集合
$U = \Spec(B) \setminus V(IB)$ への制限は，次数 $0$ に置かれた有限局所自由加群
$\mathcal{E}$ である（例えば More on Algebra, Lemma
\ref{more-algebra-lemma-ext-1-annihilated} の第四の同値条件から従う）。
$J$ の生成元 $f_1, \ldots, f_M$ を選ぶ。これにより完全列
$$
\bigoplus\nolimits_{j = 1, \ldots, M} \mathcal{O}_U \cdot f_j \to
\bigoplus\nolimits_{i = 1, \ldots, n} \mathcal{O}_U \cdot \text{d}x_i \to
\mathcal{E} \to 0
$$
が定まる。$U = \bigcup_{l = 1, \ldots, s} U_l$ を，ある整数
$n \geq m_l \geq 0$ に対して $\mathcal{E}|_{U_l}$ が階数
$r_l = n - m_l$ の自由加群になるような有限アフィン開被覆とする。
各 $U_l$ をアフィン開被覆で置き換えれば，部分集合
$T_l \subset \{1, \ldots, n\}$ が存在して，
$i \in \{1, \ldots, n\} \setminus T_l$ に対する元 $\text{d}x_i$ が
$\mathcal{E}|_{U_l}$ の基へ写ると仮定できる。同じ議論を繰り返せば，濃度 $m_l$ の
部分集合 $T'_l \subset \{1, \ldots, M\}$ が存在して，$j \in T'_l$ に対する $f_j$ が
$\mathcal{O}_{U_l} \cdot \text{d}x_i \to \mathcal{E}|_{U_l}$ の核の基へ写ると仮定できる。
最後に，開被覆 $U = \bigcup U_l$ は標準開集合による開被覆で細分できるので
（Algebra, Lemma \ref{algebra-lemma-Zariski-topology}），ある $g_l \in B$ に対して
$U_l = D(g_l)$ と仮定できる。
% P10-E0387
特に $V(g_1, \ldots, g_s) = V(IB)$ である。上の選択を用いた線形代数の議論により，
$\det_{i \in T_l, j \in T'_l}(\partial f_j/ \partial x_i)$ は
$B_{b_l}$ の可逆元へ写る。
% P10-E0388
同様に，$U_l$ 上で $\NL_{B/A}^\wedge$ の次数 $-1$ のコホモロジーが消えることから，
$J/J^2 + (f_j; j \in T')$ は $b_l$ のある冪により零化される。
% P10-E0389
各 $g_l$ を適当な冪で置き換えれば，補題の条件 (4)(a) と (4)(b) を得る。
いくつかの詳細は省略する。
\end{proof}

\begin{lemma}
\label{restricted-lemma-rig-smooth}
$A$ を Noether 環，$I$ をイデアル，$B$ を有限型 $A$-代数とする。
\begin{enumerate}
\item $\Spec(B) \to \Spec(A)$ が $\Spec(A) \setminus V(I)$ 上で smooth ならば，
$B^\wedge$ は $(A, I)$ 上 rig-smooth である。
\item $B^\wedge$ が $(A, I)$ 上 rig-smooth ならば，ある $g \in 1 + IB$ が存在して，
$\Spec(B_g)$ は $\Spec(A) \setminus V(I)$ 上で smooth である。
\end{enumerate}
\end{lemma}

\begin{proof}
以下，補題 \ref{restricted-lemma-equivalent-with-artin-smooth} を断りなく用いる。

\medskip\noindent
(1) を仮定する。$\NL_{B/A}$ の形成は局所化と可換であることを思い出そう。
Algebra, Lemma \ref{algebra-lemma-localize-NL} を参照せよ。したがって
smooth な環写像の定義そのもの（素朴な余接複体が次数 $0$ に置かれた有限射影加群に
擬同型であるという定義）から，$\NL_{B/A}$ はイデアル $IB$ に関して
More on Algebra, Lemma \ref{more-algebra-lemma-ext-1-annihilated} の第四の同値条件を
満たす（小さな詳細は省略する）。補題 \ref{restricted-lemma-NL-is-completion} により
$\NL_{B^\wedge/A}^\wedge = \NL_{B/A} \otimes_B B^\wedge$ なので，
More on Algebra, Lemma
\ref{more-algebra-lemma-base-change-property-ext-1-annihilated} によって (2) が従う。

\medskip\noindent
(2) を仮定する。表示 $B = A[x_1, \ldots, x_n]/J$ を選び，$N = J/J^2$ と置き，
$J/J^2$ 上の恒等写像が定める元
$\xi \in \Ext^1_B(\NL_{B/A}, J/J^2)$ を考える。再び
$\NL_{B^\wedge/A}^\wedge = \NL_{B/A} \otimes_B B^\wedge$ を用いると，仮定から
$$
\xi \otimes 1 \in
\Ext^1_{B^\wedge}(\NL_{B/A} \otimes_B B^\wedge, N \otimes_B B^\wedge) =
\Ext^1_{B^\wedge}(\NL_{B/A}, N) \otimes_B B^\wedge
$$
はある整数 $c \geq 0$ に対して $I^c$ により零化される。この等号は，例えば
More on Algebra, Lemma \ref{more-algebra-lemma-base-change-RHom} から従う
（ただし，はるかに簡単な More on Algebra, Lemma
\ref{more-algebra-lemma-pseudo-coherence-and-base-change-ext} から容易に導くこともできる）。
したがって $M = I^cB\xi \subset \Ext^1_B(\NL_{B/A}, N)$ は有限部分加群で，
$\Ext^1_B(\NL_{B/A}, N) \otimes_B B^\wedge$ において零へ写る。
$B \to B^\wedge$ は平坦なので，これは $M \otimes_B B^\wedge$ が零であることを意味する。
Nakayama の補題（Algebra, Lemma \ref{algebra-lemma-NAK}）により，
$M = I^cB\xi$ は $x \in IB$ を用いて $g = 1 + x$ と書ける元により零化される。
これは，各 $b \in I^cB$ に対して，次の図式を可換にする $B$-線形な点線の射が
存在することを意味する：
$$
\xymatrix{
J/J^2 \ar[r] \ar[d]^b & \bigoplus B\text{d}x_i \ar@{..>}[d]^h \\
J/J^2 \ar[r] & (J/J^2)_g
}
$$
したがって $(\NL_{B/A})_{gb}$ は有限射影加群に擬同型である。小さな詳細は省略する。
$D(B_{gb})$ において $(\NL_{B/A})_{gb} = \NL_{B_{gb}/A}$ なので，
$B_{gb}$ は $\Spec(A)$ 上 smooth である。これはすべての $b \in I^cB$ に対して
成り立つから，望みどおり $\Spec(B_g) \to \Spec(A)$ は
$\Spec(A) \setminus V(I)$ 上で smooth である。
\end{proof}

\begin{lemma}
\label{restricted-lemma-zero-ext-1-after-modding-out}
$(A_1, I_1) \to (A_2, I_2)$ を Remark \ref{restricted-remark-base-change} のものとし，
$A_1$ と $A_2$ は Noether であるとする。$B_1$ を $(A_1, I_1)$ に対する
(\ref{restricted-equation-C-prime}) の対象とし，$B_2$ を $B_1$ の基底変換とする。
$f_1 \in B_1$ とし，その像を $f_2 \in B_2$ とする。すべての
$B_1$-加群 $N_1$ に対し $\Ext^1_{B_1}(\NL_{B_1/A_1}^\wedge, N_1)$ が
$f_1$ により零化されるならば，すべての $B_2$-加群 $N_2$ に対し
$\Ext^1_{B_2}(\NL_{B_2/A_2}^\wedge, N_2)$ は $f_2$ により零化される。
\end{lemma}

\begin{proof}
補題 \ref{restricted-lemma-NL-base-change} により写像
$$
\NL_{B_1/A_1} \otimes_{B_2} B_1 \to \NL_{B_2/A_2}
$$
が存在し，これは $H^0$ 上で同型を，$H^{-1}$ 上で全射を誘導する。
% P10-E0384
したがって More on Algebra, Lemmas
\ref{more-algebra-lemma-two-term-base-change}，
\ref{more-algebra-lemma-base-change-property-ext-1-annihilated}，および
\ref{more-algebra-lemma-surjection-property-ext-1-annihilated} から結論が従う。
後二つはそれぞれ主イデアル $(f_1) \subset B_1$ と $(f_2) \subset B_2$ に適用する。
\end{proof}

\begin{lemma}
\label{restricted-lemma-base-change-rig-smooth-homomorphism}
$A_1 \to A_2$ を Noether 環の写像とする。$I_i \subset A_i$ を
$V(I_1A_2) = V(I_2)$ を満たすイデアルとする。$B_1$ を $(A_1, I_1)$ に対する
(\ref{restricted-equation-C-prime}) の対象とし，$B_2$ を Remark \ref{restricted-remark-base-change} の
意味での $B_1$ の基底変換とする。$B_1$ が $(A_1, I_1)$ 上 rig-smooth ならば，
$B_2$ は $(A_2, I_2)$ 上 rig-smooth である。
\end{lemma}

\begin{proof}
補題 \ref{restricted-lemma-zero-ext-1-after-modding-out}，定義
\ref{restricted-definition-rig-smooth-homomorphism}，および $A_2$ が Noether なので
ある $c \geq 0$ に対して $I_2^c$ が $I_1A_2$ に含まれるという事実から従う。
\end{proof}

\section{環準同型の変形}
\label{restricted-section-defos-ring-maps}

\noindent
rig-smooth 代数から環準同型を持ち上げることに関するいくつかの結果を述べる。

\begin{remark}[線形近似]
\label{restricted-remark-linear-approximation}
$A$ を環とし，$I \subset A$ を有限生成イデアルとする。
$C$ を $I$-進完備な $A$-代数とする。
$\psi : A[x_1, \ldots, x_r]^\wedge \to C$ を連続な
$A$-代数写像とする。$\delta_i \in C$（$i = 1, \ldots, r$）が
与えられているとする。このとき
$$
\psi' : A[x_1, \ldots, x_r]^\wedge \to C,\quad
x_i \longmapsto \psi(x_i) + \delta_i
$$
を考えることができる。Formal Spaces, Remark
\ref{formal-spaces-remark-universal-property} を参照せよ。このとき
$$
\psi'(g) = \psi(g) + \sum \psi(\partial g/\partial x_i)\delta_i + \xi
$$
が成り立つ。ここで誤差項は $\xi \in (\delta_i\delta_j)$ である。
これは $g$ を冪級数として書き，各項ごとに計算すれば従う。
$g$ の係数は零に収束するので，収束性は自動的である。
詳細は省略する。
\end{remark}

\begin{remark}[写像の持ち上げ]
\label{restricted-remark-improve-homomorphism}
$A$ を Noether 環とし，$I \subset A$ をイデアルとする。
$B$ を (\ref{restricted-equation-C-prime}) の対象とする。
$C$ を $I$-進完備な $A$-代数とする。
$\psi_n : B \to C/I^nC$ を $A$-代数準同型とする。
$\psi_n$ を $C/I^{2n}C$ への $A$-代数準同型に持ち上げることの
障害は元
$$
o(\psi_n) \in \Ext^1_B(\NL_{B/A}^\wedge, I^nC/I^{2n}C)
$$
であることを以下で説明する。実際，表示
$B = A[x_1, \ldots, x_r]^\wedge/J$ を選ぶ。
$\psi_n$ の持ち上げ
$\psi : A[x_1, \ldots, x_r]^\wedge \to C$ を選ぶ。
$\psi(J) \subset I^nC$ なので $\psi(J^2) \subset I^{2n}C$ であり，
したがって $B$-線形準同型
$$
o(\psi) :
J/J^2 \longrightarrow I^nC/I^{2n}C, \quad g \longmapsto \psi(g)
$$
を得る。これはもちろん $C$-線形写像
$J/J^2 \otimes_B C \to I^nC/I^{2n}C$ に延長される。
$\NL_{B/A}^\wedge = (J/J^2 \to \bigoplus B \text{d}x_i)$ なので，
次の同一視のもとで $o(\psi)$ の像として $o(\psi_n)$ を得る：
\begin{align*}
& \Ext^1_B(\NL_{B/A}^\wedge, I^nC/I^{2n}C) \\
& =
\Coker\left(\Hom_B(\bigoplus B\text{d}x_i, I^nC/I^{2n}C) \to
\Hom_B(J/J^2, I^nC/I^{2n}C)\right)
\end{align*}
この等号については More on Algebra, Lemma
\ref{more-algebra-lemma-map-out-of-almost-free} の (1) を参照せよ。

\medskip\noindent
ある整数 $n'$ が $n > n' > n/2$ を満たし，$o(\psi_n)$ の
$\Ext^1_B(\NL_{B/A}^\wedge, I^{n'}C/I^{2n'}C)$ への像が零であるとする。
このことから，$C/I^{n'}C$ への写像として $\psi_n$ と一致する
$A$-代数準同型 $\psi'_{2n'} : B \to C/I^{2n'}C$ を見いだせると主張する。
極端な場合 $n' = n$ は，先に $o(\psi_n)$ を $\psi_n$ を
$C/I^{2n}C$ へ持ち上げることの障害と呼んだ理由を説明している。
主張を証明しよう。$o(\psi_n)$ の像が零であるという仮定から，
$B$-加群写像
$$
h : \bigoplus B\text{d}x_i \longrightarrow I^{n'}C/I^{2n'}C
$$
であって，$o(\psi)$ と $h \circ \text{d}$ が
$I^{n'}C/I^{2n'}C$ への写像として一致するものを見いだせる。
ある $\delta_i \in I^{n'}C$ に対して
$h(\text{d}x_i) = \delta_i \bmod I^{2n'}C$ と書く。このとき写像
$$
\psi' : A[x_1, \ldots, x_r]^\wedge \to C,\quad
x_i \longmapsto \psi(x_i) - \delta_i
$$
を考える。冪級数による計算から $\psi'(J) \subset I^{2n'}C$ が分かる。
実際，$g \in J$ に対して
$$
\psi'(g) \equiv
\psi(g) - \sum \psi(\partial g/\partial x_i)\delta_i \equiv
o(\psi)(g) - (h \circ \text{d})(g) \equiv
0 \bmod I^{2n'}C
$$
を得る。最初の等号については Remark
\ref{restricted-remark-linear-approximation} を参照せよ。
したがって $\psi'$ は，望まれる $A$-代数準同型
$\psi'_{2n'} : B \to C/I^{2n'}C$ を誘導する。
\end{remark}

\begin{lemma}
\label{restricted-lemma-get-morphism-general-better}
次のデータが与えられているとする：
\begin{enumerate}
\item 整数 $c \geq 0$，
\item Noether 環 $A$ のイデアル $I$，
\item $(A, I)$ に対する (\ref{restricted-equation-C-prime}) の対象 $B$ であって，
任意の $B$-加群 $N$ に対して $I^c$ が
$\Ext^1_B(\NL_{B/A}^\wedge, N)$ を零化するもの，
\item Noether な $I$-進完備 $A$-代数 $C$；Local Cohomology,
Section \ref{local-cohomology-section-uniform} で得られる整数を
$d = d(\text{Gr}_I(C))$ および $q_0 = q(\text{Gr}_I(C))$ と書く，
\item 整数 $n \geq \max(q_0 + (d + 1)c, 2(d + 1)c + 1)$，および
\item $A$-代数準同型 $\psi_n : B \to C/I^nC$。
\end{enumerate}
このとき，$\psi_n \bmod I^{n - (d + 1)c} =
\varphi \bmod I^{n - (d + 1)c}$ を満たす $A$-代数写像
$\varphi : B \to C$ が存在する。
\end{lemma}

\begin{proof}
Remark \ref{restricted-remark-improve-homomorphism} の障害類
$$
o(\psi_n) \in \Ext^1_B(\NL_{B/A}^\wedge, I^nC/I^{2n}C)
$$
を考える。任意の $C/I^nC$-加群 $N$ に対して
\begin{align*}
\Ext^1_B(\NL_{B/A}^\wedge, N)
& =
\Ext^1_{C/I^nC}(\NL_{B/A}^\wedge \otimes_B^\mathbf{L} C/I^nC, N) \\
& =
\Ext^1_{C/I^nC}(\NL_{B/A}^\wedge \otimes_B C/I^nC, N)
\end{align*}
が成り立つ。最初の等号は More on Algebra, Lemma
\ref{more-algebra-lemma-upgrade-adjoint-tensor-RHom} により，二つ目は
More on Algebra, Lemma \ref{more-algebra-lemma-two-term-base-change}
により従う。特に，すべての $C/I^nC$-加群 $N$ に対して
$\Ext^1_{C/I^nC}(\NL_{B/A}^\wedge \otimes_B C/I^nC, N)$ は
$I^cC$ により零化される。したがって Local Cohomology, Lemma
\ref{local-cohomology-lemma-bound-two-term-complex} を適用すると，
$o(\psi_n)$ の次の群への像が零であることが分かる。
% P10-E0390
$$
\Ext^1_{C/I^nC}(\NL_{B/A}^\wedge \otimes_B C/I^nC, I^{n'}C/I^{2n'}C) =
\Ext^1_B(\NL_{B/A}^\wedge, I^{n'}C/I^{2n'}C) =
$$
ここで $n' = n - (d + 1)c$ である。Remark
\ref{restricted-remark-improve-homomorphism} の議論により，
$$
\psi'_{2n'} : B \to C/I^{2n'}C
$$
を得る。この写像は $\psi_n$ と法 $I^{n'}$ で一致する。
$n \geq 2(d + 1)c + 1$ なので $2n' > n$ であることに注意せよ。

\medskip\noindent
この手続きを繰り返すことができる。$n_0 = n$ および
$\psi^0 = \psi_n$ から始めると，最終的に狭義単調増加する整数列
$$
n_0 < n_1 < n_2 < \ldots
$$
と，$\psi^{i + 1}$ と $\psi^i$ が法 $I^{n_i - tc}$ で一致するような
$A$-代数準同型 $\psi^i : B \to C/I^{n_i}C$ を得る。
$C$ は $I$-進完備なので，$\varphi$ を写像
$\psi^i \bmod I^{n_i - (d + 1)c} : B \to C/I^{n_i - (d + 1)c}C$
の極限として取ることができ，補題が従う。
\end{proof}

\noindent
読者には次の節まで先へ進むことを勧める。実際，以下の二つの補題は，
そこに現れる代数 $C$ が Noether であると仮定すれば，上の結果から従う。

\begin{lemma}
\label{restricted-lemma-get-morphism-nonzerodivisor}
$I = (a)$ を Noether 環 $A$ の単項イデアルとする。
$B$ を (\ref{restricted-equation-C-prime}) の対象とする。
すべての $B$-加群 $N$ に対して
$\Ext^1_B(\NL_{B/A}^\wedge, N)$ が $a^c$ により零化されるような
整数 $c \geq 0$ が与えられているとする。
$C$ を $I$-進完備な $A$-代数とし，$a$ は $C$ 上の零因子でないとする。
$n > 2c$ とする。任意の $A$-代数写像
$\psi_n : B \to C/a^nC$ に対して，
$\psi_n \bmod a^{n - c}C = \varphi \bmod a^{n - c}C$ を満たす
$A$-代数写像 $\varphi : B \to C$ が存在する。
\end{lemma}

\begin{proof}
Remark \ref{restricted-remark-improve-homomorphism} の障害類
$$
o(\psi_n) \in \Ext^1_B(\NL_{B/A}^\wedge, a^nC/a^{2n}C)
$$
を考える。$a$ は $C$ 上の零因子でないので，$C$-加群の圏において写像
$a^c : a^nC/a^{2n}C \to a^nC/a^{2n}C$ は写像
$a^nC/a^{2n}C \to a^{n - c}C/a^{2n - c}C$ と同型である。
したがって $\NL_{B/A}^\wedge$ に関する仮定により，
類 $o(\psi_n)$ の
$$
\Ext^1_B(\NL_{B/A}^\wedge, a^{n - c}C/a^{2n - c}C)
$$
への像は零であり，したがってなおさら
$$
\Ext^1_B(\NL_{B/A}^\wedge, a^{n - c}C/a^{2n - 2c}C)
$$
への像も零である。Remark \ref{restricted-remark-improve-homomorphism} の議論により，
$$
\psi_{2n - 2c} : B \to C/a^{2n - 2c}C
$$
を得る。この写像は $\psi_n$ と法 $a^{n - c}C$ で一致する。
$n > 2c$ なので $2n - 2c > n$ であることに注意せよ。

\medskip\noindent
この手続きを繰り返すことができる。$n_0 = n$ および
$\psi^0 = \psi_n$ から始めると，最終的に狭義単調増加する整数列
$$
n_0 < n_1 < n_2 < \ldots
$$
と，$\psi^{i + 1}$ と $\psi^i$ が法 $a^{n_i - c}C$ で一致するような
$A$-代数準同型 $\psi^i : B \to C/a^{n_i}C$ を得る。
$C$ は $I$-進完備なので，$\varphi$ を写像
$\psi^i \bmod a^{n_i - c}C : B \to C/a^{n_i - c}C$
の極限として取ることができ，補題が従う。
\end{proof}

\begin{lemma}
\label{restricted-lemma-get-morphism-principal}
$I = (a)$ を Noether 環 $A$ の単項イデアルとする。
$B$ を (\ref{restricted-equation-C-prime}) の対象とする。
すべての $B$-加群 $N$ に対して
$\Ext^1_B(\NL_{B/A}^\wedge, N)$ が $a^c$ により零化されるような
整数 $c \geq 0$ が与えられているとする。
$C$ を $I$-進完備な $A$-代数とする。
$C[a^\infty] \cap a^dC = 0$ を満たす整数 $d \geq 0$ が
与えられているとする。$n > \max(2c, c + d)$ とする。
任意の $A$-代数写像 $\psi_n : B \to C/a^nC$ に対して，
$\psi_n \bmod a^{n - c} = \varphi \bmod a^{n - c}$ を満たす
$A$-代数写像 $\varphi : B \to C$ が存在する。
\end{lemma}

\noindent
$C$ が Noether ならば，ある $e \geq 0$ に対して
$C[a^\infty] = C[a^e]$ である。Artin--Rees
（Algebra, Lemma \ref{algebra-lemma-Artin-Rees}）により，
すべての $n \geq f$ に対して
$a^nC \cap C[a^\infty] \subset a^{n - f}C[a^\infty]$ となる
整数 $f$ が存在する。このとき $d = e + f$ は補題に現れる整数である。
この議論は，$C$ が (\ref{restricted-equation-C-prime}) の対象である場合に，
Lemma \ref{restricted-lemma-topologically-finite-type-Noetherian} により特に適用できる。

\begin{proof}
$C \to C'$ を $C$ の $C[a^\infty]$ による商とする。
$A$-代数 $C'$ は Algebra, Lemma
\ref{algebra-lemma-quotient-complete}，および
$\bigcap (C[a^\infty] + a^nC) = C[a^\infty]$ という事実により
$I$-進完備である。後者は，$n \geq d$ のとき和
$C[a^\infty] + a^nC$ が直和であることから従う。
$m \geq d$ に対して，図式
$$
\xymatrix{
0 \ar[r] &
C[a^\infty] \ar[r] \ar[d] &
C \ar[r] \ar[d] & C' \ar[r] \ar[d] & 0 \\
0 \ar[r] &
C[a^\infty] \ar[r] &
C/a^m C \ar[r] & C'/a^m C' \ar[r] & 0
}
$$
の各行は完全である。したがって，すべての $m \geq d$ に対して
$C$ は $C'$ と $C/a^mC$ の $C'/a^mC'$ 上のファイバー積である。
Lemma \ref{restricted-lemma-get-morphism-nonzerodivisor} により，
$\varphi'$ と $\psi_n$ が $C'/a^{n - c}C'$ への写像として一致するような
準同型 $\varphi' : B \to C'$ を選ぶことができる。
準同型 $(\varphi', \psi_n \bmod a^{n - c}C) : B \to
C' \times_{C'/a^{n - c}C'} C/a^{n - c}C$ を得る。
$n - c \geq d$ なので，これは準同型 $\varphi : B \to C$ と
同じものである。これで証明は完了する。
\end{proof}

\section{G-環上の rig-smooth 代数の代数化}
\label{restricted-section-over-G-ring}

\noindent
基礎環 $A$ が Noether G-環ならば，任意のイデアル $I \subset A$
（単項とは限らない）に関する任意の rig-smooth 代数に対して
\cite[III Theorem 7]{Elkik} を証明できる。

\begin{lemma}
\label{restricted-lemma-close-enough}
$I$ を Noether 環 $A$ のイデアルとする。$r \geq 0$ とし，
% P10-E0649
$P = A[x_1, \ldots, x_r]$ を $I$-進完備化と書く。
$P$ のある商の分解
$$
P^{\oplus t} \xrightarrow{K} P^{\oplus m}
\xrightarrow{g_1, \ldots, g_m} P \to B \to 0
$$
を考える。$B$ は $(A, I)$ 上 rig-smooth であると仮定する。
このとき，任意の複体
$$
P^{\oplus t} \xrightarrow{K'} P^{\oplus m}
\xrightarrow{g'_1, \ldots, g'_m} P
$$
であって $g_i - g'_i \in I^nP$ および
$K - K' \in I^n\text{Mat}(m \times t, P)$ を満たすものに対し，
$B' = P/(g'_1, \ldots, g'_m)$ と置けば，$A$-代数の同型
$B \to B'$ が存在するような整数 $n$ が存在する。
\end{lemma}

\begin{proof}
(A) Definition \ref{restricted-definition-rig-smooth-homomorphism} により，
すべての $B$-加群 $N$ に対して $I^c$ が
$\Ext^1_B(\NL_{B/A}^\wedge, N)$ を零化するような
$c \geq 0$ を選ぶことができる。

\medskip\noindent
(B) More on Algebra, Lemmas \ref{more-algebra-lemma-approximate-complex} および
\ref{more-algebra-lemma-approximate-complex-graded} により，
$n \geq c_1 + 1$ のとき複体
$$
P^{\oplus t} \xrightarrow{K'} P^{\oplus m}
\xrightarrow{g'_1, \ldots, g'_m} P \to B' \to 0
$$
が完全で，$\text{Gr}_I(B) \cong \text{Gr}_I(B')$ となるような定数
$c_1 = c(g_1, \ldots, g_m, K)$ が存在する。

\medskip\noindent
(C) $d_0 = d(\text{Gr}_I(B))$ および $q_0 = q(\text{Gr}_I(B))$ を
Local Cohomology, Section \ref{local-cohomology-section-uniform}
で得られる整数とする。

\medskip\noindent
(A) の $c$，(B) の $c_1$，および (C) の $q_0, d_0$ に対して
$n = \max(c_1 + 1, q_0 + (d_0 + 1)c, 2(d_0 + 1)c + 1)$
が機能すると主張する。

\medskip\noindent
$g'_1, \ldots, g'_m$ および $K'$ を補題のものとする。
% P10-E0650
$g_i = g'_i \in I^nP$ なので，標準的な $A$-代数準同型
$$
\psi_n : B \longrightarrow B'/I^nB'
$$
を得る。これは同型 $B/I^nB \to B'/I^nB'$ を誘導する。
$\text{Gr}_I(B) \cong \text{Gr}_I(B')$ なので
$d_0 = d(\text{Gr}_I(B'))$ および $q_0 = q(\text{Gr}_I(B'))$ であり，
$n \geq \max(q_0 + (1 + d_0)c, 2(d_0 + 1)c + 1)$ なので
Lemma \ref{restricted-lemma-get-morphism-general-better} を適用して
$A$-代数準同型
$$
\varphi : B \longrightarrow B'
$$
であって
$\varphi \bmod I^{n - (d_0 + 1)c}B' =
\psi_n \bmod I^{n - (d_0 + 1)c}B'$ を満たすものを得る。
$n - (d_0 + 1)c > 0$ なので，$\varphi$ は法 $I$ で
先に得た同型 $B/IB \to B'/IB'$ を誘導する $A$-代数準同型である。
証明の残りでは，これらの事実が $\varphi$ を同型にすることを示す。
読者には自分自身でその証明を見いだすことを勧める。

\medskip\noindent
具体的には，例えば $B$ と $B'$ が完備であることを用いて
Algebra, Lemma \ref{algebra-lemma-completion-generalities} の (1) を
適用すれば，$\varphi$ が全射であることが従う。
したがって $\varphi$ は全射
$\text{Gr}_I(B) \to \text{Gr}_I(B')$ を誘導する。
その始域と終域は互いに同型な Noether 環なので，この全射は同型でなければならない。
Algebra, Lemma
\ref{algebra-lemma-surjective-endo-noetherian-ring-is-iso} を参照せよ
（もちろん，$\varphi$ が先に得た同型を誘導することを直接示してもよいが，
そのためにはわずかな議論が必要となる）。
ゆえに $\varphi$ は，すべての $e \geq 0$ に対して単射
$I^eB/I^{e + 1}B \to I^eB'/I^{e + 1}B'$ を誘導する。
Krull の共通部分定理（Algebra, Lemma
\ref{algebra-lemma-intersect-powers-ideal-module-zero}）により，
任意の $b \in B$ に対して $b \in I^eB$ かつ
$b \not \in I^{e + 1}B$ となる $e \geq 0$ が存在するので，
これは $\varphi$ が単射であることを含意する。これで証明は完了する。
\end{proof}

\begin{lemma}
\label{restricted-lemma-algebraize-easy}
$I$ を Noether 環 $A$ のイデアルとする。$C^h$ を，有限型
$A$-代数 $C$ のイデアル $IC$ に関するヘンゼル化とする。
$J \subset C^h$ をイデアルとする。このとき
$B^\wedge \cong (C^h/J)^\wedge$ を満たす有限型 $A$-代数 $B$ が存在する。
\end{lemma}

\begin{proof}
More on Algebra, Lemma
\ref{more-algebra-lemma-henselization-Noetherian-pair} により，
環 $C^h$ は Noether である。$J = (g_1, \ldots, g_m)$ と書く。
環 $C^h$ は，$C/IC \to C'/IC'$ が同型となるような
エタール $C$-代数 $C'$ のフィルター付き余極限である
（More on Algebra, Lemma \ref{more-algebra-lemma-henselization}
の証明を参照せよ）。
% P10-E0391
$g_1, \ldots, g_m$ が $g'_1, \ldots, g'_m \in C'$ の像となるような
$C'$ を一つ選ぶ。$B = C'/(g'_1, \ldots, g'_m)$ と置けば，
有限型 $A$-代数を得る。もちろん
% P10-E0392
$(C, IC)$ と $C', IC')$ は同じヘンゼル化と同じ完備化をもつ。
このことから $B^\wedge = (C^h/J)^\wedge$ が容易に従う。
\end{proof}

\begin{proposition}
\label{restricted-proposition-approximate}
$I$ を Noether G-環 $A$ のイデアルとする。
$B$ を (\ref{restricted-equation-C-prime}) の対象とする。
$B$ が $(A, I)$ 上 rig-smooth ならば，有限型 $A$-代数 $C$ と
$A$-代数の同型 $B \cong C^\wedge$ が存在する。
\end{proposition}

\begin{proof}
表示 $B = A[x_1, \ldots, x_r]^\wedge/J$ を選ぶ。
$P = A[x_1, \ldots, x_r]^\wedge$ と書く。
このイデアルの生成元 $g_1, \ldots, g_m \in J$ を選ぶ。
$g_1, \ldots, g_m$ の間の関係の加群の生成元
$k_1, \ldots, k_t$，すなわち
$$
P^{\oplus t} \xrightarrow{k_1, \ldots, k_t}
P^{\oplus m} \xrightarrow{g_1, \ldots, g_m}
P \to B \to 0
$$
が分解となるようなものを選ぶ。$k_i = (k_{i1}, \ldots, k_{im})$
と書けば，
\begin{equation}
\label{restricted-equation-relations-straight-up}
\sum\nolimits_j k_{ij}g_j = 0
\end{equation}
が $i = 1, \ldots, t$ に対して成り立つ。
$K = (k_{ij})$ を成分 $k_{ij}$ をもつ $m \times t$-行列と書く。

\medskip\noindent
$A[x_1, \ldots, x_r]^h$ を組
$(A[x_1, \ldots, x_r], IA[x_1, \ldots, x_r])$ のヘンゼル化とする。
More on Algebra, Lemma \ref{more-algebra-lemma-henselization} を参照せよ。
More on Algebra, Lemma
\ref{more-algebra-lemma-henselization-Noetherian-pair} により，
$A[x_1, \ldots, x_r]^h$ を
$P = A[x_1, \ldots, x_r]^\wedge$ の部分環とみなしてよく，実際そうする。
$A$ は Noether G-環なので $A[x_1, \ldots, x_r]$ もそうである。
More on Algebra, Proposition
\ref{more-algebra-proposition-finite-type-over-G-ring} を参照せよ。
したがって，$I$ により生成されるイデアルに関して，写像
$A[x_1, \ldots, x_r]^h \to A[x_1, \ldots, x_r]^\wedge = P$
に対する近似が成り立つ。Smoothing Ring Maps, Lemma
\ref{smoothing-lemma-henselian-pair} を参照せよ。
Lemma \ref{restricted-lemma-close-enough} に現れる十分大きな整数 $n$ を選ぶ。
近似性質により，$g'_1, \ldots, g'_m \in A[x_1, \ldots, x_r]^h$
および行列
$K' = (k'_{ij}) \in \text{Mat}(m \times t, A[x_1, \ldots, x_r]^h)$
であって，$A[x_1, \ldots, x_r]^h$ において
$\sum\nolimits_j k'_{ij}g'_j = 0$ を満たし，しかも
$g_i - g'_i \in I^nP$ および
$K - K' \in I^n\text{Mat}(m \times t, P)$ を満たすものを選べる。
$n$ の選び方により，同型
$$
B \to P/(g'_1, \ldots, g'_m) =
\left(A[x_1, \ldots, x_r]^h/(g'_1, \ldots, g'_m)\right)^\wedge
$$
が存在する。Lemma \ref{restricted-lemma-algebraize-easy} により証明が完了する。
\end{proof}

\noindent
次の補題は，$A$ が G-環でなく単に Noether である場合には一般には正しくない。
実際，$(A, \mathfrak m)$ が局所環で $I = \mathfrak m$ ならば，
この補題は $A^h$ に対する Artin 近似
（Smoothing Ring Maps, Theorem
\ref{smoothing-theorem-approximation-property} におけるもの）と同値であり，
これはすべての Noether 局所環に対して成り立つわけではない。

\begin{lemma}
\label{restricted-lemma-fully-faithfulness}
% P10-E0651
$A$ を Noether G-環とする。$I \subset A$ をイデアルとする。
$B, C$ を有限型 $A$-代数とする。$I$-進完備化の任意の
$A$-代数写像 $\varphi : B^\wedge \to C^\wedge$ と任意の
$N \geq 1$ に対して，次のものが存在する：
\begin{enumerate}
\item 同型 $C/IC \to C'/IC'$ を誘導するエタール環写像 $C \to C'$，
\item $A$-代数写像 $\varphi : B \to C'$。
\end{enumerate}
これらは，$\varphi$ と $\psi$ が
$C^\wedge = (C')^\wedge$ への写像として法 $I^N$ で一致するように選べる。
\end{lemma}

\begin{proof}
$C \to C'$ は平坦であり
（Algebra, Lemma \ref{algebra-lemma-etale}），したがってすべての $n$ に対して
同型 $C/I^nC \to C'/I^nC'$ を誘導し
（More on Algebra, Lemma
\ref{more-algebra-lemma-neighbourhood-isomorphism}），ゆえに完備化上の同型を
誘導するので，補題の主張には意味がある。
$C^h$ を組 $(C, IC)$ のヘンゼル化とする。More on Algebra, Lemma
\ref{more-algebra-lemma-henselization} を参照せよ。
このとき $C^h$ は代数 $C'$ のフィルター付き余極限であり，写像
$C \to C' \to C^h$ は完備化上の同型を誘導する
（More on Algebra, Lemma
\ref{more-algebra-lemma-henselization-Noetherian-pair}）。
したがって，法 $I^N$ で $\psi$ と合同な $A$-代数写像
$B \to C^h$ が存在することを示せば十分である。
$B = A[x_1, \ldots, x_n]/(f_1, \ldots, f_m)$ と書く。
環写像 $\psi$ は，$j = 1, \ldots, m$ に対して
$f_j(\hat c_1, \ldots, \hat c_n) = 0$ を満たす元
$\hat c_1, \ldots, \hat c_n \in C^\wedge$ に対応する。
$A$ は Noether G-環なので $C$ もそうである。
More on Algebra, Proposition
\ref{more-algebra-proposition-finite-type-over-G-ring} を参照せよ。
したがって Smoothing Ring Maps, Lemma
\ref{smoothing-lemma-henselian-pair} を適用すると，
$j = 1, \ldots, m$ に対して
$f_j(c_1, \ldots, c_n) = 0$ を満たし，かつ
$\hat c_i - c_i \in I^NC^h$ を満たす元
$c_1, \ldots, c_n \in C^h$ を得る。
これが望まれる写像 $B \to C^h$ を定める。
\end{proof}

\section{rig-smooth 代数の代数化}
\label{restricted-section-algebraization-rig-smooth}

\noindent
rig-smooth 代数が特定の表示をもてば，その代数化は直接的に行えることが分かる。
この点の議論については Remark \ref{restricted-remark-discussion} も参照せよ。

\begin{lemma}
\label{restricted-lemma-presentation-rig-smooth}
$A$ を環とする。$f_1, \ldots, f_m \in A[x_1, \ldots, x_n]$ とし，
$B = A[x_1, \ldots, x_n]/(f_1, \ldots, f_m)$ と置く。
$m \leq n$ と仮定し，
$g = \det_{1 \leq i, j \leq m}(\partial f_j/\partial x_i)$ と置く。
このとき
\begin{enumerate}
\item すべての $B$-加群 $N$ に対して
$g$ は $\Ext^1_B(\NL_{B/A}, N)$ を零化する。
\item $n = m$ ならば，$\NL_{B/A}$ 上の $g$ による乗法は
$D(B)$ において $0$ である。
\end{enumerate}
\end{lemma}

\begin{proof}
% P10-E0393
$T$ を，$1 \leq i, j \leq n$ に対する
成分 $\partial f_j/\partial x_i$ をもつ $m \times m$ 行列とする。
$K \in D(B)$ を，次数 $-1$ と $0$ に項をもつ複体
$T : B^{\oplus m} \to B^{\oplus m}$ により表されるものとする。
More on Algebra, Lemmas \ref{more-algebra-lemma-silly} により，
$g : K \to K$ は $D(B)$ において零である。
$J = (f_1, \ldots, f_m)$ と置く。
$\NL_{B/A}$ は
$J/J^2 \to \bigoplus_{i = 1, \ldots, n} B\text{d}x_i$
とホモトピー同値であることを思い出そう。
Algebra, Section \ref{algebra-section-netherlander} を参照せよ。
$L$ を複体
$J/J^2 \to \bigoplus_{i = 1, \ldots, m} B\text{d}x_i$ と書く。
% P10-E0652
このとき商写像 $\NL_{B/A} \to L$ がある。
また，次数 $-1$ の項 $B^{\oplus m}$ の第 $j$ 基底元を
$J/J^2$ における $f_j$ の類へ送ることにより，複体の全射
$K \to L$ も得る。図式で書けば
$$
\NL_{B/A} \to L \leftarrow K
$$
である。More on Algebra, Lemma
\ref{more-algebra-lemma-two-term-surjection-map-zero} から，
$L$ 上の $g$ による乗法は $D(B)$ において $0$ であると結論する。
一方，識別三角
$B^{\oplus n - m}[0] \to \NL_{B/A} \to L$ から，
すべての $B$-加群 $N$ に対して
$\Ext^1_B(L, N) \to \Ext^1_B(\NL_{B/A}, N)$ は全射であり，
したがって $g$ により零化されることが分かる。これで (1) が示された。
$n = m$ ならば $\NL_{B/A} = L$ なので，(2) が成り立つ。
\end{proof}

\begin{lemma}
\label{restricted-lemma-approximate-presentation-rig-smooth}
$I$ を Noether 環 $A$ のイデアルとする。
$B$ を (\ref{restricted-equation-C-prime}) の対象とする。
$B = A[x_1, \ldots, x_r]^\wedge/J$ を表示とする。
元 $b \in B$，整数 $0 \leq m \leq r$，および
$f_1, \ldots, f_m \in J$ であって，次を満たすものが存在すると仮定する：
\begin{enumerate}
\item $\Spec(B)$ において $V(b) \subset V(IB)$ である。
\item
$\Delta = \det_{1 \leq i, j \leq m}(\partial f_j/\partial x_i)$
の $B$ における像は $b$ を割り切る。
\item $b J \subset (f_1, \ldots, f_m) + J^2$ である。
\end{enumerate}
このとき，有限型 $A$-代数 $C$ と $A$-代数の同型
$B \cong C^\wedge$ が存在する。
\end{lemma}

\begin{proof}
これらの条件は $B$ が $(A, I)$ 上 rig-smooth であることを含意する。
Lemma \ref{restricted-lemma-equivalent-with-artin-smooth} を参照せよ。
ある $b' \in B$ に対して $B$ において $b' \Delta = b$ と書く。
$I = (a_1, \ldots, a_t)$ とする。$V(b) \subset V(IB)$ なので，
$I^cB \subset bB$ を満たす整数 $c \geq 0$ が存在する。
ある $b_i \in B$ に対して $B$ において $bb_i = a_i^c$ と書く。

\medskip\noindent
整数 $n \gg 0$ を選ぶ（どの程度大きく取るかは後で分かる）。
$f_i - f'_i \in I^nA[x_1, \ldots, x_r]^\wedge$ を満たす多項式
$f'_1, \ldots, f'_m \in A[x_1, \ldots, x_r]$ を選ぶ。
$\Delta' = \det_{1 \leq i, j \leq m}(\partial f'_j/\partial x_i)$ と置き，
有限型 $A$-代数
$$
C = A[x_1, \ldots, x_r, z_1, \ldots, z_t]/
(f'_1, \ldots, f'_m,
z_1\Delta' - a_1^c, \ldots, z_t\Delta' - a_t^c)
$$
を考える。$C$ に Lemma \ref{restricted-lemma-presentation-rig-smooth} を適用する。
計算すると
$$
\det\left(
\begin{matrix}
\text{次の式の偏微分からなる行列} \\
f'_1, \ldots, f'_m, z_1\Delta' - a_1^c, \ldots, z_t\Delta' - a_t^c \\
\text{次の変数に関して} \\
x_1, \ldots, x_m, z_1, \ldots, z_t
\end{matrix}
\right) =
(\Delta')^{t + 1}
$$
を得る。したがって，すべての $C$-加群 $N$ に対して
$\Ext^1_C(\NL_{C/A}, N)$ は $(\Delta')^{t + 1}$ により零化される。
$C$ において $a_i^c$ は $\Delta'$ で割り切れるので，
$a_i^{(t + 1)c}$ もこれらの $\Ext^1$ を零化する。
したがって，すべての $C$-加群 $N$ に対して
$I^{c_1}$ は $\Ext^1_C(\NL_{C/A}, N)$ を零化する。
ここで $c_1 = 1 + t((t + 1)c - 1)$ である。
$c_1$ の正確な値は議論の残りには重要でない。
重要なのは，それが $n$ に依存しないことである。

\medskip\noindent
Lemma \ref{restricted-lemma-NL-is-completion} により
$\NL_{C^\wedge/A}^\wedge = \NL_{C/A} \otimes_C C^\wedge$ なので，
任意の $C^\wedge$-加群 $N$ に対しても，$I^{c_1}$ による乗法は
$\Ext^1_{C^\wedge}(\NL_{C^\wedge/A}^\wedge, N)$ 上で零である。
More on Algebra, Lemmas
\ref{more-algebra-lemma-base-change-property-ext-1-annihilated} および
\ref{more-algebra-lemma-two-term-base-change} を参照せよ。
特に，$C^\wedge$ は $(A, I)$ 上 rig-smooth である。

\medskip\noindent
類 $x_i$ を類 $x_i$ へ送り，類 $z_i$ を類 $b_ib'$ へ送る
全射な $A$-代数準同型
$$
\psi_n : C \longrightarrow B/I^nB
$$
が存在することに注意せよ。これが機能するのは，証明の最初の段落で
$b'$ および $b_i$ を選んだ仕方による。

\medskip\noindent
$d = d(\text{Gr}_I(B))$ および $q_0 = q(\text{Gr}_I(B))$ を
Local Cohomology, Section \ref{local-cohomology-section-uniform}
で得られる整数とする。Lemma \ref{restricted-lemma-get-morphism-general-better} により，
$n \geq \max(q_0 + (d + 1)c_1, 2(d + 1)c_1 + 1)$ と取れば，
$\psi_n$ と法 $I^{n - (d + 1)c_1}B$ で合同な
$A$-代数準同型 $\varphi : C^\wedge \to B$ を見いだせる。

\medskip\noindent
$\varphi : C^\wedge \to B$ は法 $I$ で全射なので，全射である
（例えば Algebra, Lemma
\ref{algebra-lemma-completion-generalities} を用いよ）。
証明を完了するには，
$\Ker(\varphi)/\Ker(\varphi)^2$ が $I$ のある冪により零化されることを
示せば十分である。More on Algebra, Lemma
\ref{more-algebra-lemma-quotient-by-idempotent} を参照せよ。

\medskip\noindent
$\varphi$ は全射なので，$\NL_{B/C^{\wedge}}^\wedge$ のコホモロジー加群は
$H^0(\NL_{B/C^{\wedge}}^\wedge) = 0$ および
$H^{-1}(\NL_{B/C^{\wedge}}^\wedge) = \Ker(\varphi)/\Ker(\varphi)^2$
である。Lemma \ref{restricted-lemma-exact-sequence-NL} による完全列
$$
H^{-1}(\NL_{C^\wedge/A}^\wedge \otimes_{C^\wedge} B) \to
H^{-1}(\NL_{B/A}^\wedge) \to
H^{-1}(\NL_{B/C^{\wedge}}^\wedge) \to
H^0(\NL_{C^\wedge/A}^\wedge \otimes_{C^\wedge} B) \to
H^0(\NL_{B/A}^\wedge) \to 0
$$
がある。最初の二つの加群は，$B$ と $C^\wedge$ が
$(A, I)$ 上 rig-smooth なので，$I$ のある冪により零化される。
したがって，全射
$H^0(\NL_{C^\wedge/A}^\wedge \otimes_{C^\wedge} B) \to
H^0(\NL_{B/A}^\wedge)$ の核が $I$ のある冪により零化されることを
示せば十分である。そのためには，それが $b$ のある冪により
零化されることを示せば十分である。言い換えると，
$$
H^0(\NL_{C^\wedge/A}^\wedge) \otimes_{C^\wedge} B[1/b]
\longrightarrow
H^0(\NL_{B/A}^\wedge) \otimes_B B[1/b]
$$
が同型であることを示せば十分である。しかし，どちらも
$\text{d}x_{m + 1}, \ldots, \text{d}x_r$ を基底とする階数
$r - m$ の自由 $B[1/b]$-加群である。以上で証明が完了する。
\end{proof}

\begin{remark}
\label{restricted-remark-discussion}
$I$ を Noether 環 $A$ のイデアルとする。
$B$ を $(A, I)$ 上 rig-smooth な
(\ref{restricted-equation-C-prime}) の対象とする。
\cite[Theorem 1.2]{gabber-zavyalov} において，$B$ は有限型
$A$-代数の $I$-進完備化と同型であることが示されている。
% P10-E0394
この結果は，次の部分的な結果の一覧に取って代わる：
\begin{enumerate}
\item $A$ が G-環ならば，結果は Proposition
\ref{restricted-proposition-approximate} から従う。
\item $B$ が $(A, I)$ 上 rig-エタールならば，結果は Lemma
\ref{restricted-lemma-approximate} から従う。
\item $I$ が単項ならば，結果は \cite[III Theorem 7]{Elkik} から従う。
\end{enumerate}
\end{remark}

\section{Rig-エタール代数}
\label{restricted-section-rig-etale}

\noindent
rig-smooth 代数の定義
（Definition \ref{restricted-definition-rig-smooth-homomorphism}）を踏まえれば，
次の定義は意外なものではない。

\begin{definition}
\label{restricted-definition-rig-etale-homomorphism}
$A$ を Noether 環とし，$I \subset A$ をイデアルとする。
$B$ を (\ref{restricted-equation-C-prime}) の対象とする。
ある整数 $c \geq 0$ が存在し，すべての $a \in I^c$ に対して
$\NL_{B/A}^\wedge$ 上の $a$ による乗法が $D(B)$ において零であるとき，
$B$ は {\it $(A, I)$ 上 rig-エタール} であるという。
\end{definition}

\noindent
次の補題の条件 (\ref{restricted-item-condition-artin}) は，
\cite{ArtinII} で形式修正を定義するために用いられた条件の一つである。
後でわれわれの条件と比較しやすくするため，条件の一覧にこれを加えた。

\begin{lemma}
\label{restricted-lemma-equivalent-with-artin}
$A$ を Noether 環とし，$I \subset A$ をイデアルとする。
$B$ を (\ref{restricted-equation-C-prime}) の対象とする。
$B = A[x_1, \ldots, x_r]^\wedge/J$ と書き
（Lemma \ref{restricted-lemma-topologically-finite-type-Noetherian}），
その素朴な余接複体を
$\NL_{B/A}^\wedge = (J/J^2 \to \bigoplus B\text{d}x_i)$
（\ref{restricted-equation-NL}）とする。次は同値である：
\begin{enumerate}
\item $B$ は $(A, I)$ 上 rig-エタールである。
\item
\label{restricted-item-zero-on-NL}
ある $c \geq 0$ が存在し，すべての $a \in I^c$ に対して
$\NL_{B/A}^\wedge$ 上の $a$ による乗法が $D(B)$ において零である。
\item
\label{restricted-item-zero-on-cohomology-NL}
% P10-E0653
% P10-E0654
ある $c \geq 0$ が存在し，$H^i(\NL_{B/A}^\wedge)$（$i = -1, 0$）は
$I^c$ により零化される。
\item
\label{restricted-item-zero-on-cohomology-NL-truncations}
ある $c \geq 0$ が存在し，すべての $n \geq 1$ に対して
$H^i(\NL_{B_n/A_n})$（$i = -1, 0$）は $I^c$ により零化される。
ここで $A_n = A/I^n$ および $B_n = B/I^nB$ である。
\item
\label{restricted-item-condition-artin-pre-pre}
すべての $a \in I$ に対して，次を満たす $c \geq 0$ が存在する：
\begin{enumerate}
\item $a^c$ は $H^0(\NL_{B/A}^\wedge)$ を零化する。
\item $a^c J \subset (f_1, \ldots, f_r) + J^2$ を満たす
$f_1, \ldots, f_r \in J$ が存在する。
\end{enumerate}
\item
\label{restricted-item-condition-artin-pre}
すべての $a \in I$ に対して，次を満たす
$f_1, \ldots, f_r \in J$ および $c \geq 0$ が存在する：
\begin{enumerate}
\item $\det_{1 \leq i, j \leq r}(\partial f_j/\partial x_i)$ は
$B$ において $a^c$ を割り切る。
\item $a^c J \subset (f_1, \ldots, f_r) + J^2$ である。
\end{enumerate}
\item
\label{restricted-item-condition-artin}
$J$ の生成元 $f_1, \ldots, f_t$ を選ぶと，次が成り立つ：
\begin{enumerate}
\item $A$ 上の $B$ の Jacobi イデアル，すなわち
% P10-E0395
行列 $(\partial f_j/\partial x_i)_{1 \leq i \leq r, 1 \leq j \leq t}$
の $r \times r$ 小行列式で生成される $B$ のイデアルは，
ある $c$ に対してイデアル $I^cB$ を含む。
\item $A$ 上の $B$ の Cramer イデアル，すなわち
$A[x_1, \ldots, x_r]^\wedge$-加群としての $J$ の第 $r$ Fitting
イデアルの $B$ における像で生成される $B$ のイデアルは，
ある $c$ に対して $I^cB$ を含む。
\end{enumerate}
\end{enumerate}
\end{lemma}

\begin{proof}
(1) と (\ref{restricted-item-zero-on-NL}) の同値性は Definition
\ref{restricted-definition-rig-etale-homomorphism} の言い換えである。

\medskip\noindent
(\ref{restricted-item-zero-on-NL}) と
(\ref{restricted-item-zero-on-cohomology-NL}) の同値性は
More on Algebra, Lemma
\ref{more-algebra-lemma-zero-in-derived} から従う。

\medskip\noindent
(\ref{restricted-item-zero-on-cohomology-NL}) と
(\ref{restricted-item-zero-on-cohomology-NL-truncations}) の同値性は，系
$\{\NL_{B_n/A_n}\}$ と
$\NL_{B/A}^\wedge \otimes_B B_n$ が狭義に同型であることから従う。
Lemma \ref{restricted-lemma-NL-is-limit} を参照せよ。いくつかの詳細は省略する。

\medskip\noindent
(\ref{restricted-item-zero-on-NL}) を仮定し，$a \in I$ とする。
$\NL_{B/A}^\wedge$ 上の $a^c$ による乗法が零となるように $c$ を取る。
More on Algebra, Lemma
\ref{more-algebra-lemma-map-out-of-almost-free} の (1) により，
$\text{d} \circ \alpha$ と $\alpha \circ \text{d}$ がともに
$a^c$ による乗法となる写像
$\alpha : \bigoplus B\text{d}x_i \to J/J^2$ が存在する。
$f_i \in J$ を，法 $J^2$ での類が
$\alpha(\text{d}x_i)$ に等しい元とする。
簡単な計算から (\ref{restricted-item-condition-artin-pre})(a), (b) が従う。

\medskip\noindent
(\ref{restricted-item-condition-artin-pre}) が
(\ref{restricted-item-condition-artin-pre-pre}) を含意することの確認は省略する。
これは $M \to B^{\oplus r}$ という形の $B$ 上の二項複体に関する
主張にすぎない。

\medskip\noindent
(\ref{restricted-item-condition-artin-pre-pre}) が成り立つと仮定する。
$I = (a_1, \ldots, a_t)$ とする。
$a_i^{c_i}$ に対して
(\ref{restricted-item-condition-artin-pre-pre})(a), (b) が成り立つような整数を
$c_i \geq 0$ とする。このとき $I^{\sum c_i}$ は
$H^0(\NL_{B/A}^\wedge)$ を零化する。
$f_{i, 1}, \ldots, f_{i, r} \in J$ を，$a_i$ に対する
(\ref{restricted-item-condition-artin-pre-pre})(b) のものとする。合成
$$
B^{\oplus r} \to J/J^2 \to \bigoplus B\text{d}x_i
$$
を考える。ここで第 $j$ 基底ベクトルは $J/J^2$ における
$f_{i, j}$ の類へ写される。
(\ref{restricted-item-condition-artin-pre-pre})(a), (b) により，
この合成の余核は $a_i^{2c_i}$ により零化される。
したがって，この写像は $a_i^{c_i}$ を可逆化した後に全射であり，
ゆえに同型である
（Algebra, Lemma \ref{algebra-lemma-fun}）。
したがって $B^{\oplus r} \to \bigoplus B\text{d}x_i$ の核は
$a_i$-冪捩れであり，ゆえに
$H^{-1}(\NL_{B/A}^\wedge) =
\Ker(J/J^2 \to \bigoplus B\text{d}x_i)$ は $a_i$-冪捩れである。
$B$ は Noether なので
（Lemma \ref{restricted-lemma-topologically-finite-type-Noetherian}），
$H^{-1}(\NL_{B/A}^\wedge)$ を含むすべての加群は有限である。
したがって，ある $d_i \geq 0$ に対して $a_i^{d_i}$ は
$H^{-1}(\NL_{B/A}^\wedge)$ を零化する。
このことから $I^{\sum d_i}$ は $H^{-1}(\NL_{B/A}^\wedge)$ を零化し，
(\ref{restricted-item-zero-on-cohomology-NL}) が成り立つことが分かる。

\medskip\noindent
以上により，条件
(\ref{restricted-item-zero-on-NL})，
(\ref{restricted-item-zero-on-cohomology-NL})，
(\ref{restricted-item-zero-on-cohomology-NL-truncations})，
(\ref{restricted-item-condition-artin-pre-pre})，および
(\ref{restricted-item-condition-artin-pre}) は同値である。
したがって，残るのはこれらの条件が
(\ref{restricted-item-condition-artin}) と同値であることを示すことである。
Cramer イデアル $\text{Fit}_r(J) B$ は
$J/J^2 = J \otimes_{A[x_1, \ldots, x_r]^\wedge} B$ なので
$\text{Fit}_r(J/J^2)$ に等しい。
More on Algebra, Lemma
\ref{more-algebra-lemma-fitting-ideal-basics} の (3) を参照せよ。
また，Jacobi イデアルはちょうど
$\text{Fit}_0(H^0(\NL_{B/A}^\wedge))$ である。
したがって，(\ref{restricted-item-zero-on-cohomology-NL}) と
(\ref{restricted-item-condition-artin}) の同値性は純代数的な問題であり，
次の段落でこれを論じる。

\medskip\noindent
$R$ を Noether 環とし，$I \subset R$ をイデアルとする。
$M \xrightarrow{d} R^{\oplus r}$ を二項複体とする。
次が同値であることを示さなければならない：
\begin{enumerate}
\item[(A)] $M \to R^{\oplus r}$ のコホモロジーは $I$ のある冪により
零化される。
\item[(B)] イデアル $\text{Fit}_r(M)$ および
$\text{Fit}_0(\text{Coker}(d))$ は $I$ のある冪を含む。
\end{enumerate}
$R$ は Noether なので，
% P10-E0655
部分 (2) は対応する閉部分スキームの包含として言い換えられる。
Algebra, Lemmas \ref{algebra-lemma-Zariski-topology} および
\ref{algebra-lemma-Noetherian-power} を参照せよ。
一方，$V(\text{Fit}_0(\Coker(d)))$ の補集合上では $d$ の余核は零であり，
$V(\text{Fit}_r(M))$ の補集合上では加群 $M$ は局所的に $r$ 個の元で
生成される。More on Algebra, Lemma
\ref{more-algebra-lemma-fitting-ideal-generate-locally} を参照せよ。
したがって (B) は次と同値である：
\begin{enumerate}
\item[(C)] $V(I)$ の外では $d$ の余核は零であり，加群 $M$ は
局所的に $\leq r$ 個の元で生成される。
\end{enumerate}
もちろんこれは，$M \to R^{\oplus r}$ のコホモロジーが
$\Spec(R) \setminus V(I)$ 上で消えるという条件と同値であり，
さらにこれは (A) と同値である。これで証明は完了する。
\end{proof}

\begin{lemma}
\label{restricted-lemma-rig-etale-rig-smooth}
$A$ を Noether 環とし，$I$ をイデアルとする。
$B$ を (\ref{restricted-equation-C-prime}) の対象とする。
$B$ が $(A, I)$ 上 rig-エタールならば，
$B$ は $(A, I)$ 上 rig-smooth である。
\end{lemma}

\begin{proof}
Definitions \ref{restricted-definition-rig-smooth-homomorphism} および
\ref{restricted-definition-rig-etale-homomorphism} から直ちに従う。
\end{proof}

\begin{lemma}
\label{restricted-lemma-rig-etale}
$A$ を Noether 環とし，$I$ をイデアルとする。
$B$ を有限型 $A$-代数とする。
\begin{enumerate}
\item $\Spec(B) \to \Spec(A)$ が
$\Spec(A) \setminus V(I)$ 上エタールならば，
$B^\wedge$ は Lemma \ref{restricted-lemma-equivalent-with-artin} の
同値な条件を満たす。
\item $B^\wedge$ が Lemma \ref{restricted-lemma-equivalent-with-artin} の
同値な条件を満たすならば，$\Spec(B_g)$ が
$\Spec(A) \setminus V(I)$ 上エタールとなる
$g \in 1 + IB$ が存在する。
\end{enumerate}
\end{lemma}

\begin{proof}
$B^\wedge$ が Lemma \ref{restricted-lemma-equivalent-with-artin} の
同値な条件を満たすと仮定する。
素朴な余接複体 $\NL_{B/A}$ は有限型 $B$-加群の複体なので，
$H^{-1}$ および $H^0$ は有限 $B$-加群である。
完備化は有限 $B$-加群上の完全函手であり
（Algebra, Lemma \ref{algebra-lemma-completion-flat}），
$\NL_{B^\wedge/A}^\wedge$ は複体 $\NL_{B/A}$ の完備化である
（これは表示を選べば容易に分かる）。
したがって仮定から，すべての $n$ に対して
$H^{-1}/I^nH^{-1}$ および $H^0/I^nH^0$ が $I^c$ により
零化されるような $c \geq 0$ が存在する。
Nakayama の補題（Algebra, Lemma \ref{algebra-lemma-NAK}）により，
これは $I^cH^{-1}$ と $I^cH^0$ が
$x \in IB$ によって $g = 1 + x$ と書ける元により零化されることを意味する。
$g$ を可逆化すると（これは商 $B/I^nB$ を変えない），
$\NL_{B/A}$ のコホモロジーは $I^c$ により零化される。
したがって，エタール環写像の定義により，$A \to B$ は
$V(I)$ の上にない $B$ の任意の素イデアルでエタールである。
Algebra, Definition \ref{algebra-definition-etale} を参照せよ。

\medskip\noindent
逆に，$\Spec(B) \to \Spec(A)$ が
$\Spec(A) \setminus V(I)$ 上エタールであると仮定する。
このとき，すべての $a \in I$ に対して，
% P10-E0656
$a^c$ による乗法が $\NL_{B/A}$ 上で零となるような
$c \geq 0$ が存在する。
$\NL_{B^\wedge/A}^\wedge$ は $\NL_{B/A}$ の導来完備化なので
（Lemma \ref{restricted-lemma-NL-is-limit} を参照せよ），
$B^\wedge$ は Lemma \ref{restricted-lemma-equivalent-with-artin} の
同値な条件を満たす。
\end{proof}

\begin{lemma}
\label{restricted-lemma-zero-after-modding-out}
$(A_1, I_1) \to (A_2, I_2)$ を，$A_1$ と $A_2$ が Noether である
Remark \ref{restricted-remark-base-change} のものとする。
$B_1$ を $(A_1, I_1)$ に対する
(\ref{restricted-equation-C-prime}) の対象とする。
$B_2$ を $B_1$ の基底変換とする。
$\NL^\wedge_{B_1/A_1}$ 上の $f_1 \in B_1$ による乗法が
$D(B_1)$ において零ならば，$\NL^\wedge_{B_2/A_2}$ 上の
像 $f_2 \in B_2$ による乗法は $D(B_2)$ において零である。
\end{lemma}

\begin{proof}
Lemma \ref{restricted-lemma-NL-base-change} により写像
% P10-E0657
$$
\NL_{B_1/A_1} \otimes_{B_2} B_1 \to \NL_{B_2/A_2}
$$
が存在し，これは $H^0$ 上で同型を，$H^{-1}$ 上で全射を誘導する。
% P10-E0384
したがって，More on Algebra, Lemma
\ref{more-algebra-lemma-two-term-surjection-map-zero} により結果が従う。
% P10-E0658
\end{proof}

\begin{lemma}
\label{restricted-lemma-base-change-rig-etale-homomorphism}
$A_1 \to A_2$ を Noether 環の写像とする。
$I_i \subset A_i$ を $V(I_1A_2) = V(I_2)$ を満たすイデアルとする。
$B_1$ を $(A_1, I_1)$ に対する
(\ref{restricted-equation-C-prime}) の対象とする。
$B_2$ を Remark \ref{restricted-remark-base-change} の意味での
$B_1$ の基底変換とする。
$B_1$ が $(A_1, I_1)$ 上 rig-エタールならば，
$B_2$ は $(A_2, I_2)$ 上 rig-エタールである。
\end{lemma}

\begin{proof}
Lemma \ref{restricted-lemma-zero-after-modding-out}，Definition
\ref{restricted-definition-rig-etale-homomorphism}，および $A_2$ は Noether なので
ある $c \geq 0$ に対して $I_2^c \subset I_1A_2$ となるという事実から従う。
\end{proof}

\begin{lemma}
\label{restricted-lemma-fully-faithful-etale-over-complement}
$A$ を Noether 環とし，$I \subset A$ をイデアルとする。
$B$ を，$\Spec(B) \to \Spec(A)$ が
$\Spec(A) \setminus V(I)$ 上エタールとなる有限型 $A$-代数とする。
$C$ を Noether $A$-代数とする。このとき，$I$-進完備化の任意の
$A$-代数写像 $B^\wedge \to C^\wedge$ は，一意な $A$-代数写像
$$
B \longrightarrow C^h
$$
から来る。ここで $C^h$ は More on Algebra, Lemma
\ref{more-algebra-lemma-henselization} における組 $(C, IC)$ のヘンゼル化である。
さらに，任意の $A$-代数準同型 $B \to C^h$ は，
$C/IC \to C'/IC'$ が同型となるあるエタール $C$-代数 $C'$ を経由する。
\end{lemma}

\begin{proof}
一意性は，例えば More on Algebra, Lemma
\ref{more-algebra-lemma-henselization-Noetherian-pair} により，
$C^h$ が $C^\wedge$ の部分環であることから従う。
最後の主張は，$C^h$ がこれらの $C$-代数 $C'$ の
フィルター付き余極限であることから従う。
More on Algebra, Lemma \ref{more-algebra-lemma-henselization}
の証明を参照せよ。以上を述べたので，存在性の証明に移る。

\medskip\noindent
$\varphi : B^\wedge \to C^\wedge$ を与えられた写像とする。
これは写像 $C \to B \otimes_A C$ の完備化の切断
$$
\sigma : (B \otimes_A C)^\wedge \longrightarrow C^\wedge
$$
を定める。$(A, I, B, C, \varphi)$ を
$(C, IC, B \otimes_A C, C, \sigma)$ で置き換えてよい。
このようにして $A = C$ と仮定してよいことが分かる。

\medskip\noindent
$A = C$ の場合の存在性を証明する。この場合，写像
$\varphi : B^\wedge \to A^\wedge$ は必ず全射である。
Lemmas \ref{restricted-lemma-rig-etale} および
\ref{restricted-lemma-exact-sequence-NL} により，
$\NL_{A^\wedge/\!_\varphi B^\wedge}^\wedge$ のコホモロジー群は
$I$ のある冪により零化される。
$\varphi$ は全射なので，これは
$\Ker(\varphi)/\Ker(\varphi)^2$ が $I$ のある冪により
零化されることを含意する。
したがって，$\varphi : B^\wedge \to A^\wedge$ は有限型
$B$-代数 $B \to D$ の完備化である。More on Algebra, Lemma
\ref{more-algebra-lemma-quotient-by-idempotent} を参照せよ。
ゆえに $A \to D$ は，同型 $A^\wedge \to D^\wedge$ を誘導する
有限型代数写像である。Lemma \ref{restricted-lemma-rig-etale} により，
$D$ を局所化で置き換えて，$A \to D$ が $V(I)$ の外で
エタールであると仮定してよい。
$A^\wedge \to D^\wedge$ は同型なので，
$\Spec(D) \to \Spec(A)$ は $V(ID)$ の近傍でもエタールである
（例えば More on Morphisms, Lemma
\ref{more-morphisms-lemma-check-smoothness-on-infinitesimal-nbhds} による）。
したがって $\Spec(D) \to \Spec(A)$ はエタールである。
ゆえに $D$ は $A^h$ へ写り，補題が証明された。
\end{proof}

\section{押し出し論法}
\label{restricted-section-approximation}

\noindent
この節の唯一の目標は，次の補題を証明することである。
この補題は rig-エタール代数の代数化において重要な役割を果たす。
証明には代数空間の理論を少し用いる。
本章の以前の版には，代数と若干の変形理論を用いた（はるかに長い）
証明があり，関心のある読者は Stacks project の履歴でそれを見つけられる。

\begin{lemma}
\label{restricted-lemma-lift-approximation}
$A$ を Noether 環，$I \subset A$ をイデアルとする。
$J \subset A$ を冪零イデアルとする。次の可換図式を考える。
$$
\xymatrix{
C \ar[r] & C_0 \ar@{=}[r] & C/JC \\
& B_0 \ar[u] \\
A \ar[r] \ar[uu] & A_0 \ar[u] \ar@{=}[r] & A/J
}
$$
この図式の縦の矢印はいずれも有限型であり，次を満たすとする。
\begin{enumerate}
\item $\Spec(C) \to \Spec(A)$ は $\Spec(A) \setminus V(I)$ 上エタールである。
\item $\Spec(B_0) \to \Spec(A_0)$ は
$\Spec(A_0) \setminus V(IA_0)$ 上エタールである。
\item $B_0 \to C_0$ はエタールであり，同型
$B_0/IB_0 = C_0/IC_0$ を誘導する。
\end{enumerate}
このとき，上の図式を次の可換図式へ補完できる。
$$
\xymatrix{
C \ar[r] & C/JC \\
B \ar[u] \ar[r] & B_0 \ar[u] \\
A \ar[r] \ar[u] & A/J \ar[u]
}
$$
ここで $A \to B$ は有限型，$B/JB = B_0$，$B \to C$ はエタールであり，
$\Spec(B) \to \Spec(A)$ は $\Spec(A) \setminus V(I)$ 上エタールである。
\end{lemma}

\begin{proof}
$X = \Spec(A)$，$X_0 = \Spec(A_0)$，$Y_0 = \Spec(B_0)$，
$Z = \Spec(C)$，$Z_0 = \Spec(C_0)$ と置く。さらに，
$U \subset X$，$U_0 \subset X_0$，$V_0 \subset Y_0$，
$W \subset Z$，$W_0 \subset Z_0$ を，それぞれ
$I$ の消滅集合の補集合と書く。状況を視覚化するための図は次のとおりである。
$$
\xymatrix{
Z \ar[dd] & Z_0 \ar[l] \ar[d] \\
& Y_0 \ar[d] \\
X & X_0 \ar[l]
}
\quad\quad\quad
\xymatrix{
W \ar[dd] & W_0 \ar[l] \ar[d] \\
& V_0 \ar[d] \\
U & U_0 \ar[l]
}
$$
補題の条件により，
$$
\xymatrix{
W_0 \ar[r] \ar[d] & Z_0 \ar[d] \\
V_0 \ar[r] & Y_0
}
$$
は初等識別正方形である。Derived Categories of Spaces, Definition
\ref{spaces-perfect-definition-elementary-distinguished-square}
を参照せよ。さらに，
$W_0 \to U_0$ および $V_0 \to U_0$ はエタールである。
$J$ は冪零と仮定したので，射 $X_0 \subset X$ は有限次数の肥厚である。
エタール・サイトの位相的不変性により，
$V_0 = V \times_X X_0$ を満たすスキームの一意なエタール射
$V \to X$ が存在し，与えられた射 $W_0 \to V_0$ を
$X$ 上の一意な射 $W \to V$ へ持ち上げられる。
\'Etale Morphisms, Theorem
\ref{etale-theorem-remarkable-equivalence}
を参照せよ。
$W_0 \to V_0$ は分離なので，射 $W \to V$ も分離である。
例えば More on Morphisms, Lemma
\ref{more-morphisms-lemma-deform-property} を参照せよ。
Pushouts of Spaces, Lemma
\ref{spaces-pushouts-lemma-construct-elementary-distinguished-square}
により，$X$ 上の代数空間の圏において初等識別正方形
$$
\xymatrix{
W \ar[r] \ar[d] & Z \ar[d] \\
V \ar[r] & Y
}
$$
を構成できる。初等識別正方形の基底変換は初等識別正方形なので
（Derived Categories of Spaces, Lemma
\ref{spaces-perfect-lemma-make-more-elementary-distinguished-squares}），
$$
\xymatrix{
W_0 \ar[r] \ar[d] & Z_0 \ar[d] \\
V_0 \ar[r] & Y \times_X X_0
}
$$
は初等識別正方形である。初等識別正方形は押し出しなので
（Pushouts of Spaces, Lemma
\ref{spaces-pushouts-lemma-elementary-distinguished-square-pushout}），
これらの空間を含む二つの正方形と両立する一意な同型
$Y \times_X X_0 = Y_0$ が存在する。
Limits of Spaces, Proposition \ref{spaces-limits-proposition-affine}
により，$Y$ はアフィンである。$Y = \Spec(B)$ と書く。
$B$ が所望の図式に入り，要求されたすべての性質を満たすことは明らかである。
\end{proof}

\section{rig-エタール代数の代数化}
\label{restricted-section-approximation-principal}

\noindent
主な目標は，基礎となる Noether 環 $A$ が G-環であるとは仮定せず，
かつイデアル $I \subset A$ が任意で，必ずしも単項とは限らない場合に，
rig-エタール代数の代数化を証明することである。
まず単項イデアルの場合を証明し，次いで
Section \ref{restricted-section-approximation} の結果を用いて証明を完成させる。

\begin{lemma}
\label{restricted-lemma-approximate-principal}
\begin{reference}
\cite[III Theorem 7]{Elkik} の rig-エタールの場合
\end{reference}
$A$ を Noether 環，$I = (a)$ を単項イデアルとする。
$B$ を (\ref{restricted-equation-C-prime}) の対象で，
$(A, I)$ 上 rig-エタールであるものとする。
このとき，有限型 $A$-代数 $C$ と同型
$B \cong C^\wedge$ が存在する。
\end{lemma}

\begin{proof}
表示 $B = A[x_1, \ldots, x_r]^\wedge/J$ を選ぶ。
Lemma \ref{restricted-lemma-equivalent-with-artin} の (6) により，
$c \geq 0$ と $f_1, \ldots, f_r \in J$ であって，
$\det_{1 \leq i, j \leq r}(\partial f_j/\partial x_i)$ が
$B$ において $a^c$ を割り切り，かつ
$a^c J \subset (f_1, \ldots, f_r) + J^2$
となるものを見つけられる。
したがって Lemma \ref{restricted-lemma-approximate-presentation-rig-smooth}
を適用できる。これで証明は終わるが，この場合には
Lemma \ref{restricted-lemma-get-morphism-general-better}
の使用を，はるかに容易な
Lemma \ref{restricted-lemma-get-morphism-principal}
で置き換えられることを指摘しておく。
\end{proof}

\begin{lemma}
\label{restricted-lemma-approximate}
$A$ を Noether 環とする。$I \subset A$ をイデアルとする。
$B$ を (\ref{restricted-equation-C-prime}) の対象で，
$(A, I)$ 上 rig-エタールであるものとする。
このとき，有限型 $A$-代数 $C$ と同型
$B \cong C^\wedge$ が存在する。
\end{lemma}

\begin{proof}
$I$ の生成元の個数に関する帰納法で補題を証明する。
$I = (a_1, \ldots, a_t)$ と書く。$t = 0$ なら $I = 0$ なので，
証明すべきことはない。$t = 1$ なら，補題は
Lemma \ref{restricted-lemma-approximate-principal} から従う。$t > 1$ と仮定する。

\medskip\noindent
任意の $m \geq 1$ に対し $\bar A_m = A/(a_t^m)$ と置く。
$\bar A_m$ のイデアル
$\bar I_m = (\bar a_1, \ldots, \bar a_{t - 1})$ を考える。
$V(I \bar A_m) = V(\bar I_m)$ であることに注意せよ。
$B_m = B/(a_t^m)$ を，写像
$(A, I) \to (\bar A_m, \bar I_m)$ に関する $B$ の基底変換とする。
Remark \ref{restricted-remark-take-bar} を参照せよ。
Lemma \ref{restricted-lemma-base-change-rig-etale-homomorphism}
により，$B_m$ は $(\bar A_m, \bar I_m)$ 上 rig-エタールである。

\medskip\noindent
$t$ に関する帰納法の仮定により，有限型
$\bar A_m$-代数 $C_m$ と写像 $C_m \to B_m$ であって，
$\bar I_m$ に関する完備化が同型
$C_m^\wedge \cong B_m$ を誘導するものを見つけられる。
Lemma \ref{restricted-lemma-rig-etale} により，
$\Spec(C_m) \to \Spec(\bar A_m)$ が
$\Spec(\bar A_m) \setminus V(\bar I_m)$ 上エタールであると
仮定してよい。

\medskip\noindent
% P10-E0659
さらに，前段落のような $A_m \to C_m \to B_m$ を，
$C_m/(a_t^{m - 1}) \to C_{m - 1}$ が，与えられた
$A$-代数構造および
$B_{m - 1} = B_m/(a_t^{m - 1})$ への写像と両立する同型となるように
選べることを主張する。まず
$A_1 \to C_1 \to B_1$ の選択を一つ固定する。
所望の性質をもつ
$C_{m - 1} \to C_{m - 2} \to \ldots \to C_1$
が得られたと仮定する。
$C_m/(a_t^{m - 1})$ は
$\Spec(\bar A_{m - 1}) \setminus V(\bar I_{m - 1})$
上エタールであることに注意せよ。
従って Lemma \ref{restricted-lemma-fully-faithful-etale-over-complement}
により，$\bar I_{m - 1}$ を法として同型を誘導するエタール拡大
$C_{m - 1} \to C'_{m - 1}$ と，
完備化上で同型
$B_m/(a_t^{m - 1}) \to B_{m - 1}$ を誘導する
$\bar A_{m - 1}$-代数写像
$C_m/(a_t^{m - 1}) \to C'_{m - 1}$ が存在する。
$C_m/(a_t^{m - 1}) \to C'_{m - 1}$ は，
Morphisms, Lemma \ref{morphisms-lemma-etale-permanence}
により $V(\bar I_{m - 1})$ の補集合上エタールであり，
$V(\bar I_{m - 1})$ 上では完備化上の同型を誘導するので，
そこでもエタールである（例えば More on Morphisms, Lemma
\ref{more-morphisms-lemma-check-smoothness-on-infinitesimal-nbhds}
による）。従って
$C_m/(a_t^{m - 1}) \to C'_{m - 1}$ はエタールである。
エタール射の位相的不変性
（\'Etale Morphisms, Theorem
\ref{etale-theorem-remarkable-equivalence}）により，
エタール環写像 $C_m \to C'_m$ であって，
$C_m/(a_t^{m - 1}) \to C'_{m - 1}$ が
$C_m/(a_t^{m - 1}) \to C'_m/(a_t^{m - 1})$
と同型になるものが存在する。
$C'_m$ の $\bar I_m$-進完備化は $C_m$ の
$\bar I_m$-進完備化，すなわち $B_m$ に等しいことに注意せよ
（細部は省略する）。図式
$$
\xymatrix{
 & C'_m \ar[r] & C'_m/(a_t^{m - 1}) \\
C''_m \ar@{..>}[ru] \ar@{..>}[rr] & & C_{m - 1} \ar[u] \\
 & \bar A_m \ar[r] \ar[uu] \ar@{..>}[lu] & \bar A_{m - 1} \ar[u]
}
$$
に Lemma \ref{restricted-lemma-lift-approximation} を適用する。
% P10-E0660
これにより，$C_{m - 1}$ を $\bar A_m$ 上の代数へ持ち上げた
「持ち上げ」$C''_m$ で，所望の性質をすべてもつものが存在する。

\medskip\noindent
構成により，$(C_m)$ は単項イデアル $(a_t)$ に対する圏
(\ref{restricted-equation-C}) の対象である。
従って逆極限 $B' = \lim C_m$ は $(a_t)$-進完備な
$A$-代数で，$B'/a_t B'$ は $A/(a_t)$ 上有限型である。
Lemma \ref{restricted-lemma-topologically-finite-type} を参照せよ。
構成により $C_m = B'/(a_t^m)$ であり，$C_m$ の
$I$-進完備化は $B_m$，$B'$ の $I$-進完備化は $B$ と同型である。
% P10-E0659
% P10-E0662
各 $m$ に対し，複体 $\NL_{C_m/A_m}$ は，構成により
$V(IC_m) \subset \Spec(C_m)$ に台をもつ有限生成コホモロジー加群を
もつ。従って，これらの加群は $I$ のある冪で零化されるので，
$I$-進完備である。$\NL^\wedge_{B_m/A_m}$ は
$\NL_{C_m/A_m}$ の $I$-進完備化であるから
（Lemma \ref{restricted-lemma-NL-is-completion} を参照），
写像 $\NL_{C_m/A_m} \to \NL_{B_m/A_m}$ は
コホモロジー上の同型を誘導する。
% P10-E0661
従って，$B$ が $A$ 上 rig-エタールであることから，
すべての $m$ に対して $\NL_{C_m/A_m}$ のコホモロジーを零化する
$I$ の，添字に依存しないある冪が存在する。
Lemma \ref{restricted-lemma-equivalent-with-artin} を参照せよ。
同じ補題により，$B'$ は $(A, (a_t))$ 上 rig-エタールである。
最後に，$(A, (a_t))$ 上の $B'$ に
Lemma \ref{restricted-lemma-approximate-principal} を適用すれば証明が終わる。
\end{proof}

\begin{lemma}
\label{restricted-lemma-approximate-by-etale-over-complement}
$A$ を Noether 環とする。$I \subset A$ をイデアルとする。
$B$ を $I$-進完備な $A$-代数で，
$A/I \to B/IB$ が有限型であるものとする。
Lemma \ref{restricted-lemma-equivalent-with-artin} の同値な条件は，
次の条件とも同値である。
\begin{enumerate}
\item[(8)]
\label{restricted-item-algebraize}
有限型 $A$-代数 $C$ であって，
$\Spec(C) \to \Spec(A)$ が $\Spec(A) \setminus V(I)$
上エタールであり，かつ $B \cong C^\wedge$ となるものが存在する。
\end{enumerate}
\end{lemma}

\begin{proof}
Lemmas \ref{restricted-lemma-equivalent-with-artin}, \ref{restricted-lemma-approximate}, and
\ref{restricted-lemma-rig-etale} を組み合わせればよい。細部を一つ省略する。
\end{proof}

\section{有限型射}
\label{restricted-section-finite-type}

\noindent
Formal Spaces, Section \ref{formal-spaces-section-finite-type}
では，形式代数空間の有限型射を定義した。
この節では，有限生成の定義イデアルをもつ adic 位相環の，
それに対応する型の連続環写像を調べる。
読者にはこの節を読み飛ばすことを強く勧める。

\begin{lemma}
\label{restricted-lemma-finite-type}
$A$ と $B$ を，有限生成の定義イデアルをもつ adic 位相環とする。
$\varphi : A \to B$ を連続環準同型とする。次は同値である。
\begin{enumerate}
\item $\varphi$ は adic であり，$B$ は $A$ 上位相的有限型である。
\item $\varphi$ は taut であり，$B$ は $A$ 上位相的有限型である。
\item 定義イデアル $I \subset A$ であって，$B$ 上の位相が
$I$-進位相となるものが存在し，かつ定義イデアル $I' \subset A$ であって
$A/I' \to B/I'B$ が有限型となるものが存在する。
\item すべての定義イデアル $I \subset A$ に対して，$B$ 上の位相は
$I$-進位相であり，$A/I \to B/IB$ は有限型である。
\item 定義イデアル $I \subset A$ であって，$B$ 上の位相が
$I$-進位相となり，$B$ が圏
(\ref{restricted-equation-C-prime}) に属するものが存在する。
\item すべての定義イデアル $I \subset A$ に対して，$B$ 上の位相は
$I$-進位相であり，$B$ は圏 (\ref{restricted-equation-C-prime}) に属する。
\item 位相的 $A$-代数として，$B$ は
$A\{x_1, \ldots, x_r\}$ の閉イデアルによる商である。
\item 位相的 $A$-代数として，$B$ は
$A[x_1, \ldots, x_r]^\wedge$ の閉イデアルによる商である。
ここで $A[x_1, \ldots, x_r]^\wedge$ は，
$A$ のある定義イデアルに関する $A[x_1, \ldots, x_r]$ の完備化である。
% P10-E0663
\item ここにさらに追加する。
\end{enumerate}
さらに，これらの同値な条件は，Formal Spaces, Remark
\ref{formal-spaces-remark-variant-adic-star}
で定義された $\text{WAdm}^{adic*}$ の射の局所的性質を定める。
\end{lemma}

\begin{proof}
taut 環準同型は Formal Spaces, Definition
\ref{formal-spaces-definition-taut} で定義されている。
adic 環準同型は Formal Spaces, Definition
\ref{formal-spaces-definition-adic-homomorphism} で定義されている。
補題は Formal Spaces, Lemmas
\ref{formal-spaces-lemma-quotient-restricted-power-series},
\ref{formal-spaces-lemma-quotient-restricted-power-series-admissible}, and
\ref{formal-spaces-lemma-adic-homomorphism}
を組み合わせれば従う。細部は省略する。
% P10-E0664
念のため述べると，位相環
$A[x_1, \ldots, x_n]^\wedge$ と $A\{x_1, \ldots, x_r\}$
の間に相違はない。Formal Spaces, Remark
\ref{formal-spaces-remark-I-adic-completion-and-restricted-power-series}
を参照せよ。
\end{proof}

\begin{remark}
\label{restricted-remark-NL-well-defined-topological}
$A \to B$ を $\text{WAdm}^{adic*}$ の射で，adic かつ位相的有限型な
ものとする（Lemma \ref{restricted-lemma-finite-type} を参照）。
$B = A\{x_1, \ldots, x_r\}/J$ と書く。このとき，次のように置ける
\footnote{実際，この構成は Formal Spaces, Lemma
\ref{formal-spaces-lemma-quotient-restricted-power-series}
の同値な条件を満たす $\text{WAdm}^{count}$ の射に対して成り立つ。}
$$
\NL_{B/A}^\wedge = \left(J/J^2 \longrightarrow \bigoplus B\text{d}x_i\right)
$$
Lemma \ref{restricted-lemma-NL-up-to-homotopy} の証明とまったく同様に，
読者はこの $B$-加群の複体がホモトピー圏 $K(B)$ において
（一意な同型を除いて）良定義であることを示せる。
さて，$A$ が Noether 環で $I \subset A$ が定義イデアルなら，
この構成は Equation (\ref{restricted-equation-NL}) で，
Section \ref{restricted-section-naive-cotangent-complex} において定義した
$(A, I)$ 上の $B$ の素朴な余接複体を再現する。これは単に
% P10-E0664
$A[x_1, \ldots, x_n]^\wedge$ が $A\{x_1, \ldots, x_r\}$
と一致するからである。Formal Spaces, Remark
\ref{formal-spaces-remark-I-adic-completion-and-restricted-power-series}
を参照せよ。
特に，依然として $A$ が adic Noether 位相環であるとき，
対象 $\NL_{B/A}^\wedge$ は定義イデアル $I \subset A$
の選択に依存しないことが分かる。
\end{remark}

\begin{lemma}
\label{restricted-lemma-base-change-finite-type}
Lemma \ref{restricted-lemma-finite-type} で定義された
$\textit{WAdm}^{adic*}$ の射の性質 $P$ を考える。
このとき $P$ は，Formal Spaces, Remark
\ref{formal-spaces-remark-base-change-variant-adic-star}
で定義された基底変換で安定である。
\end{lemma}

\begin{proof}
主張が意味をもつことは Lemma \ref{restricted-lemma-finite-type} による。
これが真であることを見るため，$\textit{WAdm}^{adic*}$ における射
$B \to A$ と $B \to C$ があり，ある閉イデアル $J$ に対して
位相的 $B$-代数として
$A = B\{x_1, \ldots, x_r\}/J$
であると仮定する。
% P10-E0665
すると $A \widehat{\otimes}_B C$ は
$C\{x_1, \ldots, x_r\}/J'$ と同型である。ここで
$J'$ は $JC\{x_1, \ldots, x_r\}$ の閉包である。
若干の細部は省略する。
\end{proof}

\begin{lemma}
\label{restricted-lemma-composition-finite-type}
Lemma \ref{restricted-lemma-finite-type} で定義された
$\textit{WAdm}^{adic*}$ の射の性質 $P$ を考える。
このとき $P$ は，Formal Spaces, Remark
\ref{formal-spaces-remark-composition-variant-adic-star}
で定義された合成で安定である。
\end{lemma}

\begin{proof}
主張が意味をもつことは Lemma \ref{restricted-lemma-finite-type} による。
これを証明する最も容易な方法は，(a) adic 位相環の間の
adic 環写像の合成は adic であり，(b) 連続環写像の合成は
% P10-E0666
位相的有限型であるという性質を保つことを示すことである。
細部は省略する。
\end{proof}

\noindent
次の補題は，adic* 形式代数空間の射が局所有限型であることと，
それがエタール局所的に Lemma \ref{restricted-lemma-finite-type}
で述べた型の位相環写像により与えられることが同値である，と述べる。

\begin{lemma}
\label{restricted-lemma-finite-type-morphisms}
$S$ をスキームとする。$f : X \to Y$ を，$S$ 上の
局所 adic* 形式代数空間の射とする。次は同値である。
\begin{enumerate}
\item 任意の可換図式
$$
\xymatrix{
U \ar[d] \ar[r] & V \ar[d] \\
X \ar[r] & Y
}
$$
で，$U$ と $V$ がアフィン形式代数空間であり，$U \to X$ と
$V \to Y$ が代数空間により表現可能かつエタールであるものに対し，
射 $U \to V$ は，adic かつ位相的有限型な
$\textit{WAdm}^{adic*}$ の射に対応する。
\item Formal Spaces, Definition
\ref{formal-spaces-definition-formal-algebraic-space}
におけるような被覆 $\{Y_j \to Y\}$ が存在し，各 $j$ に対して
Formal Spaces, Definition
\ref{formal-spaces-definition-formal-algebraic-space}
におけるような被覆
$\{X_{ji} \to Y_j \times_Y X\}$ が存在して，各
$X_{ji} \to Y_j$ が adic かつ位相的有限型な
$\textit{WAdm}^{adic*}$ の射に対応する。
\item Formal Spaces, Definition
\ref{formal-spaces-definition-formal-algebraic-space}
におけるような被覆 $\{X_i \to X\}$ が存在し，各 $i$ に対して
分解 $X_i \to Y_i \to Y$ が存在する。ここで $Y_i$ は
アフィン形式代数空間，$Y_i \to Y$ は代数空間により表現可能かつ
エタールであり，$X_i \to Y_i$ は adic かつ位相的有限型な
$\textit{WAdm}^{adic*}$ の射に対応する。
\item $f$ は局所有限型である。
\end{enumerate}
\end{lemma}

\begin{proof}
Lemma \ref{restricted-lemma-finite-type} の (1) と (2) の同値性，および
Formal Spaces, Lemma
\ref{formal-spaces-lemma-finite-type-local-property}
から直ちに従う。
\end{proof}

\section{被約化上の有限型性}
\label{restricted-section-finite-type-red}

\noindent
この節では，局所可算添字付き形式代数空間の射
$X \to Y$ であって，$X_{red} \to Y_{red}$ が局所有限型であるものに
ついて少し論じる。これを代数的条件へ翻訳する。この代数的条件を
理解するには，次のことを念頭に置くとよい。
\begin{itemize}
\item $A$ が弱許容位相環なら，位相的冪零元の集合
$\mathfrak a \subset A$ は開な根基イデアルであり，
$\text{Spf}(A)_{red} = \Spec(A/\mathfrak a)$ である。
\end{itemize}
Formal Spaces, Definition
\ref{formal-spaces-definition-weakly-admissible},
Lemma \ref{formal-spaces-lemma-topologically-nilpotent}, and
Example \ref{formal-spaces-example-reduction-affine-formal-spectrum}
を参照せよ。

\begin{lemma}
\label{restricted-lemma-finite-type-red}
$\text{WAdm}^{count}$ の射 $\varphi : A \to B$ に対し，性質
$P(\varphi)=$「誘導された環準同型
$A/\mathfrak a \to B/\mathfrak b$ が有限型である」を考える。
ここで $\mathfrak a \subset A$ と $\mathfrak b \subset B$ は
位相的冪零元からなるイデアルである。
このとき $P$ は，Formal Spaces, Situation
\ref{formal-spaces-situation-local-property}
で定義された局所的性質である。
\end{lemma}

\begin{proof}
Formal Spaces, Situation
\ref{formal-spaces-situation-local-property}
におけるような可換図式
$$
\xymatrix{
B \ar[r] & (B')^\wedge \\
A \ar[r] \ar[u]^\varphi & (A')^\wedge \ar[u]_{\varphi'}
}
$$
を考える。この図式に $\text{Spf}$ を取ると，横の矢印が
代数空間により表現可能かつエタールであるような
アフィン形式代数空間の図式
$$
\xymatrix{
\text{Spf}(B) \ar[d] &
\text{Spf}((B')^\wedge) \ar[l] \ar[d] \\
\text{Spf}(A) &
\text{Spf}((A')^\wedge) \ar[l]
}
$$
を得る。これは Formal Spaces, Lemma
\ref{formal-spaces-lemma-etale} による。
従って，横の矢印がエタールであるようなアフィン・スキームの可換図式
$$
\xymatrix{
\text{Spf}(B)_{red} \ar[d]^f &
\text{Spf}((B')^\wedge)_{red} \ar[l]^g \ar[d]^{f'} \\
\text{Spf}(A)_{red} &
\text{Spf}((A')^\wedge)_{red} \ar[l]
}
$$
を得る。これは Formal Spaces, Lemma
\ref{formal-spaces-lemma-reduction-smooth} による。
Formal Spaces, Example
\ref{formal-spaces-example-reduction-affine-formal-spectrum} および
Lemma \ref{formal-spaces-lemma-etale-surjective}
により，Formal Spaces, Situation
\ref{formal-spaces-situation-local-property}
の条件 (1), (2), (3) は，次の主張へ翻訳される。
\begin{enumerate}
\item $f$ が局所有限型なら，$f'$ は局所有限型である。
\item $f'$ が局所有限型で $g$ が全射なら，$f$ は局所有限型である。
\item $T_i \to S$, $i = 1, \ldots, n$ が局所有限型なら，
$\coprod_{i = 1, \ldots, n} T_i \to S$ は局所有限型である。
\end{enumerate}
性質 (1) と (2) は，局所有限型であることがエタール位相において
始域と終域の双方で局所的であることから従う。
Morphisms of Spaces, Section
\ref{spaces-morphisms-section-finite-type}
の議論を参照せよ。性質 (3) は定義から直ちに従う。
\end{proof}

\begin{lemma}
\label{restricted-lemma-base-change-finite-type-red}
Lemma \ref{restricted-lemma-finite-type-red} で定義された
$\textit{WAdm}^{count}$ の射の性質 $P$ を考える。
このとき $P$ は基底変換で安定である
（Formal Spaces, Situation
\ref{formal-spaces-situation-base-change-local-property}）。
\end{lemma}

\begin{proof}
主張が意味をもつことは Lemma \ref{restricted-lemma-finite-type-red} による。
これが真であることを見るため，$\textit{WAdm}^{count}$ における射
$B \to A$ と $B \to C$ があり，
$B/\mathfrak b \to A/\mathfrak a$ が有限型であると仮定する。
ここで $\mathfrak b \subset B$ と $\mathfrak a \subset A$ は
位相的冪零元からなるイデアルである。
$A$ と $B$ は弱許容なので，イデアル
$\mathfrak a$ と $\mathfrak b$ は開である。
$\mathfrak c \subset C$ を位相的冪零元からなる（開）イデアルとする。
すると全射
$A \widehat{\otimes}_B C \to
A/\mathfrak a \otimes_{B/\mathfrak b} C/\mathfrak c$
が得られ，その核は弱定義イデアルなので位相的冪零元のみからなる
（Formal Spaces, Lemma
\ref{formal-spaces-lemma-completed-tensor-product}
の証明と比較せよ）。すでに
$C/\mathfrak c \to A/\mathfrak a \otimes_{B/\mathfrak b} C/\mathfrak c$
は $B/\mathfrak b \to A/\mathfrak a$ の基底変換として有限型なので，
結論を得る。
\end{proof}

\begin{lemma}
\label{restricted-lemma-composition-finite-type-red}
Lemma \ref{restricted-lemma-finite-type-red} で定義された
$\textit{WAdm}^{count}$ の射の性質 $P$ を考える。
このとき $P$ は合成で安定である
（Formal Spaces, Situation
\ref{formal-spaces-situation-composition-local-property}）。
\end{lemma}

\begin{proof}
省略する。ヒント：有限型環写像の合成は有限型である。
\end{proof}

\begin{lemma}
\label{restricted-lemma-finite-type-finite-type-red}
$\varphi : A \to B$ を $\textit{WAdm}^{count}$ の射とする。
$\varphi$ が taut かつ位相的有限型なら，$\varphi$ は
Lemma \ref{restricted-lemma-finite-type-red} で定義された条件を満たす。
\end{lemma}

\begin{proof}
これは定義から容易に従う。
\end{proof}

\begin{lemma}
\label{restricted-lemma-Noetherian-finite-type-red}
$\varphi : A \to B$ を $\textit{WAdm}^{Noeth}$ の射で，
Lemma \ref{restricted-lemma-finite-type-red} で定義された条件を満たすものとする。
このとき $A \to B$ は位相的有限型である。
\end{lemma}

\begin{proof}
$\mathfrak b \subset B$ を位相的冪零元からなるイデアルとする。
$B/\mathfrak b$ の $A$ 上の生成元へ写る
$b_1, \ldots, b_r \in B$ を選ぶ。
イデアル $\mathfrak b$ の生成元
$b_{r + 1}, \ldots, b_s$ を選ぶ。写像
$$
\varphi : A[x_1, \ldots, x_s] \longrightarrow B, \quad
x_i \longmapsto b_i
$$
の像が稠密であることを主張する。実際，ある $n \geq 0$ に対して
$b \in \mathfrak b^n$ なら，多重指数
$E = (e_{r + 1}, \ldots, e_s)$ で
$|E| = \sum e_i = n$ となるものと $b_E \in B$ を用いて，
$b = \sum b_E b_{r + 1}^{e_{r + 1}} \ldots b_s^{e_s}$
と書ける。次に，
$b_E = f_E(b_1, \ldots, b_r) + b'_E$
と書ける。ここで $b'_E \in \mathfrak b$ かつ
$f_E \in A[x_1, \ldots, x_r]$ である。
合わせると
$b \in \Im(\varphi) + \mathfrak b^{n + 1}$
を得る。帰納法により，すべての $n \geq 0$ に対して
$B = \Im(\varphi) + \mathfrak b^n$ である。
% P10-E0396
$\mathfrak b$ は $B$ の定義イデアルなので，これは所望の主張を含意する。
\end{proof}

\begin{lemma}
\label{restricted-lemma-Noetherian-adic-finite-type-red}
$\varphi : A \to B$ を $\textit{WAdm}^{Noeth}$ の射とする。
$\varphi$ が adic なら，次は同値である。
\begin{enumerate}
\item $\varphi$ は Lemma \ref{restricted-lemma-finite-type-red}
で定義された条件を満たす。
\item $\varphi$ は Lemma \ref{restricted-lemma-finite-type}
で定義された条件を満たす。
\end{enumerate}
\end{lemma}

\begin{proof}
省略する。ヒント：(1) $\Rightarrow$ (2) の証明には
Lemma \ref{restricted-lemma-Noetherian-finite-type-red} を用いる。
\end{proof}

\begin{lemma}
\label{restricted-lemma-finite-type-red-morphisms}
$S$ をスキームとする。$f : X \to Y$ を，$S$ 上の
局所可算添字付き形式代数空間の射とする。次は同値である。
\begin{enumerate}
\item 任意の可換図式
$$
\xymatrix{
U \ar[d] \ar[r] & V \ar[d] \\
X \ar[r] & Y
}
$$
で，$U$ と $V$ がアフィン形式代数空間であり，$U \to X$ と
$V \to Y$ が代数空間により表現可能かつエタールであるものに対し，
射 $U \to V$ は，Lemma \ref{restricted-lemma-finite-type-red}
で定義された性質を満たす $\textit{WAdm}^{count}$ の射に対応する。
\item Formal Spaces, Definition
\ref{formal-spaces-definition-formal-algebraic-space}
におけるような被覆 $\{Y_j \to Y\}$ が存在し，各 $j$ に対して
Formal Spaces, Definition
\ref{formal-spaces-definition-formal-algebraic-space}
におけるような被覆
$\{X_{ji} \to Y_j \times_Y X\}$ が存在して，各
$X_{ji} \to Y_j$ が，Lemma \ref{restricted-lemma-finite-type-red}
で定義された性質を満たす $\textit{WAdm}^{count}$ の射に対応する。
\item Formal Spaces, Definition
\ref{formal-spaces-definition-formal-algebraic-space}
におけるような被覆 $\{X_i \to X\}$ が存在し，各 $i$ に対して
分解 $X_i \to Y_i \to Y$ が存在する。ここで $Y_i$ は
アフィン形式代数空間，$Y_i \to Y$ は代数空間により表現可能かつ
エタールであり，$X_i \to Y_i$ は，Lemma
\ref{restricted-lemma-finite-type-red} で定義された性質を満たす
$\textit{WAdm}^{count}$ の射に対応する。
\item 射 $f_{red} : X_{red} \to Y_{red}$ は局所有限型である。
\end{enumerate}
\end{lemma}

\begin{proof}
(1), (2), (3) の同値性は，Lemma \ref{restricted-lemma-finite-type-red}
および Formal Spaces, Lemma
\ref{formal-spaces-lemma-property-defines-property-morphisms}
の適用から従う。(2) における $Y_j$ と $X_{ji}$ を取る。このとき，
\begin{itemize}
\item 族 $\{Y_{j, red} \to Y_{red}\}$ と
$\{X_{ji, red} \to X_{red}\}$ は，
アフィン・スキームによるエタール被覆である。
これは Formal Spaces, Lemma
\ref{formal-spaces-lemma-reduction-formal-algebraic-space}
の証明の議論から，または直接 Formal Spaces, Lemma
\ref{formal-spaces-lemma-reduction-smooth} から従う。
\item $X_{ji} \to Y_j$ が $\textit{WAdm}^{count}$ の射
$B_j \to A_{ji}$ に対応するなら，
$X_{ji, red} \to Y_{j, red}$ は環写像
$B_j/\mathfrak b_j \to A_{ji}/\mathfrak a_{ji}$ に対応する。
ここで $\mathfrak b_j$ と $\mathfrak a_{ji}$ は
位相的冪零元からなるイデアルである。これは Formal Spaces, Example
\ref{formal-spaces-example-reduction-affine-formal-spectrum}
から従う。従って $X_{ji, red} \to Y_{j, red}$ が局所有限型であることと，
$B_j \to A_{ji}$ が Lemma \ref{restricted-lemma-finite-type-red}
で定義された性質を満たすことは同値である。
\end{itemize}
これらの注意により，(2) と (4) の同値性が従う。実際，
局所有限型であることは，代数空間の射の性質であって，
始域と終域の双方でエタール局所的である。
Morphisms of Spaces, Section
\ref{spaces-morphisms-section-finite-type}
の議論を参照せよ。
\end{proof}

\section{平坦射}
\label{restricted-section-flat}

\noindent
この節では，局所 Noether 形式代数空間の平坦射を定義する。

\begin{lemma}
\label{restricted-lemma-flat-axioms}
$\textit{WAdm}^{Noeth}$ の射上の性質
$P(\varphi)=$「$\varphi$ は平坦である」は，Formal Spaces, Remark
\ref{formal-spaces-remark-variant-Noetherian}
で定義された局所的性質である。
\end{lemma}

\begin{proof}
主張の意味を思い出そう。まず，$\textit{WAdm}^{Noeth}$ は，
対象が adic Noether 位相環，射が連続環準同型である圏である。
次の可換図式を考える。
$$
\xymatrix{
B \ar[r] & (B')^\wedge \\
A \ar[r] \ar[u]^\varphi & (A')^\wedge \ar[u]_{\varphi'}
}
$$
これは次の条件を満たすとする。$A$ と $B$ は adic Noether 位相環，
$A \to A'$ と $B \to B'$ はエタール環写像，
ある定義イデアル $I \subset A$ に対して
$(A')^\wedge = \lim A'/I^nA'$，
ある定義イデアル $J \subset B$ に対して
$(B')^\wedge = \lim B'/J^nB'$ であり，
$\varphi : A \to B$ と
$\varphi' : (A')^\wedge \to (B')^\wedge$ は連続である。
Formal Spaces, Lemma
\ref{formal-spaces-lemma-completion-in-sub}
により，$(A')^\wedge$ と $(B')^\wedge$ は
adic Noether 位相環であることに注意せよ。示すべきことは次である。
\begin{enumerate}
\item $\varphi$ が平坦 $\Rightarrow \varphi'$ が平坦である。
% P10-E0667
\item $B \to B'$ が忠実平坦なら，
$\varphi'$ が平坦 $\Rightarrow \varphi$ が平坦である。
\item $i = 1, \ldots, n$ に対して $A \to B_i$ が平坦なら，
$A \to \prod_{i = 1, \ldots, n} B_i$ は平坦である。
\end{enumerate}
Noether 環の完備化は平坦であることを，以下では断りなく用いる
（Algebra, Lemma \ref{algebra-lemma-completion-flat}）。
もちろん $A \to A'$ と $B \to B'$ は平坦なので，特に
図式の横の矢印は平坦である。

\medskip\noindent
(1) の証明。$\varphi$ が平坦なら，合成
$A \to (A')^\wedge \to (B')^\wedge$ は平坦である。
従って More on Flatness, Lemma
\ref{flat-lemma-etale-flat-up-down} により，
$A' \to (B')^\wedge$ は平坦である。
そこで More on Algebra, Lemma
\ref{more-algebra-lemma-flat-after-completion}
を $R = A'$，イデアル $I(A')$，および
$M = (B')^\wedge = M^\wedge$ に適用すると，
$(A')^\wedge \to (B')^\wedge$ が平坦であることが分かる。

\medskip\noindent
(2) の証明。$\varphi'$ が平坦で，
$B \to B'$ が忠実平坦であると仮定する。
このとき合成 $A \to (A')^\wedge \to (B')^\wedge$ は平坦である。
また Formal Spaces, Lemma
\ref{formal-spaces-lemma-etale-surjective}
により，$B \to (B')^\wedge$ は忠実平坦である。
従って Algebra, Lemma
\ref{algebra-lemma-flatness-descends-more-general}
により，$\varphi : A \to B$ は平坦である。

\medskip\noindent
(3) の証明。省略する。
\end{proof}

\begin{lemma}
\label{restricted-lemma-base-change-flat-continuous}
Lemma \ref{restricted-lemma-flat-axioms} で定義された
$\textit{WAdm}^{Noeth}$ の射の性質を $P$ と書く。
Lemma \ref{restricted-lemma-finite-type-red} で定義された性質を
$\textit{WAdm}^{Noeth}$ の射の性質とみなし，$Q$ と書く。
Lemma \ref{restricted-lemma-finite-type} で定義された性質を
$\textit{WAdm}^{Noeth}$ の射の性質とみなし，$R$ と書く。このとき，
\begin{enumerate}
\item $P$ は $Q$ による基底変換で安定であり
（Formal Spaces, Remark
\ref{formal-spaces-remark-base-change-variant-variant-Noetherian}），
\item $P + R$ は基底変換で安定である
（Formal Spaces, Remark
\ref{formal-spaces-remark-base-change-variant-Noetherian}）。
\end{enumerate}
\end{lemma}

\begin{proof}
性質 $P$, $Q$, $R$ はいずれも $\textit{WAdm}^{Noeth}$
の射の局所的性質なので，主張は意味をもつ。
$\varphi : B \to A$ と $\psi : B \to C$ を
$\textit{WAdm}^{Noeth}$ の射とする。
% P10-E0668
$Q(\varphi)$ または $Q(\psi)$ のいずれかが成り立つなら，
Formal Spaces, Lemma
\ref{formal-spaces-lemma-completed-tensor-product}
により $A \widehat{\otimes}_B C$ は Noether である。
$R$ は $Q$ を含意するので
（Lemma \ref{restricted-lemma-finite-type-finite-type-red}），
これは (1), (2) の双方の場合に成り立つ。
これが最初に確認すべきことである。残るのは，
$C \to A \widehat{\otimes}_B C$ が平坦であることを示すことである。

\medskip\noindent
(1) の証明。定義イデアル $I \subset A$ と $J \subset B$ を固定する。
Lemma \ref{restricted-lemma-Noetherian-finite-type-red} により，
環写像 $B \to C$ は位相的有限型である。
従ってすべての $n \geq 1$ に対して
$B \to C/J^n$ は有限型である。
従って環 $A \otimes_B C/J^n$ は Noether である
（これは $A$ 上有限型であり，$A$ は Noether だからである）。
そこで $A \otimes_B C/J^n$ の $I$-進完備化
$A \widehat{\otimes}_B C/J^n$ は $C/J^n$ 上平坦である。
実際，$C/J^n \to A \otimes_B C/J^n$ は
$B \to A$ の基底変換として平坦であり，
$A \otimes_B C/J^n \to A \widehat{\otimes}_B C/J^n$ は
Algebra, Lemma \ref{algebra-lemma-completion-flat}
により平坦である。
$A \widehat{\otimes}_B C/J^n =
(A \widehat{\otimes}_B C)/J^n(A \widehat{\otimes}_B C)$
であることに注意せよ。細部は省略する。
従って $M = A \widehat{\otimes}_B C$ は，
$J$-進位相に関して完備な $C$-加群であり，すべての $n \geq 1$
に対して $M/J^nM$ は $C/J^n$ 上平坦である。
More on Algebra, Lemma
\ref{more-algebra-lemma-limit-flat}
により，これは $M$ が $C$ 上平坦であることを含意する。

\medskip\noindent
(2) の証明。この場合 $B \to A$ は adic なので，単に
$A \widehat{\otimes}_B C = \lim A \otimes_B C/J^n$ である。
Formal Spaces, Lemma
\ref{formal-spaces-lemma-completed-tensor-product}
を $C/J^n$ で $C$ を置き換えて適用すると，
環 $A \otimes_B C/J^n$ は Noether である。
前と同じ仕方で結論を得る。
\end{proof}

\begin{lemma}
\label{restricted-lemma-composition-flat-continuous}
Lemma \ref{restricted-lemma-flat-axioms} で定義された
$\textit{WAdm}^{Noeth}$ の射の性質を $P$ と書く。
このとき $P$ は合成で安定である
（Formal Spaces, Remark
\ref{formal-spaces-remark-composition-variant-Noetherian}）。
\end{lemma}

\begin{proof}
これは平坦環写像の合成が平坦であることから従う。
\end{proof}

\begin{definition}
\label{restricted-definition-flat}
$S$ をスキームとする。$f : X \to Y$ を，$S$ 上の局所 Noether
形式代数空間の射とする。任意の可換図式
$$
\xymatrix{
U \ar[d] \ar[r] & V \ar[d] \\
X \ar[r] & Y
}
$$
で，$U$ と $V$ がアフィン形式代数空間であり，$U \to X$ と
$V \to Y$ が代数空間により表現可能かつエタールであるものに対し，
射 $U \to V$ が adic Noether 位相環の平坦写像に対応するとき，
$f$ は{\it 平坦}であるという。
\end{definition}

\noindent
この条件を始域と終域でエタール局所的に確認できることを証明しよう。

\begin{lemma}
\label{restricted-lemma-flat-morphisms}
$S$ をスキームとする。$f : X \to Y$ を，$S$ 上の
局所 Noether 形式代数空間の射とする。次は同値である。
\begin{enumerate}
\item $f$ は平坦である。
\item 任意の可換図式
$$
\xymatrix{
U \ar[d] \ar[r] & V \ar[d] \\
X \ar[r] & Y
}
$$
で，$U$ と $V$ がアフィン形式代数空間であり，$U \to X$ と
$V \to Y$ が代数空間により表現可能かつエタールであるものに対し，
射 $U \to V$ は $\textit{WAdm}^{Noeth}$ の平坦写像に対応する。
\item Formal Spaces, Definition
\ref{formal-spaces-definition-formal-algebraic-space}
におけるような被覆 $\{Y_j \to Y\}$ が存在し，各 $j$ に対して
Formal Spaces, Definition
\ref{formal-spaces-definition-formal-algebraic-space}
におけるような被覆
$\{X_{ji} \to Y_j \times_Y X\}$ が存在して，各
$X_{ji} \to Y_j$ が $\textit{WAdm}^{Noeth}$ の平坦写像に対応する。
\item Formal Spaces, Definition
\ref{formal-spaces-definition-formal-algebraic-space}
におけるような被覆 $\{X_i \to X\}$ が存在し，各 $i$ に対して
分解 $X_i \to Y_i \to Y$ が存在する。ここで $Y_i$ は
アフィン形式代数空間，$Y_i \to Y$ は代数空間により表現可能かつ
エタールであり，$X_i \to Y_i$ は
$\textit{WAdm}^{Noeth}$ の平坦写像に対応する。
\end{enumerate}
\end{lemma}

\begin{proof}
(1) と (2) の同値性は Definition \ref{restricted-definition-flat} である。
(2), (3), (4) の同値性は，Lemma \ref{restricted-lemma-flat-axioms}
により平坦であることが $\text{WAdm}^{Noeth}$ の射の局所的性質である
ことと，Formal Spaces, Lemma
\ref{formal-spaces-lemma-property-defines-property-morphisms}
の，局所 Noether
% P10-E0669
形式代数空間の間の射に対する変種を適用することから従う。
この変種は Formal Spaces, Remark
\ref{formal-spaces-remark-variant-Noetherian} で述べられている。
\end{proof}

\begin{lemma}
\label{restricted-lemma-base-change-flat}
$S$ をスキームとする。$f : X \to Y$ と $g : Z \to Y$
を，$S$ 上の局所 Noether 形式代数空間の射とする。
\begin{enumerate}
\item $f$ が平坦で，$g_{red} : Z_{red} \to Y_{red}$ が
局所有限型なら，基底変換 $X \times_Y Z \to Z$ は平坦である。
\item $f$ が平坦かつ局所有限型なら，基底変換
$X \times_Y Z \to Z$ は平坦かつ局所有限型である。
\end{enumerate}
\end{lemma}

\begin{proof}
(1) は Formal Spaces, Remark
\ref{formal-spaces-remark-base-change-variant-variant-Noetherian},
Lemma \ref{restricted-lemma-base-change-flat-continuous} の (1),
Lemma \ref{restricted-lemma-flat-morphisms}, and
Lemma \ref{restricted-lemma-finite-type-red-morphisms}
を組み合わせれば従う。

\medskip\noindent
(2) は Formal Spaces, Remark
\ref{formal-spaces-remark-base-change-variant-Noetherian},
Lemma \ref{restricted-lemma-base-change-flat-continuous} の (2),
Lemma \ref{restricted-lemma-flat-morphisms}, and
Lemma \ref{restricted-lemma-finite-type-morphisms}
を組み合わせれば従う。
\end{proof}

\begin{lemma}
\label{restricted-lemma-composition-flat}
$S$ をスキームとする。$f : X \to Y$ と $g : Y \to Z$
を，$S$ 上の局所 Noether 形式代数空間の射とする。
$f$ と $g$ が平坦なら，$g \circ f$ も平坦である。
\end{lemma}

\begin{proof}
Formal Spaces, Remark
\ref{formal-spaces-remark-composition-variant-Noetherian}
と Lemma \ref{restricted-lemma-composition-flat-continuous}
を組み合わせればよい。
\end{proof}

\begin{lemma}
\label{restricted-lemma-representable-flat}
% P10-E0397
$S$ をスキームとする。$f : X \to Y$ を，$S$ 上の
局所 Noether 形式代数空間の射とする。
$f$ が代数空間により表現可能であり，Bootstrap, Definition
\ref{bootstrap-definition-property-transformation}
の意味で平坦なら，$f$ は Definition
\ref{restricted-definition-flat} の意味で平坦である。
\end{lemma}

\begin{proof}
これは健全性確認であり，証明は自明であるはずだが，完全にはそうでない。
読者には証明を読み飛ばすことを勧める。
$f$ が代数空間により表現可能であり，Bootstrap, Definition
\ref{bootstrap-definition-property-transformation}
の意味で平坦であると仮定する。可換図式
$$
\xymatrix{
U \ar[d] \ar[r] & V \ar[d] \\
X \ar[r] & Y
}
$$
で，$U$ と $V$ がアフィン形式代数空間であり，$U \to X$ と
$V \to Y$ が代数空間により表現可能かつエタールであるものを考える。
このとき射 $U \to V$ は，Formal Spaces, Lemma
\ref{formal-spaces-lemma-representable-local-property}
により，$\textit{WAdm}^{Noeth}$ の taut 写像
$B \to A$ に対応する。
これは $B \to A$ が adic であることを意味することに注意せよ
（Formal Spaces, Lemma
\ref{formal-spaces-lemma-adic-homomorphism}）。
特に，任意の定義イデアル $J \subset B$ に対して，
$A$ 上の位相は $J$-進位相であり，図式
$$
\xymatrix{
\Spec(A/J^nA) \ar[r] \ar[d] & \Spec(B/J^n) \ar[d] \\
U \ar[r] & V
}
$$
は Cartesian である。

\medskip\noindent
% P10-E0398
$T$ をスキームとする射 $T \to V$ を取る。このとき
\begin{align*}
X \times_Y T \to T\text{ は平坦}
& \Rightarrow
U \times_Y T \to T\text{ は平坦} \\
& \Rightarrow
U \times_V V \times_Y T \to T\text{ は平坦} \\
& \Rightarrow
U \times_V V \times_Y T \to V \times_Y T\text{ は平坦} \\
& \Rightarrow
U \times_V T \to T\text{ は平坦}
\end{align*}
最初の主張は $f$ に関する仮定である。
% P10-E0399
第1の含意は，$U \to X$ がエタール，従って平坦であり，
代数空間の平坦射の合成が平坦であることから従う。
第2の含意は
$U \times_Y T = U \times_V V \times_Y T$
であることから従う。第3の含意は More on Flatness, Lemma
\ref{flat-lemma-etale-flat-up-down}
から従う。第4の含意は，射
$T \to V \times_Y T$ により引き戻せることから従う。
従って $U \to V$ は Bootstrap, Definition
\ref{bootstrap-definition-property-transformation}
の意味で平坦である。連続環写像 $B \to A$ の言葉では，
これは環写像 $B/J^n \to A/J^nA$ が平坦であることを意味する
（上の図式を参照）。

\medskip\noindent
最後に，例えば More on Algebra, Lemma
\ref{more-algebra-lemma-limit-flat}
により，$B \to A$ が平坦であると結論できる。
\end{proof}

\section{Rig-閉点}
\label{restricted-section-rig-points}

\noindent
後の節で rig-平坦性を検査するためにこれを使えるだけの理論を展開する。
局所 Noether 形式スキームの場合などを論じている \cite{BL-I} には，
さらに多くの理論がある。

\begin{lemma}
\label{restricted-lemma-rig-point}
$A$ を Noether adic 位相環とする。
$\mathfrak q \subset A$ を素イデアルとする。次の条件は同値である。
\begin{enumerate}
\item ある定義イデアル $I \subset A$ に対して
$I \not \subset \mathfrak q$ であり，$\mathfrak q$ はこの性質に関して極大である，
\item ある定義イデアル $I \subset A$ に対して，素イデアル
$\mathfrak q$ は $\Spec(A) \setminus V(I)$ の閉点を定める，
\item 任意の定義イデアル $I \subset A$ に対して
$I \not \subset \mathfrak q$ であり，$\mathfrak q$ はこの性質に関して極大である，
\item 任意の定義イデアル $I \subset A$ に対して，素イデアル
$\mathfrak q$ は $\Spec(A) \setminus V(I)$ の閉点を定める，
\item $\dim(A/\mathfrak q) = 1$ であり，ある定義イデアル
$I \subset A$ に対して $I \not \subset \mathfrak q$ である，
\item $\dim(A/\mathfrak q) = 1$ であり，任意の定義イデアル
$I \subset A$ に対して $I \not \subset \mathfrak q$ である，
\item $\dim(A/\mathfrak q) = 1$ であり，$A/\mathfrak q$ 上に誘導される位相は
自明でない，
\item $A/\mathfrak q$ は $1$ 次元 Noether 完備局所整域であり，その極大イデアルは
$A$ の任意の定義イデアルの像の根基である，および
% P10-E0670
\item さらにここに加える。
\end{enumerate}
\end{lemma}

\begin{proof}
(1) と (2) が同値であることは明らかであり，同じ理由により
(3) と (4) も同値である。
$V(I)$ は $A$ の定義イデアル $I$ の選択によらないので，
(2) と (4) が同値であることも分かる。

\medskip\noindent
同値な条件 (1) -- (4) が成り立つと仮定する。
$\dim(A/\mathfrak q) > 1$ ならば，
$\dim((A/\mathfrak q)_\mathfrak m) > 1$ となる極大イデアル
$\mathfrak q \subset \mathfrak m \subset A$ を選べる。
このとき $\Spec((A/\mathfrak q)_\mathfrak m) - V(I(A/\mathfrak q)_\mathfrak m)$ は，
Algebra, Lemma
\ref{algebra-lemma-Noetherian-local-domain-dim-2-infinite-opens}
により無限集合となる。これは $\mathfrak q$ が
$\Spec(A) \setminus V(I)$ で閉じていることに反する。
従って (6) が成り立つ。(6) が (5) を含意することは自明である。

\medskip\noindent
逆に，(5) が成り立つと仮定する。$I \subset A$ を定義イデアルとする。
$A/\mathfrak q$ は $I(A/\mathfrak q)$ に関して完備なので
（例えば Algebra, Lemma \ref{algebra-lemma-completion-tensor} による），
Algebra, Lemma \ref{algebra-lemma-radical-completion} により，
$\Spec(A/\mathfrak q)$ のすべての閉点は $V(IA/\mathfrak q)$ に含まれる。
$\dim(A/\mathfrak q) = 1$ かつ $I \not \subset \mathfrak q$ であることから，
二つのことが従う：(a) $V(IA/\mathfrak q)$ は
$\Spec(A/\mathfrak q)$ の一般点と異なるすべての点を含み，
(b) $V(IA/\mathfrak q)$ は（有限）離散集合でなければならない。
(a) より $\mathfrak q$ は $\Spec(A) \setminus V(I)$ の閉点であり，
(2) が成り立つと結論する。

\medskip\noindent
引き続き (5) を仮定する。この有限離散空間
$V(IA/\mathfrak q)$ は，例えば More on Algebra, Lemma
\ref{more-algebra-lemma-irreducible-henselian-pair-connected}
（および完備対が Hensel 対であること；More on Algebra, Lemma
\ref{more-algebra-lemma-complete-henselian} を参照）により，一点集合でなければならない。
従って (8) が成り立つ。逆に，(8) が (5) を含意することは明らかである。

\medskip\noindent
以上により (1) -- (6) と (8) は同値である。
これらが (7) とも同値であることの確認は省略する。
\end{proof}

\noindent
このような素イデアルを扱いやすくするため，次の非標準的な用語を導入する。

\begin{definition}
\label{restricted-definition-rig-closed}
$A$ を Noether adic 位相環とする。
$\mathfrak q \subset A$ を素イデアルとする。Lemma
\ref{restricted-lemma-rig-point} の同値な条件が満たされるとき，
$\mathfrak q$ は {it rig-閉} であるという。
\end{definition}

\noindent
相対的な状況においてさえ rig-閉素イデアルが十分多く存在することを
本質的に述べる補題をいくつか必要とする。

\begin{lemma}
\label{restricted-lemma-rig-closed-point-relative-residue-field}
$\varphi : A \to B$ を $\textit{WAdm}^{Noeth}$ の射とする。
$\mathfrak a \subset A$ と $\mathfrak b \subset B$ を
位相的冪零元のイデアルとする。
$A/\mathfrak a \to B/\mathfrak b$ が有限型であると仮定する。
$\mathfrak q \subset B$ を rig-閉とする。
局所環 $B/\mathfrak q$ の剰余体 $\kappa$ は有限型
$A/\mathfrak a$-代数である。
\end{lemma}

\begin{proof}
$\mathfrak q$ を含む一意な極大イデアルを
$\mathfrak q \subset \mathfrak m \subset B$ とする。
このとき $\mathfrak b \subset \mathfrak m$ である。従って
$A/\mathfrak a \to B/\mathfrak b \to B/\mathfrak m = \kappa$ は
有限型である。
\end{proof}

\begin{lemma}
\label{restricted-lemma-rig-closed-point-relative}
$\varphi : A \to B$ を，adic かつ位相的有限型である
$\textit{WAdm}^{Noeth}$ の射とする。
$\mathfrak q \subset B$ を rig-閉とする。
$\mathfrak p = \varphi^{-1}(\mathfrak q) \subset A$ とおく。
$\mathfrak a \subset A$ を位相的冪零元のイデアルとする。
次の条件は同値である。
\begin{enumerate}
\item $B/\mathfrak q$ の剰余体 $\kappa$ は $A/\mathfrak a$ 上有限である，
\item $\mathfrak p \subset A$ は rig-閉である，
\item $A/\mathfrak p \subset B/\mathfrak q$ は環の有限拡大である。
\end{enumerate}
\end{lemma}

\begin{proof}
(1) を仮定する。$B/\mathfrak q$ は $1$ 次元 Noether 局所環であり，
その位相は極大イデアルから来る adic 位相であることを思い出そう。
$\varphi$ は adic なので，$A \to B/\mathfrak q$ は adic である。
従って $\varphi(\mathfrak a)$ は $B/\mathfrak q$ の非零イデアルである。
% P10-E0400
従って $B/\mathfrak q + \varphi(\mathfrak a)$ は有限長をもつ。
従って，仮定により $B/\mathfrak q + \varphi(\mathfrak a)$ は
$A/\mathfrak a$-加群として有限である。
Algebra, Lemma \ref{algebra-lemma-finite-over-complete-ring} により，
$B/\mathfrak q$ は $A$ 上有限である。従って (3) が成り立つ。

\medskip\noindent
(3) を仮定する。このとき Algebra, Lemma
\ref{algebra-lemma-integral-overring-surjective} により，
$\Spec(B/\mathfrak q) \to \Spec(A/\mathfrak p)$ は全射である。
従って (2) が成り立つ。

\medskip\noindent
(2) を仮定する。$A/\mathfrak p$ の剰余体を $\kappa'$ と書く。
Lemma \ref{restricted-lemma-rig-closed-point-relative-residue-field}
（および Lemma \ref{restricted-lemma-finite-type-finite-type-red}）により，
拡大 $\kappa/\kappa'$ は代数として有限生成である。
Hilbert の零点定理（Algebra, Lemma
\ref{algebra-lemma-field-finite-type-over-domain}）により，
$\kappa/\kappa'$ は有限拡大である。従って (1) が成り立つ。
\end{proof}

\begin{lemma}
\label{restricted-lemma-rig-closed-jacobson}
$\varphi : A \to B$ を，adic かつ位相的有限型である
$\textit{WAdm}^{Noeth}$ の射とする。
$\mathfrak q \subset B$ を rig-閉とする。
ある定義イデアル $I \subset A$ に対して $A/I$ が Jacobson ならば，
$\mathfrak p = \varphi^{-1}(\mathfrak q) \subset A$ は rig-閉である。
\end{lemma}

\begin{proof}
Lemma \ref{restricted-lemma-rig-closed-point-relative-residue-field}
（Lemma \ref{restricted-lemma-finite-type-finite-type-red} と組み合わせる）により，
$B/\mathfrak q$ の剰余体 $\kappa$ は $A/\mathfrak a$ 上有限型である。
$A/\mathfrak a$ は Jacobson なので，Algebra, Lemma
\ref{algebra-lemma-silly-jacobson} により，$\kappa$ は
$A/\mathfrak a$ 上有限である。Lemma
\ref{restricted-lemma-rig-closed-point-relative} により結論を得る。
\end{proof}

\begin{lemma}
\label{restricted-lemma-rig-closed-point-in-image}
$\varphi : A \to B$ を，adic かつ位相的有限型である
$\textit{WAdm}^{Noeth}$ の射とする。
$\mathfrak p \subset A$ を rig-閉とする。
$\mathfrak a \subset A$ と $\mathfrak b \subset B$ を
位相的冪零元のイデアルとする。$\varphi$ が平坦ならば，
次の条件は同値である。
\begin{enumerate}
\item $A/\mathfrak p$ の極大イデアルは
$\Spec(B/\mathfrak b) \to \Spec(A/\mathfrak a)$ の像に属する，
\item $\mathfrak p = \varphi^{-1}(\mathfrak q)$ となる
rig-閉素イデアル $\mathfrak q \subset B$ が存在する。
\end{enumerate}
また，これらが成り立つならば，$\varphi$，$\mathfrak p$，$\mathfrak q$ は
Lemma \ref{restricted-lemma-rig-closed-point-relative} の結論を満たす。
\end{lemma}

\begin{proof}
(2) $\Rightarrow$ (1) は直ちに従う。(1) を仮定する。
$\mathfrak q$ の存在を示すため，$A$ を $A/\mathfrak p$ で，
$B$ を $B/\mathfrak p B$ で置き換えてよい（詳細の一部は省略する）。
従って $(A, \mathfrak m, \kappa)$ は $1$ 次元 Noether 完備局所環，
$\mathfrak m = \mathfrak a$，$\mathfrak p = (0)$ であると仮定してよい。
条件 (1) は，単に
$B_0 = B \otimes_A \kappa = B/\mathfrak m B = B/\mathfrak a B$
が非零であることを述べる。$B_0$ は $\kappa$ 上有限型であることに注意せよ。
従って $\dim(B_0)$ に関する帰納法を使える。
$\dim(B_0) = 0$ ならば，任意の極小素イデアル
$\mathfrak q \subset B$ が求めるものである
（$A \to B$ の平坦性により，$\mathfrak q$ は
$\mathfrak p = (0)$ の上にある）。
$\dim(B_0) > 0$ ならば，$b \in B$ で，その像
$b_0 \in B_0$ が非零因子かつ非単元となるものを取れる；Algebra, Lemma
\ref{algebra-lemma-dim-not-zero-exists-nonzerodivisor-nonunit} を参照。
Algebra, Lemma \ref{algebra-lemma-grothendieck} により，
環 $B' = B/bB$ は $A$ 上平坦である。
$B'_0 = B' \otimes_A \kappa = B_0/(b_0)$ は非零なので，
$B'$ に帰納法の仮定を適用して結論を得る。
補題の最後の主張は Lemma
\ref{restricted-lemma-rig-closed-point-relative} から明らかである。
\end{proof}

\noindent
いくつか記法を導入する。

\begin{definition}
\label{restricted-definition-completed-principal-localization}
$A$ を，有限生成な定義イデアルをもつ adic 位相環とする。
$f \in A$ とする。$A$ の $f$ における主局所化
$A_f = A[1/f]$ を，$A$ の任意の定義イデアルに関して完備化したものを，
$A$ の {it 完備化主局所化} $A_{\{f\}}$ と呼ぶ。
\end{definition}

\noindent
念のため述べると，$f$ が位相的冪零ならば，$A_{\{f\}}$ は零環である。

\begin{lemma}
\label{restricted-lemma-rig-closed-point-in-localization}
$A$ を adic Noether 位相環とする。
$\mathfrak p \subset A$ を素イデアルとする。
$f \in A$ を，$A/\mathfrak p$ で単元に写る元とする。このとき
$$
\mathfrak p A_{\{f\}} =
\mathfrak p(A_f)^\wedge =
\mathfrak p \otimes_A (A_f)^\wedge =
(\mathfrak p_f)^\wedge
$$
は素イデアルであり，その商は
$$
A/\mathfrak p = (A/\mathfrak p) \otimes_A (A_f)^\wedge =
(A_f)^\wedge / \mathfrak p (A_f)^\wedge = A_{\{f\}}/\mathfrak p A_{\{f\}}
$$
である。
\end{lemma}

\begin{proof}
$A_f$ は Noether なので，環写像 $A \to A_f \to (A_f)^\wedge$ は平坦である。
任意の有限 $A$-加群 $M$ に対して，
$M \otimes_A (A_f)^\wedge$ は $M_f$ の完備化である。
$f$ が $M$ 上の単元ならば，$M_f = M$ は既に完備である。
Algebra, Section \ref{algebra-section-completion-noetherian} の議論を参照せよ。
これらの観察から結果は容易に従う。
\end{proof}

\begin{lemma}
\label{restricted-lemma-rig-closed-point-after-localization}
$\varphi : A \to B$ を，adic かつ位相的有限型である
$\textit{WAdm}^{Noeth}$ の射とする。
$\mathfrak q \subset B$ を rig-閉とする。
$B/\mathfrak q$ で単元に写る $f \in A$ が存在して，図式
$$
\vcenter{
\xymatrix{
B \ar[r] &
B_{\{f\}} \\
A \ar[r] \ar[u]_\varphi &
A_{\{f\}} \ar[u]_{\varphi_{\{f\}}}
}
}
\quad\text{素イデアルを伴い}\quad
\vcenter{
\xymatrix{
\mathfrak q \ar@{-}[r] \ar@{-}[d] &
\mathfrak q' \ar@{-}[d] \ar@{=}[r] &
\mathfrak q B_{\{f\}} \\
\mathfrak p \ar@{-}[r] &
\mathfrak p'
}
}
$$
を得る。ここで $\mathfrak p'$ は rig-閉である。すなわち，
写像 $A_{\{f\}} \to B_{\{f\}}$ と素イデアル
$\mathfrak q'$，$\mathfrak p'$ は，Lemma
\ref{restricted-lemma-rig-closed-point-relative} の同値な条件を満たす。
\end{lemma}

\begin{proof}
$\mathfrak q'$ の記述については Lemma
\ref{restricted-lemma-rig-closed-point-in-localization} を参照されたい。
この補題が主張しているのは，$f$ を適切に選べば素イデアル
$\mathfrak p'$ が $\dim((A_f)^\wedge/\mathfrak p') = 1$ を満たすということだけである。
Lemma \ref{restricted-lemma-rig-closed-point-relative} により，これはさらに，
$B/\mathfrak q = (B_f)^\wedge/\mathfrak q'$ の剰余体 $\kappa$ が
$(A_f)^\wedge/\mathfrak a' = (A/\mathfrak a)_f$ 上有限であることを意味する。
Lemma \ref{restricted-lemma-rig-closed-point-relative-residue-field} により，
$A/\mathfrak a \to \kappa$ は有限型代数準同型である。
Algebra, Lemma \ref{algebra-lemma-field-finite-type-over-domain}
の形の Hilbert の零点定理により，$\kappa$ で単元に写り，
$\kappa$ が $A_f$ 上有限となる $f \in A$ を取れる。これで証明が完了する。
\end{proof}

\begin{lemma}
\label{restricted-lemma-rig-closed-point-variables}
$A$ を Noether adic 位相環とする。
$A\{x_1, \ldots, x_n\}$ を $A$ 上の制限冪級数環とする。
$\mathfrak q \subset A\{x_1, \ldots, x_n\}$ を素イデアルとする。
$\mathfrak q' = A[x_1, \ldots, x_n] \cap \mathfrak q$ および
$\mathfrak p = A \cap \mathfrak q$ とおく。
$\mathfrak q$ と $\mathfrak p$ が rig-閉ならば，写像
$$
A[x_1, \ldots, x_n]_{\mathfrak q'}
\to
A\{x_1, \ldots, x_n\}_\mathfrak q
$$
は，それぞれの極大イデアルに関する完備化の間の同型を定める。
\end{lemma}

\begin{proof}
Lemma \ref{restricted-lemma-rig-closed-point-relative} により，環写像
$A/\mathfrak p \to A\{x_1, \ldots, x_n\}/\mathfrak q$ は有限である。
各 $m \geq 1$ に対して，加群 $\mathfrak q^m/\mathfrak q^{m + 1}$ は，
有限 $A\{x_1, \ldots, x_n\}/\mathfrak q$-加群なので，$A$ 上有限である。
% P10-E0401
従って $A\{x_1, \ldots x_n\}/\mathfrak q^m$ は有限 $A$-加群である。
従って $A[x_1, \ldots, x_n] \to A\{x_1, \ldots, x_n\}/\mathfrak q^m$ は
全射である（像は稠密な $A$-部分加群だからである）。
これより，すべての $m$ に対して
$A[x_1, \ldots, x_n]/(\mathfrak q')^m \to A\{x_1, \ldots, x_n\}/\mathfrak q^m$
が同型であることが直ちに従う。補題はこれから容易に従う。
ヒント：$\mathfrak p$ に含まれない位相的冪零元 $g \in A$ を取る。
このとき完備化の写像は
% P10-E0671
$$
\lim_m \left(A[x_1, \ldots, x_n]/(\mathfrak q')^m\right)_g
\longrightarrow
\left(A\{x_1, \ldots, x_n\}/\mathfrak q^m\right)_g
$$
である。詳細の一部は省略する。
\end{proof}

\begin{lemma}
\label{restricted-lemma-rig-closed-point-etale}
$\varphi : A \to B$ を $\textit{WAdm}^{Noeth}$ の射とする。
$\varphi$ は adic，位相的有限型かつ平坦であり，ある（従って任意の）
定義イデアル $I \subset A$ に対して $A/I \to B/IB$ は
\'etale であると仮定する。
$\mathfrak q \subset B$ を rig-閉とし，
$\mathfrak p = A \cap \mathfrak q$ も rig-閉であると仮定する。このとき
$\mathfrak p B_\mathfrak q = \mathfrak q B_\mathfrak q$ である。
\end{lemma}

\begin{proof}
$A/\mathfrak p$ という $1$ 次元完備局所環の剰余体を $\kappa$ とする。
$A/I \to B/IB$ は \'etale なので，$B \otimes_A \kappa$ は
$\kappa$ の有限分離拡大の有限直積である；Algebra, Lemma
\ref{algebra-lemma-etale-over-field} を参照。
その因子の一つが $B/\mathfrak q$ の剰余体である。
Algebra, Lemma \ref{algebra-lemma-finite-over-complete-ring} により，
$B/\mathfrak p B$ は有限 $A/\mathfrak p$-代数である。また平坦でもある。
以上を合わせると，Algebra, Lemma
\ref{algebra-lemma-characterize-etale} により，
$A/\mathfrak p  \to B /\mathfrak p B$ は有限 \'etale である。
従って $B/\mathfrak p B$ は被約であり，これは補題の主張を含意する
（詳細は省略する）。
\end{proof}

\begin{lemma}
\label{restricted-lemma-fibre-regular}
$A$ を adic Noether 位相環とする。
$\mathfrak p \subset A$ を rig-閉素イデアルとする。
任意の $n \geq 1$ に対して，環写像
$$
A/\mathfrak p
\longrightarrow
A\{x_1, \ldots, x_n\} \otimes_A A/\mathfrak p =
A/\mathfrak p\{x_1, \ldots, x_n\}
$$
は正則である。特に，代数
$A\{x_1, \ldots, x_n\} \otimes_A \kappa(\mathfrak p)$ は
$\kappa(\mathfrak p)$ 上幾何学的に正則である。
\end{lemma}

\begin{proof}
正則環写像に関するいくつかの事実を使う。読者は More on Algebra, Section
\ref{more-algebra-section-regular} にそれらを見いだせる。
$A/\mathfrak p$ は Noether 完備局所環なので excellent である
（More on Algebra, Proposition
\ref{more-algebra-proposition-ubiquity-excellent}）。
従って $A/\mathfrak p[x_1, \ldots, x_n]$ も excellent である
（同じ参照による）。従って More on Algebra, Lemma
\ref{more-algebra-lemma-map-G-ring-to-completion-regular} により，
$A/\mathfrak p[x_1, \ldots, x_n] \to A/\mathfrak p\{x_1, \ldots, x_n\}$
は正則環準同型である。もちろん
$A/\mathfrak p \to A/\mathfrak p[x_1, \ldots, x_n]$ は滑らかであり，
従って正則である。正則環写像の合成は正則なので，証明は完了する。
\end{proof}

\section{Rig-平坦準同型}
\label{restricted-section-rig-flat-homomorphisms}

\noindent
この節では，adic Noether 位相環の rig-平坦準同型を定義する。

\begin{lemma}
\label{restricted-lemma-naively-rig-flat-continuous}
$\varphi : A \to B$ を $\textit{WAdm}^{adic*}$ の射とする
（Formal Spaces, Section \ref{formal-spaces-section-morphisms-rings}）。
$\varphi$ は adic であると仮定する。次の条件は同値である：
\begin{enumerate}
\item 任意の位相的冪零元 $f \in A$ に対して，$B_f$ は $A$ 上平坦である，
\item 任意の位相的冪零元 $g \in B$ に対して，$B_g$ は $A$ 上平坦である，
\item 定義イデアルを含まない任意の素イデアル
$\mathfrak q \subset B$ に対して，$B_\mathfrak q$ は $A$ 上平坦である，
\item 任意の rig-閉素イデアル $\mathfrak q \subset B$ に対して，
$B_\mathfrak q$ は $A$ 上平坦である，および
% P10-E0672
\item さらにここに加える。
\end{enumerate}
\end{lemma}

\begin{proof}
定義と Algebra, Lemma \ref{algebra-lemma-flat-localization} から従う。
\end{proof}

\begin{definition}
\label{restricted-definition-naively-rig-flat}
$\varphi : A \to B$ を adic Noether 位相環の間の連続環準同型とする。
すなわち，$\varphi$ は $\textit{WAdm}^{Noeth}$ の射である。
$\varphi$ が adic，位相的有限型であり，かつ Lemma
\ref{restricted-lemma-naively-rig-flat-continuous} の同値な条件を満たすとき，
$\varphi$ は {it 素朴に rig-平坦} であるという。
\end{definition}

\noindent
次の例は，この概念が「局所化」されないことを示す。

\begin{example}
\label{restricted-example-recompletion-not-rig-flat}
Examples, Lemma \ref{examples-lemma-nonreduced-recompletion} により，
主イデアル $I = (a)$ に関して完備な $2$ 次元 Noether 局所整域
$(A, \mathfrak m)$ と元 $f \in \mathfrak m$，$f \not \in I$ で，
次の性質をもつものが存在する：環 $A_{\{f\}}[1/a]$ は非被約である。
ここで $A_{\{f\}}$ は主局所化 $A_f$ の $I$-進完備化
$(A_f)^\wedge$ である。念のため述べると，環
$A_{\{f\}}[1/a]$ は非零である。
$B = A_{\{f\}}/ \text{nil}(A_{\{f\}})$ をその冪零根基による商とする。
$A \to B$ は adic かつ位相的有限型であることに注意せよ。
実際，$x$ を $B$ における $1/f$ の像へ送る写像により，
$B$ は $A\{x\} = A[x]^\wedge$ の商である。
$a$ を含まない $B$ の任意の素イデアル $\mathfrak q$ は，
$(0) \subset A$ の上になければならない\footnote{実際，$B$ は
$a$-進完備なので，$a \in \mathfrak q'$ となる
$\mathfrak q \subset \mathfrak q' \subset B$ を取れる。
このとき $\mathfrak p' = A \cap \mathfrak q'$ は $a$ を含むが
$f$ を含まないので，高さ $1$ の素イデアルである。
さらに，$a \not \in \mathfrak p$ であるから，
$\mathfrak p = A \cap \mathfrak q$ は $\mathfrak p'$ に真に含まれる。
$\dim(A) = 2$ なので，$\mathfrak p = (0)$ である。}。
従って $B_\mathfrak q$ は $A$ の分数体上の加群なので，$A$ 上平坦である。
従って $A \to B$ は素朴に rig-平坦である。
他方，写像
$$
A_{\{f\}}
\longrightarrow
B_{\{f\}} = (B_f)^\wedge = B = A_{\{f\}} /\text{nil}(A_{\{f\}})
$$
は，$a$ を可逆にした後も平坦でない。というのも，非自明な全射
$A_{\{f\}}[1/a] \to A_{\{f\}}[1/a]/\text{nil}(A_{\{f\}}[1/a])$
を得るからである。従って
$A_{\{f\}} \to B_{\{f\}}^\wedge$ は素朴に rig-平坦でない。
\end{example}

\noindent
次の定義を使えば，この問題を容易に回避できることが分かる。

\begin{definition}
\label{restricted-definition-rig-flat-continuous-homomorphism}
$\varphi : A \to B$ を adic Noether 位相環の間の連続環準同型とする。
すなわち，$\varphi$ は $\textit{WAdm}^{Noeth}$ の射である。
$\varphi$ が adic，位相的有限型であり，かつすべての $f \in A$ に対して
誘導される写像
$$
A_{\{f\}} \longrightarrow B_{\{f\}}
$$
が素朴に rig-平坦であるとき（Definition
\ref{restricted-definition-naively-rig-flat}），$\varphi$ は {it rig-平坦} であるという。
\end{definition}

\noindent
上の定義で $f = 1$ とおけば，rig-平坦性が素朴な rig-平坦性を
含意することが分かる。この例は逆が偽であることを示す。
しかし次の補題が示すように，多くの状況では rig-平坦性と
その素朴な版との差を気にする必要はない。

\begin{lemma}
\label{restricted-lemma-rig-flat-naive}
$\varphi : A \to B$ を $\textit{WAdm}^{Noeth}$ の射とする。
ある（同値なことに任意の）定義イデアル $I \subset A$ に対して
$A/I$ が Jacobson であり，$\varphi$ が素朴に rig-平坦ならば，
$\varphi$ は rig-平坦である。
\end{lemma}

\begin{proof}
$\varphi$ は素朴に rig-平坦であると仮定する。
まず，仮定から明らかに従うことをいくつか述べる。
すなわち，$f \in A$ とする。このとき
$A, B, A_{\{f\}}, B_{\{f\}}$ は Noether adic 位相環である。
写像 $A \to A_{\{f\}} \to B_{\{f\}}$ と
$A \to B \to B_{\{f\}}$ は adic かつ位相的有限型である。
環写像 $A \to A_{\{f\}}$ と $B \to B_{\{f\}}$ は，
$A \to A_f$ と $B \to B_f$，および Algebra, Lemma
\ref{algebra-lemma-completion-flat} により平坦な完備化写像の合成なので，
平坦である。各環 $A, B, A_{\{f\}}, B_{\{f\}}$ の $I$ による商は
$A/I$ 上有限型であり，従ってやはり Jacobson である
（Algebra, Proposition \ref{algebra-proposition-Jacobson-permanence}）。

\medskip\noindent
$\mathfrak q' \subset B_{\{f\}}$ を rig-閉とする。
Lemma \ref{restricted-lemma-naively-rig-flat-continuous} により，
$(B_{\{f\}})_{\mathfrak q'}$ が $A_{\{f\}}$ 上平坦であることを
示せば十分である。Lemma \ref{restricted-lemma-rig-closed-jacobson} により，
$\mathfrak q'$ の下にある素イデアル
$\mathfrak q \subset B$，$\mathfrak p' \subset A_{\{f\}}$，
$\mathfrak p \subset A$ は rig-閉である。図式
$$
\xymatrix{
B_\mathfrak q \ar[r] &
(B_{\{f\}})_{\mathfrak q'} \\
A_\mathfrak p \ar[u] \ar[r] &
(A_{\{f\}})_{\mathfrak p'} \ar[u]
}
$$
に $M = B_\mathfrak q$ として Algebra, Lemma
\ref{algebra-lemma-yet-another-variant-local-criterion-flatness}
を適用する。まだ確認していない唯一の仮定は，$\mathfrak p$ が
$(A_{\{f\}})_{\mathfrak p'}$ の極大イデアルを生成することである。
これは Lemma \ref{restricted-lemma-rig-closed-point-in-localization} から従う。
ここで，$f$ が $A/\mathfrak p$ の単元に写ることを示すために，
$\mathfrak p$ と $\mathfrak p'$ が rig-閉であることを用いる
（Jacobson という仮定なしには，証明のこの段階だけが成り立たない）。
実際，これは $A/\mathfrak p \to A_{\{f\}}/\mathfrak p'$ が
局所環の有限包含であり（Lemma
\ref{restricted-lemma-rig-closed-point-relative}），しかも $f$ は後者で単元に
写ることを述べている。
\end{proof}

\begin{lemma}
\label{restricted-lemma-rig-flat-base-change}
$\varphi : A \to B$ と $A \to C$ を $\textit{WAdm}^{Noeth}$ の射とする。
$\varphi$ は rig-平坦であり，$A \to C$ は adic かつ位相的有限型であると
仮定する。このとき $C \to B \widehat{\otimes}_A C$ は rig-平坦である。
\end{lemma}

\begin{proof}
$\varphi$ は rig-平坦であると仮定する。$f \in C$ を元とする。
$C_{\{f\}} \to B \widehat{\otimes}_A C_{\{f\}}$ が
素朴に rig-平坦であることを示さなければならない。
$C$ を $C_{\{f\}}$ で置き換えられるので，
% P10-E0673
$C \to B \widehat{\otimes}_A C$ が素朴に rig-平坦であることを
示せば十分である。

\medskip\noindent
$A \to C$ が全射である場合，またはより一般に $C$ が有限
$A$-加群である場合，完備 Noether 環上の有限加群は完備なので，
$B \otimes_A C = B \widehat{\otimes}_A C$ である；Algebra, Lemma
\ref{algebra-lemma-completion-tensor} を参照。
通常の平坦性の基底変換（Algebra, Lemma
\ref{algebra-lemma-flat-base-change}）により，この場合には
$\varphi$ の素朴な rig-平坦性から
% P10-E0674
$C \to B \times_A C$ の素朴な rig-平坦性が従う。

\medskip\noindent
一般の場合には，$A \to C$ を
$A \to A\{x_1, \ldots, x_n\} \to C$ と分解できる。
ここで $A\{x_1, \ldots, x_n\}$ は制限冪級数環であり，
$A\{x_1, \ldots, x_n\} \to C$ は全射である。
従って $C = A\{x_1, \ldots, x_n\}$ の場合に
% P10-E0675
$C \to B \widehat{\otimes}_A B$ が素朴に rig-平坦であることを
示せば十分である。
$A\{x_1, \ldots, x_n\} = A\{x_1, \ldots, x_{n - 1}\}\{x_n\}$
なので，$n$ に関する帰納法により，次の段落で論じる場合に帰着する。

\medskip\noindent
ここでは $C = A\{x\}$ である。
$B \widehat{\otimes}_A C = B\{x\}$ であることに注意せよ。
$A\{x\} \to B\{x\}$ が素朴に rig-平坦であることを示さなければならない。
$\mathfrak q \subset B\{x\}$ を rig-閉素イデアルとする。
$B\{x\}_{\mathfrak q}$ が $A\{x\}$ 上平坦であることを示さなければならない。
$\mathfrak p = A \cap \mathfrak q$ とおく。
Lemma \ref{restricted-lemma-rig-closed-point-after-localization} により，
$f \in A$ で，$f$ が $B\{x\}/\mathfrak q$ で単元に写り，誘導される
$A_{\{f\}}$ の素イデアル $\mathfrak p'$ が rig-閉となるものを取れる。
以下では $A_{\{f\}}\{x\} = A\{x\}_{\{f\}}$ を用いる。
$B$ についても同様である。詳細は省略する。図式
$$
\xymatrix{
(B\{x\})_{\mathfrak q} \ar[r] &
(B_{\{f\}}\{x\})_{\mathfrak q'} \\
A\{x\} \ar[r] \ar[u] &
A_{\{f\}}\{x\} \ar[u]
}
$$
を考える。左の縦射が平坦であることを示したい。
上の横射は，局所環の局所準同型であり，かつ
$B_{\{f\}}\{x\}$ が Noether 環 $B\{x\}$ の局所化の完備化なので
平坦であることから，忠実平坦である。同様に，下の横射は平坦である。
従って右の縦射が平坦であることを示せば十分である。
これは次の段落で論じる場合に帰着する。

\medskip\noindent
ここでは $C = A\{x\}$ であり，rig-閉素イデアル
$\mathfrak q \subset B\{x\}$ が与えられ，
$\mathfrak p = A \cap \mathfrak q$ も rig-閉であるとする。
Lemma \ref{restricted-lemma-rig-closed-point-relative} により，中間の素イデアル
$B \cap \mathfrak q$ と $A\{x\} \cap \mathfrak q$ も rig-閉である。
Noether 局所環の局所準同型からなる図式
$$
\xymatrix{
(B[x])_{B[x] \cap \mathfrak q} \ar[r] &
(B\{x\})_{\mathfrak q} \\
(A[x])_{A[x] \cap \mathfrak q} \ar[r] \ar[u] &
(A\{x\})_{A\{x\} \cap \mathfrak q} \ar[u]
}
$$
を考える。Lemma \ref{restricted-lemma-rig-closed-point-variables} により，
横射は完備化の間の同型を定める。左の縦射が平坦であることは既に分かっている
（$A \to B$ は素朴に rig-平坦であり，従って
$A[x] \to B[x]$ は定義イデアルが定める閉部分の外で平坦だからである）。
従って最後に More on Algebra, Lemma
\ref{more-algebra-lemma-flat-completion} により結論を得る。
\end{proof}

\begin{lemma}
\label{restricted-lemma-rig-flat-local-etale}
可換図式
$$
\xymatrix{
B \ar[r] & B' \\
A \ar[r] \ar[u]^\varphi & A' \ar[u]_{\varphi'}
}
$$
を考える。これは $\textit{WAdm}^{Noeth}$ の図式であり，すべての射は
adic かつ位相的有限型であるとする。
$A \to A'$ と $B \to B'$ は平坦であると仮定する。
$I \subset A$ を定義イデアルとする。
$\varphi$ が rig-平坦であり，$A/I \to A'/IA'$ が \'etale ならば，
$\varphi'$ は rig-平坦である。
\end{lemma}

\begin{proof}
% P10-E0676
$f \in A'$ が与えられたとき，補題の仮定は図式
$$
\xymatrix{
B \ar[r] & (B')_{\{f\}} \\
A \ar[r] \ar[u]^\varphi & (A')_{\{f\}} \ar[u]
}
$$
についても成り立つ。従って $\varphi'$ が素朴に rig-平坦であることを
示せば十分である。

\medskip\noindent
rig-閉素イデアル $\mathfrak q' \subset B'$ を取る。
$(B')_{\mathfrak q'}$ が $A'$ 上平坦であることを示さなければならない。
$B'/\mathfrak q'$ で単元に写り，$A_{\{f\}}$ の誘導される素イデアル
$\mathfrak p''$ が rig-閉となる $f \in A$ を選べる；Lemma
\ref{restricted-lemma-rig-closed-point-after-localization} を参照。
正確には，ここで $\mathfrak q'' = \mathfrak q' B'_{\{f\}}$ および
$\mathfrak p'' = A_{\{f\}} \cap \mathfrak q''$ である。図式
% P10-E0678
$$
\xymatrix{
B'_{\mathfrak q'} \ar[r] &
(B'_{\{f\}})_{\mathfrak q''} \\
A \ar[r] \ar[u] &
A_{\{f\}} \ar[u]
}
$$
を考える。左の縦射が平坦であることを示したい。
上の横射は，局所環の局所準同型であり，かつ
$B'_{\{f\}}$ が Noether 環 $B'_f$ の局所化の完備化なので
平坦であることから，忠実平坦である。同様に，下の横射は平坦である。
従って右の縦射が平坦であることを示せば十分である。
最後に，補題のすべての仮定は図式
$$
\xymatrix{
B_{\{f\}} \ar[r] &
B'_{\{f\}} \\
A_{\{f\}} \ar[r] \ar[u] &
A'_{\{f\}} \ar[u]
}
$$
についても成り立つ。これは次の段落で論じる場合への帰着を与える。

\medskip\noindent
rig-閉素イデアル $\mathfrak q' \subset B'$ を取り，
$\mathfrak p = A \cap \mathfrak q'$ も rig-閉であると仮定する。
このとき素イデアル $\mathfrak q = B \cap \mathfrak q'$ と
$\mathfrak p' = A' \cap \mathfrak q'$ も rig-閉である；Lemma
\ref{restricted-lemma-rig-closed-point-relative} を参照。
図式
$$
\xymatrix{
B_\mathfrak q \ar[r] &
B'_{\mathfrak q'} \\
A_\mathfrak p \ar[u] \ar[r] &
A'_{\mathfrak p'} \ar[u]
}
$$
に $M = B_\mathfrak q$ として Algebra, Lemma
\ref{algebra-lemma-yet-another-variant-local-criterion-flatness}
を適用する。まだ確認していない唯一の仮定は，$\mathfrak p$ が
$A'_{\mathfrak p'}$ の極大イデアルを生成することである。
これは Lemma \ref{restricted-lemma-rig-closed-point-etale} から従う。
\end{proof}

\begin{lemma}
\label{restricted-lemma-rig-flat-local-down}
可換図式
$$
\xymatrix{
B \ar[r] & B' \\
A \ar[r] \ar[u]^\varphi & A' \ar[u]_{\varphi'}
}
$$
を考える。これは $\textit{WAdm}^{Noeth}$ の図式であり，すべての射は
adic かつ位相的有限型であるとする。
% P10-E0677
$A \to A'$ は平坦，$B \to B'$ は忠実平坦であると仮定する。
$\varphi'$ が rig-平坦ならば，$\varphi$ は rig-平坦である。
\end{lemma}

\begin{proof}
% P10-E0676
$f \in A$ が与えられたとき，補題の仮定は図式
$$
\xymatrix{
B_{\{f\}} \ar[r] & (B')_{\{f\}} \\
A_{\{f\}} \ar[r] \ar[u]^\varphi & (A')_{\{f\}} \ar[u]
}
$$
についても成り立つ
（忠実平坦性に関する条件を確認するには，次を用いる：
$B \to B'$ が忠実平坦であることは，$B \to B'$ が平坦であり，かつ
ある定義イデアル $I \subset A$ に対して
$\Spec(B'/IB') \to \Spec(B/IB)$ が全射であることと同値である）。
従って $\varphi$ が素朴に rig-平坦であることを示せば十分である。
しかし，$\varphi'$ が素朴に rig-平坦であり，
$\Spec(B') \to \Spec(B)$ が全射であることは分かっている。
結果はこれから直ちに従う。
\end{proof}

\noindent
最後に，rig-平坦性が局所的性質であることを示せる。

\begin{lemma}
\label{restricted-lemma-rig-flat-axioms}
$\textit{WAdm}^{Noeth}$ の射上の性質
$P(\varphi)=$「$\varphi$ は rig-平坦である」は，Formal Spaces, Remark
\ref{formal-spaces-remark-variant-adic-star} で定義される局所的性質である。
\end{lemma}

\begin{proof}
この主張が何を意味するかを思い出そう。まず，
$\textit{WAdm}^{Noeth}$ は，対象が adic Noether 位相環，射が連続環準同型である
圏である。次の条件を満たす可換図式
$$
\xymatrix{
B \ar[r] & (B')^\wedge \\
A \ar[r] \ar[u]^\varphi & (A')^\wedge \ar[u]_{\varphi'}
}
$$
を考える：
$A$ と $B$ は adic Noether 位相環であり，
$A \to A'$ と $B \to B'$ は \'etale 環写像であり，
ある定義イデアル $I \subset A$ に対して
$(A')^\wedge = \lim A'/I^nA'$ であり，
ある定義イデアル $J \subset B$ に対して
$(B')^\wedge = \lim B'/J^nB'$ であり，
$\varphi : A \to B$ と
$\varphi' : (A')^\wedge \to (B')^\wedge$ は連続である。
Formal Spaces, Lemma
\ref{formal-spaces-lemma-completion-in-sub} により，
$(A')^\wedge$ と $(B')^\wedge$ は adic Noether 位相環であることに注意せよ。
次を示さなければならない。
\begin{enumerate}
\item $\varphi$ が rig-平坦 $\Rightarrow \varphi'$ は rig-平坦である，
\item $B \to B'$ が忠実平坦ならば，
$\varphi'$ が rig-平坦 $\Rightarrow \varphi$ は rig-平坦である，および
\item $i = 1, \ldots, n$ に対して $A \to B_i$ が rig-平坦ならば，
$A \to \prod_{i = 1, \ldots, n} B_i$ は rig-平坦である。
\end{enumerate}
adic かつ位相的有限型であるという性質は条件 (1)，(2)，(3) を満たす；
Lemma \ref{restricted-lemma-finite-type} を参照。
従って性質「rig-平坦」について (1)，(2)，(3) を確認するときには，
環写像がすべて adic かつ位相的有限型であると既に仮定してよい。
このとき (1) と (2) は Lemma
\ref{restricted-lemma-rig-flat-local-etale} と
\ref{restricted-lemma-rig-flat-local-down} から従う。
(3) の自明な証明は省略する。
\end{proof}

\begin{lemma}
\label{restricted-lemma-composition-rig-flat-continuous}
$\textit{WAdm}^{Noeth}$ の射上の性質
$P(\varphi)=$「$\varphi$ は rig-平坦である」は，Formal Spaces, Remark
\ref{formal-spaces-remark-composition-variant-Noetherian}
で定義される意味で合成について安定である。
\end{lemma}

\begin{proof}
Lemma \ref{restricted-lemma-rig-flat-axioms} により，この主張には意味がある。
真であることを見るため，$\textit{WAdm}^{Noeth}$ における
rig-平坦射 $A \to B$ と $B \to C$ が与えられたと仮定する。
Lemma \ref{restricted-lemma-composition-finite-type} により，
$A \to C$ は adic かつ位相的有限型である。
証明を終えるには，すべての $f \in A$ に対して写像
$A_{\{f\}} \to C_{\{f\}}$ が素朴に rig-平坦であることを示せばよい。
$A_{\{f\}} \to B_{\{f\}}$ と $B_{\{f\}} \to C_{\{f\}}$ は
素朴に rig-平坦なので，素朴に rig-平坦な写像の合成が素朴に
rig-平坦であることを示せば十分である。
これは Algebra, Lemma \ref{algebra-lemma-composition-flat} の帰結である。
\end{proof}

\section{Rig-平坦射}
\label{restricted-section-rig-flat-morphisms}

\noindent
% P10-E0679
この節では Section \ref{restricted-section-rig-flat-homomorphisms} の結果を用いて，
局所 Noether 代数空間の rig-平坦射を定義する。

\begin{definition}
\label{restricted-definition-rig-flat}
$S$ をスキームとする。$f : X \to Y$ を，$S$ 上の
局所 Noether 形式代数空間の射とする。
アフィン形式代数空間 $U$，$V$ をもつ任意の可換図式
$$
\xymatrix{
U \ar[d] \ar[r] & V \ar[d] \\
X \ar[r] & Y
}
$$
で，$U \to X$ と $V \to Y$ が代数空間により表現可能かつ
\'etale であるものに対し，射 $U \to V$ が adic Noether 位相環の
rig-平坦写像に対応するとき，$f$ は {it rig-平坦} であるという。
\end{definition}

\noindent
この条件を始域と終域上で \'etale 局所的に検査できることを証明しよう。

\begin{lemma}
\label{restricted-lemma-rig-flat-morphisms}
$S$ をスキームとする。$f : X \to Y$ を，$S$ 上の
局所 Noether 形式代数空間の射とする。次の条件は同値である。
\begin{enumerate}
\item $f$ は rig-平坦である，
\item アフィン形式代数空間 $U$，$V$ をもつ任意の可換図式
$$
\xymatrix{
U \ar[d] \ar[r] & V \ar[d] \\
X \ar[r] & Y
}
$$
で，$U \to X$ と $V \to Y$ が代数空間により表現可能かつ
\'etale であるものに対し，射 $U \to V$ は
$\textit{WAdm}^{Noeth}$ の rig-平坦写像に対応する，
\item Formal Spaces, Definition
\ref{formal-spaces-definition-formal-algebraic-space}
の意味での被覆 $\{Y_j \to Y\}$ が存在し，各 $j$ に対して
Formal Spaces, Definition
\ref{formal-spaces-definition-formal-algebraic-space}
の意味での被覆 $\{X_{ji} \to Y_j \times_Y X\}$ が存在して，
各 $X_{ji} \to Y_j$ は $\textit{WAdm}^{Noeth}$ の
rig-平坦写像に対応する，および
\item Formal Spaces, Definition
\ref{formal-spaces-definition-formal-algebraic-space}
の意味での被覆 $\{X_i \to X\}$ が存在し，各 $i$ に対して
分解 $X_i \to Y_i \to Y$ が存在する。ここで $Y_i$ は
アフィン形式代数空間，$Y_i \to Y$ は代数空間により表現可能かつ
\'etale であり，$X_i \to Y_i$ は $\textit{WAdm}^{Noeth}$ の
rig-平坦写像に対応する。
\end{enumerate}
\end{lemma}

\begin{proof}
(1) と (2) の同値性は Definition \ref{restricted-definition-rig-flat} である。
(2)，(3)，(4) の同値性は，Lemma
\ref{restricted-lemma-rig-flat-axioms} により rig-平坦性が
$\text{WAdm}^{Noeth}$ の射の局所的性質であること，および
% P10-E0679
局所 Noether 代数空間の間の射に対する Formal Spaces, Lemma
\ref{formal-spaces-lemma-property-defines-property-morphisms}
の変種を適用することから従う。この変種は Formal Spaces, Remark
\ref{formal-spaces-remark-variant-Noetherian}
で述べられている。
\end{proof}

\begin{lemma}
\label{restricted-lemma-base-change-rig-flat}
$S$ をスキームとする。$f : X \to Y$ と $g : Z \to Y$ を，
$S$ 上の局所 Noether 形式代数空間の射とする。
$f$ が rig-平坦で $g$ が局所有限型ならば，基底変換
$X \times_Y Z \to Z$ は rig-平坦である。
\end{lemma}

\begin{proof}
Formal Spaces, Remark
\ref{formal-spaces-remark-base-change-variant-variant-Noetherian}
および Formal Spaces, Section
\ref{formal-spaces-section-adic} の議論により，これは Lemma
\ref{restricted-lemma-rig-flat-base-change} から従う。
\end{proof}

\begin{lemma}
\label{restricted-lemma-composition-rig-flat}
$S$ をスキームとする。$f : X \to Y$ と $g : Y \to Z$ を，
$S$ 上の局所 Noether 形式代数空間の射とする。
$f$ と $g$ が rig-平坦ならば，$g \circ f$ も rig-平坦である。
\end{lemma}

\begin{proof}
Formal Spaces, Remark
\ref{formal-spaces-remark-composition-variant-Noetherian}
により，これは Lemma
\ref{restricted-lemma-composition-rig-flat-continuous} から従う。
\end{proof}

\section{Rig-smooth 準同型}
\label{restricted-section-rig-smooth-homomorphisms}

\noindent
% P10-E0680
この節では，局所 Noether 形式代数空間の rig-smooth 射を導入するために
必要となる，adic Noether 位相環の rig-smooth 準同型の性質をいくつか証明する。

\begin{lemma}
\label{restricted-lemma-rig-smooth-continuous}
$A \to B$ を $\textit{WAdm}^{Noeth}$ の射とする
（Formal Spaces, Section \ref{formal-spaces-section-morphisms-rings}）。
次の条件は同値である：
\begin{enumerate}
% P10-E0681
\item[(a)] $A \to B$ は Lemma \ref{restricted-lemma-finite-type} の
同値な条件を満たし，定義イデアル $I \subset B$ で，
$B$ が $(A, I)$ 上 rig-smooth となるものが存在する，および
\item[(b)] $A \to B$ は Lemma \ref{restricted-lemma-finite-type} の
同値な条件を満たし，任意の定義イデアル $I \subset A$ に対して，
代数 $B$ は $(A, I)$ 上 rig-smooth である。
\end{enumerate}
\end{lemma}

\begin{proof}
% P10-E0682
$I$ と $I'$ を $A$ の定義イデアルとする。このとき，
$I^c \subset I'$ かつ $(I')^c \subset I$ となる整数 $c \geq 0$ が存在する。
従って，$B$ が $(A, I)$ 上 rig-smooth であることと，
$B$ が $(A, I')$ 上 rig-smooth であることは同値である。
これは Definition \ref{restricted-definition-rig-smooth-homomorphism}，
包含 $I^c \subset I'$ と $(I')^c \subset I$，および Remark
\ref{restricted-remark-NL-well-defined-topological} により素朴な余接複体
$\NL_{B/A}^\wedge$ が $A$ の定義イデアルの選択によらないことから従う。
\end{proof}

\begin{definition}
\label{restricted-definition-rig-smooth-continuous-homomorphism}
$\varphi : A \to B$ を adic Noether 位相環の間の連続環準同型とする。
すなわち，$\varphi$ は $\textit{WAdm}^{Noeth}$ の射である。
Lemma \ref{restricted-lemma-rig-smooth-continuous} の同値な条件が成り立つとき，
$\varphi$ は {it rig-smooth} であるという。
\end{definition}

\noindent
これは局所的性質を定める。

\begin{lemma}
\label{restricted-lemma-rig-smooth-axioms}
$\textit{WAdm}^{Noeth}$ の射上の性質
$P(\varphi)=$「$\varphi$ は rig-smooth である」は，Formal Spaces, Remark
\ref{formal-spaces-remark-variant-Noetherian}
で定義される局所的性質である。
\end{lemma}

\begin{proof}
この主張が何を意味するかを思い出そう。まず，
$\textit{WAdm}^{Noeth}$ は，対象が adic Noether 位相環，
射が連続環準同型である圏である。次の条件を満たす可換図式
$$
\xymatrix{
B \ar[r] & (B')^\wedge \\
A \ar[r] \ar[u]^\varphi & (A')^\wedge \ar[u]_{\varphi'}
}
$$
を考える：
$A$ と $B$ は adic Noether 位相環であり，
$A \to A'$ と $B \to B'$ は \'etale 環写像であり，
ある定義イデアル $I \subset A$ に対して
$(A')^\wedge = \lim A'/I^nA'$ であり，
ある定義イデアル $J \subset B$ に対して
$(B')^\wedge = \lim B'/J^nB'$ であり，
$\varphi : A \to B$ と
$\varphi' : (A')^\wedge \to (B')^\wedge$ は連続である。
Formal Spaces, Lemma
\ref{formal-spaces-lemma-completion-in-sub} により，
$(A')^\wedge$ と $(B')^\wedge$ は adic Noether 位相環であることに注意せよ。
次を示さなければならない。
\begin{enumerate}
\item $\varphi$ が rig-smooth $\Rightarrow \varphi'$ は rig-smooth である，
\item $B \to B'$ が忠実平坦ならば，
$\varphi'$ が rig-smooth $\Rightarrow \varphi$ は rig-smooth である，および
\item $i = 1, \ldots, n$ に対して $A \to B_i$ が rig-smooth ならば，
$A \to \prod_{i = 1, \ldots, n} B_i$ は rig-smooth である。
\end{enumerate}
Lemma \ref{restricted-lemma-finite-type} の同値な条件は，条件 (1)，(2)，(3) を満たす。
従って性質「rig-smooth」について (1)，(2)，(3) を確認するときには，
いずれの場合も環写像が Lemma
\ref{restricted-lemma-finite-type} の同値な条件を満たすと既に仮定してよい。

\medskip\noindent
定義イデアル $I \subset A$ を取る。上の注意により，図式中の各環の位相は
$I$-進位相であり，$B$，$(A')^\wedge$，$(B')^\wedge$ は
$(A, I)$ に対する圏（\ref{restricted-equation-C-prime}）に属する。
$A \to A'$ と $B \to B'$ は \'etale なので，複体
$\NL_{A'/A}$ と $\NL_{B'/B}$ は零であり，従って Lemma
\ref{restricted-lemma-NL-is-completion} により
$\NL_{(A')^\wedge/A}^\wedge$ と
$\NL_{(B')^\wedge/B}^\wedge$ も零である。
Lemma \ref{restricted-lemma-exact-sequence-NL} を
$A \to (A')^\wedge \to (B')^\wedge$ に適用すると，同型
$$
H^i(\NL_{(B')^\wedge/(A')^\wedge}^\wedge) \to H^i(\NL_{(B')^\wedge/A}^\wedge)
$$
を得る。従って
% P10-E0683
$\NL_{(B')^\wedge/A}^\wedge \to \NL_{(B')^\wedge/(A')^\wedge}$
は擬同型である。環写像 $B/I^nB \to B'/I^nB'$ は \'etale であり，
従って局所完全交叉である
（Algebra, Lemma \ref{algebra-lemma-etale-standard-smooth}）。
従って Lemma \ref{restricted-lemma-exact-sequence-NL} と
\ref{restricted-lemma-transitive-lci-at-end} を
$A \to B \to (B')^\wedge$ に適用して，同型
$$
H^i(\NL_{B/A}^\wedge \otimes_B (B')^\wedge) \to
H^i(\NL_{(B')^\wedge/A}^\wedge)
$$
を得る。従って
$\NL_{B/A}^\wedge \otimes_B (B')^\wedge
\to \NL_{(B')^\wedge/A}^\wedge$
は擬同型である。これら二つの観察を合わせると，
$$
\NL_{(B')^\wedge/(A')^\wedge}^\wedge \cong
\NL_{B/A}^\wedge \otimes_B (B')^\wedge
$$
を $D((B')^\wedge)$ において得る。以上の準備を終え，実際の証明に入る。

\medskip\noindent
(1) の証明。$\varphi$ は rig-smooth であると仮定する。
任意の $B$-加群 $N$ に対して
$\Ext^1_B(\NL_{B/A}^\wedge, N)$ が $I^c$ で零化されるような
$c \geq 0$ が存在する。More on Algebra, Lemma
\ref{more-algebra-lemma-two-term-base-change} と
\ref{more-algebra-lemma-base-change-property-ext-1-annihilated}
により，この性質は $B \to (B')^\wedge$ による基底変換で保たれる。
従って，すべての $(B')^\wedge$-加群 $N$ に対して
$\Ext^1_{(B')^\wedge}(\NL_{(B')^\wedge/(A')^\wedge}^\wedge, N)$
は $I^c(A')^\wedge$ で零化される。これは $\varphi'$ が
rig-smooth であることを述べる。これで (1) が証明された。

\medskip\noindent
(2) を証明するため，$B \to B'$ は忠実平坦であり，
$\varphi'$ は rig-smooth であると仮定する。
任意の $(B')^\wedge$-加群 $N'$ に対して
$\Ext^1_{(B')^\wedge}(\NL_{(B')^\wedge/(A')^\wedge}^\wedge, N')$
が $I^c(B')^\wedge$ で零化されるような $c \geq 0$ が存在する。
合成 $B \to B' \to (B')^\wedge$ は平坦である
（Algebra, Lemma \ref{algebra-lemma-completion-flat}）。
従って任意の $B$-加群 $N$ に対して
$$
\Ext^1_B(\NL_{B/A}^\wedge, N) \otimes_B (B')^\wedge =
\Ext^1_{(B')^\wedge}(\NL_{B/A}^\wedge \otimes_B (B')^\wedge,
N \otimes_B (B')^\wedge)
$$
が More on Algebra, Lemma
\ref{more-algebra-lemma-base-change-RHom} の part (3) により成り立つ
（細部の一部は省略する）。従って，この加群は $I^c$ で零化される。
しかし，$B \to B'$ が忠実平坦であるという仮定により，
$B \to (B')^\wedge$ は実際に忠実平坦である
（Formal Spaces, Lemma
\ref{formal-spaces-lemma-etale-surjective}）。従って
$\Ext^1_B(\NL_{B/A}^\wedge, N)$ は $I^c$ で零化される。
故に $\varphi$ は rig-smooth である。これで (2) が証明された。

\medskip\noindent
(3) を証明するために $B = \prod_{i = 1, \ldots, n} B_i$ とおくと，
$\NL_{B/A}^\wedge$ は，$B$-加群の複体とみなした
複体 $\NL_{B_i/A}^\wedge$ の直和であることを注意するだけでよい。
\end{proof}

\begin{lemma}
\label{restricted-lemma-base-change-rig-smooth-continuous}
$\textit{WAdm}^{Noeth}$ の射上の性質
$P(\varphi)=$「$\varphi$ は rig-smooth である」と
$Q(\varphi)$=「$\varphi$ は adic である」を考える。
このとき $P$ は，Formal Spaces, Remark
\ref{formal-spaces-remark-base-change-variant-variant-Noetherian}
で定義される意味で $Q$ による基底変換について安定である。
\end{lemma}

\begin{proof}
Lemma \ref{restricted-lemma-rig-smooth-continuous} により，この主張には意味がある。
真であることを見るため，$\textit{WAdm}^{Noeth}$ の射
$B \to A$ と $B \to C$ が与えられ，
$B \to A$ は rig-smooth，$B \to C$ は adic であると仮定する
（Formal Spaces, Definition
\ref{formal-spaces-definition-adic-homomorphism}）。
$A$ と $C$ の位相が $I$-進位相となるような定義イデアル
$I \subset B$ を選べる。この状況では Lemma
\ref{restricted-lemma-base-change-rig-smooth-homomorphism} により，
$A \widehat{\otimes}_B C$ は $(C, IC)$ 上 rig-smooth である。
\end{proof}

\begin{lemma}
\label{restricted-lemma-composition-rig-smooth-continuous}
$\textit{WAdm}^{Noeth}$ の射上の性質
$P(\varphi)=$「$\varphi$ は rig-smooth である」は，Formal Spaces, Remark
\ref{formal-spaces-remark-composition-variant-Noetherian}
で定義される意味で合成について安定である。
\end{lemma}

\begin{proof}
読者には，以下の証明を読むのではなく自分自身の証明を見つけることを強く勧める。
Lemma \ref{restricted-lemma-rig-smooth-continuous} により，この主張には意味がある。
真であることを見るため，$\textit{WAdm}^{Noeth}$ における rig-smooth 射
$A \to B$ と $B \to C$ が与えられたと仮定する。
$C$ と $B$ の位相が $I$-進位相となるような定義イデアル
$I \subset A$ を選べる。Lemma \ref{restricted-lemma-exact-sequence-NL} により，
完全列
$$
\xymatrix{
C \otimes_B H^0(\NL_{B/A}^\wedge) \ar[r] &
H^0(\NL_{C/A}^\wedge) \ar[r] &
H^0(\NL_{C/B}^\wedge) \ar[r] & 0 \\
H^{-1}(\NL_{B/A}^\wedge \otimes_B C) \ar[r] &
H^{-1}(\NL_{C/A}^\wedge) \ar[r] &
H^{-1}(\NL_{C/B}^\wedge) \ar[llu]
}
$$
を得る。
$H^{-1}(\NL_{B/A}^\wedge \otimes_B C)$ と
$H^{-1}(\NL_{C/B}^\wedge)$ は $I$ のある冪で零化されることに注意せよ。
これは Lemma \ref{restricted-lemma-equivalent-with-artin-smooth} の part (2) と
More on Algebra, Lemma
\ref{more-algebra-lemma-two-term-base-change} および
\ref{more-algebra-lemma-base-change-property-ext-1-annihilated}
を組み合わせれば従う（$B \to C$ による基底変換を扱うためである）。
従って $H^{-1}(\NL_{C/A}^\wedge)$ は $I$ のある冪で零化される。
次に，Lemma \ref{restricted-lemma-equivalent-with-artin-smooth} の part (2) における
rig-smooth 代数の特徴づけ（これはさらに More on Algebra, Lemma
\ref{more-algebra-lemma-ext-1-annihilated} の part (5) を参照する）により，
$f_1, \ldots, f_s \in IB$ と $g_1, \ldots, g_t \in IC$ で，
$V(f_1, \ldots, f_s) = V(IB)$ および
$V(g_1, \ldots, g_t) = V(IC)$ となり，しかも
$H^0(\NL_{B/A}^\wedge)_{f_i}$ が有限射影 $B_{f_i}$-加群，
$H^0(\NL_{C/B}^\wedge)_{g_j}$ が有限射影 $C_{g_j}$-加群となるものを選べる。
次数 $-1$ のコホモロジーは $f_ig_j$ における局所化で消えるので，
短完全列
$$
0 \to
(C \otimes_B H^0(\NL_{B/A}^\wedge))_{f_ig_j} \to
H^0(\NL_{C/A}^\wedge)_{f_ig_j} \to
H^0(\NL_{C/B}^\wedge)_{f_ig_j} \to 0
$$
を得る。従って
% P10-E0684
$H^0(\NL_{C/A}^\wedge)_{f_ig_j}$ は，同じものの拡大として
有限射影 $C_{f_ig_j}$-加群である。
従って Lemma \ref{restricted-lemma-equivalent-with-artin-smooth} の part (2) の判定法と，
% P10-E0685
それを通じて More on Algebra, Lemma
\ref{more-algebra-lemma-ext-1-annihilated} の part (4) の判定法により，
$C$ は $(A, I)$ 上 rig-smooth であると結論する。
\end{proof}

\noindent
次の補題は，rig-smooth 準同型が「rig-syntomic」または
「rig-flat$+$rig-lci」であると述べるものと解釈できる。

\begin{lemma}
\label{restricted-lemma-rig-smooth-rig-flat}
$\varphi : A \to B$ を $\textit{WAdm}^{Noeth}$ の射とする。
$\varphi$ が rig-smooth ならば $\varphi$ は rig-平坦であり，
任意の表示 $B = A\{x_1, \ldots, x_n\}/J$ と，
定義イデアルを含まない素イデアル
$J \subset \mathfrak q \subset A\{x_1, \ldots, x_n\}$ に対して，
イデアル $J_\mathfrak q \subset A\{x_1, \ldots, x_n\}_\mathfrak q$ は
正則列で生成される。
\end{lemma}

\begin{proof}
$f \in A$ とする。$\varphi$ が rig-平坦であることを示すには，
$\varphi_{\{f\}} : A_{\{f\}} \to B_{\{f\}}$ が
素朴に rig-平坦であることを示さなければならない。
$\varphi_{\{f\}}$ を $\varphi$ の基底変換とみなし，Lemma
\ref{restricted-lemma-base-change-rig-smooth-continuous} を使うか，または
rig-smooth であることが局所的性質であること
（Lemma \ref{restricted-lemma-rig-smooth-axioms}）を使えば，
$\varphi_{\{f\}}$ は rig-smooth であることが分かる。
従って $\varphi$ が素朴に rig-平坦であることを示せば十分である。

\medskip\noindent
表示 $B = A\{x_1, \ldots, x_n\}/J$ を選ぶ。
補題の第2の部分を確認するには，定義イデアルを含まないという性質に関して
極大な $\mathfrak q$ に対し，
% P10-E0686
$J \subset \mathfrak q$ について
$J_\mathfrak q \subset A\{x_1, \ldots, x_n\}_\mathfrak q$ が
正則列で生成されることを確認すれば十分である；Algebra, Lemma
\ref{algebra-lemma-regular-sequence-in-neighbourhood} を参照
（これは，$V(J)$ の素イデアルのうち $J$ を生成する正則列が存在するものの
集合が開であることを示す）。
言い換えると，$\mathfrak q$ は $A\{x_1, \ldots, x_n\}$ で
rig-閉であると仮定してよい。また $B$ が素朴に rig-平坦であることを
確認するにも，対応する局所化 $B_\mathfrak q$ が $A$ 上平坦であることを
確認すれば十分である。

\medskip\noindent
$\mathfrak q \subset A\{x_1, \ldots, x_n\}$ を
$J \subset \mathfrak q$ を満たす rig-閉素イデアルとする。
Lemma \ref{restricted-lemma-rig-closed-point-after-localization} により，
$A\{x_1, \ldots, x_n\}/\mathfrak q$ で単元に写る $f \in A$ で，
誘導される $A_{\{f\}}$ の素イデアル $\mathfrak p'$ が rig-閉となるものを
選べる。以下では
$A_{\{f\}}\{x_1, \ldots, x_n\} = A\{x_1, \ldots, x_n\}_{\{f\}}$
を用いる。詳細は省略する。図式
$$
\xymatrix{
A\{x_1, \ldots, x_n\}_{\mathfrak q} / J_\mathfrak q \ar[r] &
A_{\{f\}}\{x_1, \ldots, x_n\}_{\mathfrak q'}/
J A_{\{f\}}\{x_1, \ldots, x_n\}_{\mathfrak q'} \\
A\{x_1, \ldots, x_n\}_{\mathfrak q} \ar[r] \ar[u] &
A_{\{f\}}\{x_1, \ldots, x_n\}_{\mathfrak q'} \ar[u] \\
A \ar[r] \ar[u] &
A_{\{f\}} \ar[u]
}
$$
を考える。中央の横射は，局所環の局所準同型であり，かつ
$A_{\{f\}}\{x_1, \ldots, x_n\}$ が Noether 環
$A\{x_1, \ldots, x_n\}$ の局所化の完備化なので平坦であることから，
忠実平坦である。同様に，下の横射は平坦である。
従って $J_\mathfrak q$ が正則列で生成され，
$A \to A\{x_1, \ldots, x_n\}_{\mathfrak q} / J_\mathfrak q$
が平坦であることを示すには，
$J A_{\{f\}}\{x_1, \ldots, x_n\}_{\mathfrak q'}$ についての
同じ主張と，写像
$A_{\{f\}} \to A_{\{f\}}\{x_1, \ldots, x_n\}_{\mathfrak q'}/
J A_{\{f\}}\{x_1, \ldots, x_n\}_{\mathfrak q'}$
についての同じ主張を証明すれば十分である。
正則列に関する主張については，Algebra, Lemma
\ref{algebra-lemma-flat-increases-depth} または More on Algebra, Lemma
\ref{more-algebra-lemma-flat-descent-regular-ideal} を参照せよ。
最後に，$A_{\{f\}} \to B_{\{f\}}$ が rig-smooth であることは
既に見た。これは次の段落で論じる場合への帰着を与える。

\medskip\noindent
$\mathfrak q \subset A\{x_1, \ldots, x_n\}$ を
$J \subset \mathfrak q$ を満たす rig-閉素イデアルとし，さらに
$\mathfrak p = A \cap \mathfrak q$ も rig-閉であるとする。
Lemma \ref{restricted-lemma-equivalent-with-artin-smooth} による
rig-smooth 代数の特徴づけにより，変数 $x_1, \ldots, x_n$ を並べ替えれば，
$m \geq 0$ と $f_1, \ldots, f_m \in J$ で次を満たすものを取れる。
\begin{enumerate}
\item $J_\mathfrak q$ は $f_1, \ldots, f_m$ で生成される，および
\item $\det_{1 \leq i, j \leq m}(\partial f_j/ \partial x_i)$ は
$A\{x_1, \ldots, x_n\}_\mathfrak q$ の単元に写る。
\end{enumerate}
Lemma \ref{restricted-lemma-fibre-regular} により，ファイバー環
$$
F = A\{x_1, \ldots, x_n\} \otimes_A \kappa(\mathfrak p)
$$
は正則である。$A$-導分 $\partial / \partial x_i$ は
（一意的に）導分 $D_i : F \to F$ に延長されることに注意せよ。
More on Algebra, Lemma
\ref{more-algebra-lemma-quotient-sequence-regular} により，
$f_1, \ldots, f_m$ は $F_\mathfrak q$ の正則列に写る。
$A \to A\{x_1, \ldots, x_n\}$ の平坦性と Algebra, Lemma
\ref{algebra-lemma-grothendieck-regular-sequence} により，これは
% P10-E0687
$f_1, \ldots, f_m$ が
$A\{x_1, \ldots, x_m\}_\mathfrak q$ の正則列に写り，
これらの元による商が $A$ 上平坦であることを示す。これで証明が完了する。
\end{proof}

\begin{lemma}
\label{restricted-lemma-exact-sequence-NL-rig-smooth}
$A \to B \to C$ を，adic かつ位相的有限型である
$\textit{WAdm}^{Noeth}$ の射とする。
$B \to C$ が rig-smooth ならば，写像
$$
H^{-1}(\NL_{B/A}^\wedge \otimes_B C) \to H^{-1}(\NL_{C/A}^\wedge)
$$
の核（Lemma \ref{restricted-lemma-exact-sequence-NL} を参照）は，
ある定義イデアルで零化される。
\end{lemma}

\begin{proof}
$\overline{\mathfrak q} \subset C$ を，定義イデアルを含まない
素イデアルとする。問題の加群は有限なので，写像
$$
H^{-1}(\NL_{B/A}^\wedge \otimes_B C)_{\overline{\mathfrak q}} \to
H^{-1}(\NL_{C/A}^\wedge)_{\overline{\mathfrak q}}
$$
が単射であることを示せば十分である。
Lemma \ref{restricted-lemma-exact-sequence-NL} の証明と同様に，表示
$B = A\{x_1, \ldots, x_r\}/J$，
$C = B\{y_1, \ldots, y_s\}/J'$，および
$C = A\{x_1, \ldots, x_r, y_1, \ldots, y_s\}/K$
を選ぶ。Lemma \ref{restricted-lemma-exact-sequence-NL} の証明中の図式を見ると，
$J/J^2 \otimes_B C \to K/K^2$ が，
$\overline{\mathfrak q}$ に対応する素イデアル
$\mathfrak q \subset A\{x_1, \ldots, x_r, y_1, \ldots, y_s\}$
における局所化の後で単射であることを示せば十分である。
More on Algebra, Lemma
\ref{more-algebra-lemma-transitive-lci-at-end} とその証明を比較されたい。
これは $J/KJ \to K/K^2$ が $\mathfrak q$ における局所化の後で
単射となることと同じである。同値なことに，
$J_\mathfrak q \cap K^2_\mathfrak q = (KJ)_\mathfrak q$
を示さなければならない。
Lemma \ref{restricted-lemma-rig-smooth-rig-flat} により，
$(K/J)_\mathfrak q = J'_\mathfrak q$ は正則列で生成される。
従って，望む交叉の性質は More on Algebra, Lemma
\ref{more-algebra-lemma-conormal-sequence-H1-regular-ideal}
から従う（正則列で生成されるイデアルは $H_1$-正則であることも用いる；
More on Algebra, Section \ref{more-algebra-section-ideals} を参照）。
\end{proof}

\section{Rig-smooth 射}
\label{restricted-section-rig-smooth-morphisms}

\noindent
% P10-E0688
この節では Section \ref{restricted-section-rig-smooth-homomorphisms} の結果を用いて，
局所 Noether 代数空間の rig-smooth 射を定義する。

\begin{definition}
\label{restricted-definition-rig-smooth}
$S$ をスキームとする。$f : X \to Y$ を，$S$ 上の
局所 Noether 形式代数空間の射とする。
アフィン形式代数空間 $U$，$V$ をもつ任意の可換図式
$$
\xymatrix{
U \ar[d] \ar[r] & V \ar[d] \\
X \ar[r] & Y
}
$$
で，$U \to X$ と $V \to Y$ が代数空間により表現可能かつ
\'etale であるものに対し，射 $U \to V$ が adic Noether 位相環の
rig-smooth 写像に対応するとき，$f$ は {it rig-smooth} であるという。
\end{definition}

\noindent
この条件を始域と終域上で \'etale 局所的に検査できることを証明しよう。

\begin{lemma}
\label{restricted-lemma-rig-smooth-morphisms}
$S$ をスキームとする。$f : X \to Y$ を，$S$ 上の
局所 Noether 形式代数空間の射とする。次の条件は同値である。
\begin{enumerate}
\item $f$ は rig-smooth である，
\item アフィン形式代数空間 $U$，$V$ をもつ任意の可換図式
$$
\xymatrix{
U \ar[d] \ar[r] & V \ar[d] \\
X \ar[r] & Y
}
$$
で，$U \to X$ と $V \to Y$ が代数空間により表現可能かつ
\'etale であるものに対し，射 $U \to V$ は
$\textit{WAdm}^{Noeth}$ の rig-smooth 写像に対応する，
\item Formal Spaces, Definition
\ref{formal-spaces-definition-formal-algebraic-space}
の意味での被覆 $\{Y_j \to Y\}$ が存在し，各 $j$ に対して
Formal Spaces, Definition
\ref{formal-spaces-definition-formal-algebraic-space}
の意味での被覆 $\{X_{ji} \to Y_j \times_Y X\}$ が存在して，
各 $X_{ji} \to Y_j$ は $\textit{WAdm}^{Noeth}$ の
rig-smooth 写像に対応する，および
\item Formal Spaces, Definition
\ref{formal-spaces-definition-formal-algebraic-space}
の意味での被覆 $\{X_i \to X\}$ が存在し，各 $i$ に対して
分解 $X_i \to Y_i \to Y$ が存在する。ここで $Y_i$ は
アフィン形式代数空間，$Y_i \to Y$ は代数空間により表現可能かつ
\'etale であり，$X_i \to Y_i$ は $\textit{WAdm}^{Noeth}$ の
rig-smooth 写像に対応する。
\end{enumerate}
\end{lemma}

\begin{proof}
(1) と (2) の同値性は Definition \ref{restricted-definition-rig-smooth} である。
(2)，(3)，(4) の同値性は，Lemma
\ref{restricted-lemma-rig-smooth-axioms} により rig-smooth であることが
$\text{WAdm}^{Noeth}$ の射の局所的性質であること，および
% P10-E0688
局所 Noether 代数空間の間の射に対する Formal Spaces, Lemma
\ref{formal-spaces-lemma-property-defines-property-morphisms}
の変種を適用することから従う。この変種は Formal Spaces, Remark
\ref{formal-spaces-remark-variant-Noetherian}
で述べられている。
\end{proof}

\begin{lemma}
\label{restricted-lemma-base-change-rig-smooth}
$S$ をスキームとする。$f : X \to Y$ と $g : Z \to Y$ を，
$S$ 上の局所 Noether 形式代数空間の射とする。
$f$ が rig-smooth で $g$ が adic ならば，基底変換
$X \times_Y Z \to Z$ は rig-smooth である。
\end{lemma}

\begin{proof}
Formal Spaces, Remark
\ref{formal-spaces-remark-base-change-variant-variant-Noetherian}
および Formal Spaces, Section
\ref{formal-spaces-section-adic} の議論により，これは Lemma
\ref{restricted-lemma-base-change-rig-smooth-continuous} から従う。
\end{proof}

\begin{lemma}
\label{restricted-lemma-composition-rig-smooth}
$S$ をスキームとする。$f : X \to Y$ と $g : Y \to Z$ を，
$S$ 上の局所 Noether 形式代数空間の射とする。
$f$ と $g$ が rig-smooth ならば，$g \circ f$ も rig-smooth である。
\end{lemma}

\begin{proof}
Formal Spaces, Remark
\ref{formal-spaces-remark-composition-variant-Noetherian}
により，これは Lemma
\ref{restricted-lemma-composition-rig-smooth-continuous} から従う。
\end{proof}

\begin{lemma}
\label{restricted-lemma-rig-smooth-rig-flat-morphism}
$S$ をスキームとする。$f : X \to Y$ を，$S$ 上の
局所 Noether 形式代数空間の射とする。
$f$ が rig-smooth ならば，$f$ は rig-flat である。
\end{lemma}

\begin{proof}
Lemma \ref{restricted-lemma-rig-smooth-rig-flat} と定義から直ちに従う。
\end{proof}

\section{Rig-\'etale 準同型}
\label{restricted-section-rig-etale-homomorphisms}

\noindent
% P10-E0689
本節では，局所 Noether 代数空間の rig-\'etale 射を導入するために必要となる，
adic Noether 位相環の rig-\'etale 準同型の性質をいくつか証明する。

\begin{lemma}
\label{restricted-lemma-rig-etale-continuous}
$A \to B$ を $\textit{WAdm}^{Noeth}$ の射とする
（Formal Spaces, Section \ref{formal-spaces-section-morphisms-rings}）。
次の条件は同値である。
\begin{enumerate}
\item[(a)] $A \to B$ は Lemma \ref{restricted-lemma-finite-type} の同値な条件を満たし，
% P10-E0690
定義イデアル $I \subset B$ が存在して，$B$ は $(A, I)$ 上 rig-\'etale である，および
\item[(b)] $A \to B$ は Lemma \ref{restricted-lemma-finite-type} の同値な条件を満たし，
すべての定義イデアル $I \subset A$ に対して，代数 $B$ は
$(A, I)$ 上 rig-\'etale である。
\end{enumerate}
\end{lemma}

\begin{proof}
% P10-E0691
$I$ と $I'$ を $A$ の定義イデアルとする。このとき，ある整数 $c \geq 0$ が存在して
$I^c \subset I'$ かつ $(I')^c \subset I$ となる。従って，$B$ が
$(A, I)$ 上 rig-\'etale であることと，$B$ が $(A, I')$ 上
rig-\'etale であることは同値である。
これは Definition \ref{restricted-definition-rig-etale-homomorphism}，包含
$I^c \subset I'$ と $(I')^c \subset I$，および Remark
\ref{restricted-remark-NL-well-defined-topological} により素朴な余接複体
$\NL_{B/A}^\wedge$ が $A$ の定義イデアルの選択に依存しないことから従う。
\end{proof}

\begin{definition}
\label{restricted-definition-rig-etale-continuous-homomorphism}
$\varphi : A \to B$ を adic Noether 位相環の間の連続環準同型とする。すなわち，
$\varphi$ は $\textit{WAdm}^{Noeth}$ の射である。Lemma
\ref{restricted-lemma-rig-etale-continuous} の同値な条件が成り立つとき，
$\varphi$ は {it rig-etale} であるという。
\end{definition}

\noindent
これは局所的性質を定める。

\begin{lemma}
\label{restricted-lemma-rig-etale-axioms}
$\textit{WAdm}^{Noeth}$ の射上の性質
$P(\varphi)=$``$\varphi$ は rig-\'etale である'' は，Formal Spaces, Remark
\ref{formal-spaces-remark-variant-Noetherian} で定義された意味で局所的性質である。
\end{lemma}

\begin{proof}
この証明は Lemma \ref{restricted-lemma-rig-smooth-axioms} の証明とまったく同じである。
主張の意味を振り返ろう。まず，$\textit{WAdm}^{Noeth}$ は，対象が
adic Noether 位相環，射が連続環準同型である圏である。次の可換図式を考える。
$$
\xymatrix{
B \ar[r] & (B')^\wedge \\
A \ar[r] \ar[u]^\varphi & (A')^\wedge \ar[u]_{\varphi'}
}
$$
これは次の条件を満たすとする。$A$ と $B$ は adic Noether 位相環，
$A \to A'$ と $B \to B'$ は \'etale 環写像，ある定義イデアル
$I \subset A$ に対して $(A')^\wedge = \lim A'/I^nA'$，ある定義イデアル
$J \subset B$ に対して $(B')^\wedge = \lim B'/J^nB'$ であり，
$\varphi : A \to B$ と $\varphi' : (A')^\wedge \to (B')^\wedge$ は連続である。
Formal Spaces, Lemma \ref{formal-spaces-lemma-completion-in-sub} により，
$(A')^\wedge$ と $(B')^\wedge$ は adic Noether 位相環であることに注意せよ。
示すべきことは次である。
\begin{enumerate}
\item $\varphi$ が rig-\'etale $\Rightarrow \varphi'$ は rig-\'etale である，
\item $B \to B'$ が忠実平坦ならば，$\varphi'$ が rig-\'etale
$\Rightarrow \varphi$ は rig-\'etale である，および
\item $i = 1, \ldots, n$ に対して $A \to B_i$ が rig-\'etale ならば，
$A \to \prod_{i = 1, \ldots, n} B_i$ は rig-\'etale である。
\end{enumerate}
Lemma \ref{restricted-lemma-finite-type} の同値な条件は (1)，(2)，(3) を満たす。
従って性質 ``rig-\'etale'' について (1)，(2)，(3) を確認する際には，
いずれの場合も環写像が Lemma \ref{restricted-lemma-finite-type} の同値な条件を満たすと
あらかじめ仮定してよい。

\medskip\noindent
定義イデアル $I \subset A$ を選ぶ。上の注意により，図式中の各環の位相は
$I$-進位相であり，$B$，$(A')^\wedge$，$(B')^\wedge$ は
$(A, I)$ に対する圏 (\ref{restricted-equation-C-prime}) に属する。
$A \to A'$ と $B \to B'$ は \'etale なので，複体
$\NL_{A'/A}$ と $\NL_{B'/B}$ は零であり，従って Lemma
\ref{restricted-lemma-NL-is-completion} により
$\NL_{(A')^\wedge/A}^\wedge$ と $\NL_{(B')^\wedge/B}^\wedge$ は零である。
$A \to (A')^\wedge \to (B')^\wedge$ に Lemma
\ref{restricted-lemma-exact-sequence-NL} を適用すると，同型
$$
H^i(\NL_{(B')^\wedge/(A')^\wedge}^\wedge) \to H^i(\NL_{(B')^\wedge/A}^\wedge)
$$
を得る。従って
% P10-E0692
$\NL_{(B')^\wedge/A}^\wedge \to \NL_{(B')^\wedge/(A')^\wedge}$
は擬同型である。環写像 $B/I^nB \to B'/I^nB'$ は \'etale であり，従って
局所完全交叉である（Algebra, Lemma
\ref{algebra-lemma-etale-standard-smooth}）。従って
$A \to B \to (B')^\wedge$ に Lemmas
\ref{restricted-lemma-exact-sequence-NL} と \ref{restricted-lemma-transitive-lci-at-end} を適用でき，同型
$$
H^i(\NL_{B/A}^\wedge \otimes_B (B')^\wedge) \to
H^i(\NL_{(B')^\wedge/A}^\wedge)
$$
を得る。従って
$\NL_{B/A}^\wedge \otimes_B (B')^\wedge \to \NL_{(B')^\wedge/A}^\wedge$
は擬同型である。これら二つの観察を合わせると，$D((B')^\wedge)$ において
$$
\NL_{(B')^\wedge/(A')^\wedge}^\wedge \cong
\NL_{B/A}^\wedge \otimes_B (B')^\wedge
$$
を得る。準備が整ったので，実際の証明を始める。

\medskip\noindent
(1) の証明。$\varphi$ が rig-\'etale であると仮定する。このとき，ある
$c \geq 0$ が存在し，$a \in I^c$ による乗法が $D(B)$ において
$\NL_{B/A}^\wedge$ 上零となる。この性質は $B \to (B')^\wedge$ による
基底変換で保たれる。More on Algebra, Lemmas
\ref{more-algebra-lemma-two-term-base-change} を参照せよ。上の同型により，
$\varphi'$ は rig-\'etale である。これで (1) が示された。

\medskip\noindent
(2) を示すため，$B \to B'$ は忠実平坦で，$\varphi'$ は rig-\'etale であると仮定する。
このとき，ある $c \geq 0$ が存在し，$a \in I^c$ による乗法が
$D((B')^\wedge)$ において
$\NL_{(B')^\wedge/(A')^\wedge}^\wedge$ 上零となる。
上の同型により，
% P10-E0693
$a^c$ は $\NL_{B/A}^\wedge \otimes_B (B')^\wedge$ のコホモロジー加群を零化する。
合成 $B \to (B')^\wedge$ は，$B \to B'$ が忠実平坦であるという仮定と
Formal Spaces, Lemma \ref{formal-spaces-lemma-etale-surjective} により忠実平坦である。
従って $\NL_{B/A}^\wedge$ のコホモロジー加群は $I^c$ により零化される。
Lemma \ref{restricted-lemma-equivalent-with-artin} から $\varphi$ は rig-\'etale である。
これで (2) が示された。

\medskip\noindent
(3) を示すには，$B = \prod_{i = 1, \ldots, n} B_i$ と置き，
$\NL_{B/A}^\wedge$ が $B$-加群の複体とみなした複体
$\NL_{B_i/A}^\wedge$ の直和であることを観察すればよい。
\end{proof}

\begin{lemma}
\label{restricted-lemma-base-change-rig-etale-continuous}
$\textit{WAdm}^{Noeth}$ の射上の性質
$P(\varphi)=$``$\varphi$ は rig-\'etale である'' と
$Q(\varphi)$=``$\varphi$ は adic である'' を考える。このとき $P$ は，
Formal Spaces, Remark
\ref{formal-spaces-remark-base-change-variant-variant-Noetherian}
で定義された意味で $Q$ による基底変換について安定である。
\end{lemma}

\begin{proof}
Lemma \ref{restricted-lemma-rig-etale-continuous} により主張には意味がある。
これが真であることを見るため，$\textit{WAdm}^{Noeth}$ の射
$B \to A$ と $B \to C$ があり，$B \to A$ は rig-\'etale，
$B \to C$ は adic であると仮定する（Formal Spaces, Definition
\ref{formal-spaces-definition-adic-homomorphism}）。
このとき，$A$ と $C$ の位相が $I$-進位相となるような定義イデアル
$I \subset B$ を選べる。この状況では，Lemma
\ref{restricted-lemma-base-change-rig-etale-homomorphism} により
$A \widehat{\otimes}_B C$ が $(C, IC)$ 上 rig-\'etale であることが直ちに従う。
\end{proof}

\begin{lemma}
\label{restricted-lemma-composition-rig-etale-continuous}
$\textit{WAdm}^{Noeth}$ の射上の性質
$P(\varphi)=$``$\varphi$ は rig-\'etale である'' は，Formal Spaces, Remark
\ref{formal-spaces-remark-composition-variant-Noetherian}
で定義された意味で合成について安定である。
\end{lemma}

\begin{proof}
Lemma \ref{restricted-lemma-rig-etale-continuous} により主張には意味がある。
これが真であることを見るため，$\textit{WAdm}^{Noeth}$ における
rig-\'etale 射 $A \to B$ と $B \to C$ があると仮定する。
このとき，$C$ と $B$ の位相が $I$-進位相となるような定義イデアル
$I \subset A$ を選べる。Lemma \ref{restricted-lemma-exact-sequence-NL} により，完全列
$$
\xymatrix{
C \otimes_B H^0(\NL_{B/A}^\wedge) \ar[r] &
H^0(\NL_{C/A}^\wedge) \ar[r] &
H^0(\NL_{C/B}^\wedge) \ar[r] & 0 \\
H^{-1}(\NL_{B/A}^\wedge \otimes_B C) \ar[r] &
H^{-1}(\NL_{C/A}^\wedge) \ar[r] &
H^{-1}(\NL_{C/B}^\wedge) \ar[llu]
}
$$
を得る。ある $c \geq 0$ が存在し，すべての $a \in I$ に対して，
$a^c$ による乗法は $D(B)$ において $\NL_{B/A}^\wedge$ 上零であり，
$D(C)$ において $\NL_{C/B}^\wedge$ 上零である。当然ながら，
$a^c$ による乗法は $D(C)$ において
$\NL_{B/A}^\wedge \otimes_B C$ 上も零である。従って
% P10-E0694
$H^0(\NL_{B/A}^\wedge) \otimes_A C$，
$H^0(\NL_{C/B}^\wedge)$，
$H^{-1}(\NL_{B/A}^\wedge \otimes_B C)$，および
$H^{-1}(\NL_{C/B}^\wedge)$ は $a^c$ により零化される。
完全列から，$a^{2c}$ による乗法は
$H^0(\NL_{C/A}^\wedge)$ と $H^{-1}(\NL_{C/A}^\wedge)$ 上零である。
Lemma \ref{restricted-lemma-equivalent-with-artin} により，望みどおり $C$ は
$(A, I)$ 上 rig-\'etale である。
\end{proof}

\begin{lemma}
\label{restricted-lemma-permanence-rig-etale-continuous}
$\textit{WAdm}^{Noeth}$ の射上の性質
$P(\varphi)=$``$\varphi$ は rig-\'etale である'' は，Formal Spaces, Remark
\ref{formal-spaces-remark-permanence-variant-Noetherian}
で定義された意味で消去性質をもつ。
\end{lemma}

\begin{proof}
Lemma \ref{restricted-lemma-rig-etale-continuous} により主張には意味がある。
これが真であることを見るため，$\textit{WAdm}^{Noeth}$ の写像
$A \to B$ と $B \to C$ があり，$A \to C$ と $A \to B$ は
rig-\'etale であると仮定する。$B \to C$ が rig-\'etale であることを示さねばならない。
このとき，$C$ と $B$ の位相が $I$-進位相となるような定義イデアル
$I \subset A$ を選べる。Lemma \ref{restricted-lemma-exact-sequence-NL} により，完全列
$$
\xymatrix{
C \otimes_B H^0(\NL_{B/A}^\wedge) \ar[r] &
H^0(\NL_{C/A}^\wedge) \ar[r] &
H^0(\NL_{C/B}^\wedge) \ar[r] & 0 \\
H^{-1}(\NL_{B/A}^\wedge \otimes_B C) \ar[r] &
H^{-1}(\NL_{C/A}^\wedge) \ar[r] &
H^{-1}(\NL_{C/B}^\wedge) \ar[llu]
}
$$
を得る。ある $c \geq 0$ が存在し，すべての $a \in I$ に対して，
$a^c$ による乗法は $D(B)$ において $\NL_{B/A}^\wedge$ 上零であり，
$D(C)$ において $\NL_{C/A}^\wedge$ 上零である。従って
% P10-E0695
$H^0(\NL_{B/A}^\wedge) \otimes_A C$，
$H^0(\NL_{C/A}^\wedge)$，および
$H^{-1}(\NL_{C/A}^\wedge)$ は $a^c$ により零化される。
完全列から，$a^{2c}$ による乗法は
$H^0(\NL_{C/B}^\wedge)$ と $H^{-1}(\NL_{C/B}^\wedge)$ 上零である。
Lemma \ref{restricted-lemma-equivalent-with-artin} により，望みどおり $C$ は
$(B, IB)$ 上 rig-\'etale である。
\end{proof}

\section{Rig-\'etale 射}
\label{restricted-section-rig-etale-morphisms}

\noindent
本節では，Section \ref{restricted-section-rig-etale-homomorphisms} で得た結果を用いて，
% P10-E0696
局所 Noether 代数空間の rig-\'etale 射を定義する。

\begin{definition}
\label{restricted-definition-rig-etale}
$S$ をスキームとする。$f : X \to Y$ を，$S$ 上の局所 Noether 形式代数空間の射とする。
アフィン形式代数空間 $U$，$V$ をもつ任意の可換図式
$$
\xymatrix{
U \ar[d] \ar[r] & V \ar[d] \\
X \ar[r] & Y
}
$$
で，$U \to X$ と $V \to Y$ が代数空間により表現可能かつ \'etale であるものに対し，
射 $U \to V$ が adic Noether 位相環の rig-\'etale 写像に対応するとき，
$f$ は {\it rig-\'etale} であるという。
\end{definition}

\noindent
この条件を始域と終域上で \'etale 局所的に検査できることを証明しよう。

\begin{lemma}
\label{restricted-lemma-rig-etale-morphisms}
$S$ をスキームとする。$f : X \to Y$ を，$S$ 上の
局所 Noether 形式代数空間の射とする。次の条件は同値である。
\begin{enumerate}
\item $f$ は rig-\'etale である，
\item アフィン形式代数空間 $U$，$V$ をもつ任意の可換図式
$$
\xymatrix{
U \ar[d] \ar[r] & V \ar[d] \\
X \ar[r] & Y
}
$$
で，$U \to X$ と $V \to Y$ が代数空間により表現可能かつ \'etale であるものに対し，
射 $U \to V$ は $\textit{WAdm}^{Noeth}$ の rig-\'etale 写像に対応する，
\item Formal Spaces, Definition
\ref{formal-spaces-definition-formal-algebraic-space}
の意味での被覆 $\{Y_j \to Y\}$ が存在し，各 $j$ に対して
Formal Spaces, Definition
\ref{formal-spaces-definition-formal-algebraic-space}
の意味での被覆 $\{X_{ji} \to Y_j \times_Y X\}$ が存在して，
各 $X_{ji} \to Y_j$ は $\textit{WAdm}^{Noeth}$ の
rig-\'etale 写像に対応する，および
\item Formal Spaces, Definition
\ref{formal-spaces-definition-formal-algebraic-space}
の意味での被覆 $\{X_i \to X\}$ が存在し，各 $i$ に対して
分解 $X_i \to Y_i \to Y$ が存在する。ここで $Y_i$ は
アフィン形式代数空間，$Y_i \to Y$ は代数空間により表現可能かつ
\'etale であり，$X_i \to Y_i$ は $\textit{WAdm}^{Noeth}$ の
rig-\'etale 写像に対応する。
\end{enumerate}
\end{lemma}

\begin{proof}
(1) と (2) の同値性は Definition \ref{restricted-definition-rig-etale} である。
(2)，(3)，(4) の同値性は，Lemma \ref{restricted-lemma-rig-etale-axioms} により
rig-\'etale であることが $\text{WAdm}^{Noeth}$ の射の局所的性質であること，および
% P10-E0696
局所 Noether 代数空間の間の射に対する Formal Spaces, Lemma
\ref{formal-spaces-lemma-property-defines-property-morphisms}
の変種を適用することから従う。この変種は Formal Spaces, Remark
\ref{formal-spaces-remark-variant-Noetherian} で述べられている。
\end{proof}

\noindent
なお，rig-\'etale 射は局所有限型である。

\begin{lemma}
\label{restricted-lemma-rig-etale-finite-type}
局所 Noether 形式代数空間の rig-\'etale 射は局所有限型である。
\end{lemma}

\begin{proof}
Lemma \ref{restricted-lemma-rig-etale-axioms} の性質 $P$ は，Formal Spaces, Lemma
\ref{formal-spaces-lemma-quotient-restricted-power-series}
の同値な条件 (a)，(b)，(c)，(d) を含意する。従ってこれは Formal Spaces, Lemma
\ref{formal-spaces-lemma-finite-type-local-property} から従う。
\end{proof}

\begin{lemma}
\label{restricted-lemma-rig-etale-rig-smooth-morphism}
局所 Noether 形式代数空間の rig-\'etale 射は rig-smooth である。
\end{lemma}

\begin{proof}
定義と Lemma \ref{restricted-lemma-rig-etale-rig-smooth} から従う。
\end{proof}

\begin{lemma}
\label{restricted-lemma-base-change-rig-etale}
$S$ をスキームとする。$f : X \to Y$ と $g : Z \to Y$ を，
$S$ 上の局所 Noether 形式代数空間の射とする。
$f$ が rig-\'etale で $g$ が adic ならば，基底変換
$X \times_Y Z \to Z$ は rig-\'etale である。
\end{lemma}

\begin{proof}
Formal Spaces, Remark
\ref{formal-spaces-remark-base-change-variant-variant-Noetherian}
および Formal Spaces, Section
\ref{formal-spaces-section-adic} の議論により，これは Lemma
\ref{restricted-lemma-base-change-rig-etale-continuous} から従う。
\end{proof}

\begin{lemma}
\label{restricted-lemma-composition-rig-etale}
$S$ をスキームとする。$f : X \to Y$ と $g : Y \to Z$ を，
$S$ 上の局所 Noether 形式代数空間の射とする。
$f$ と $g$ が rig-\'etale ならば，$g \circ f$ も rig-\'etale である。
\end{lemma}

\begin{proof}
Formal Spaces, Remark
\ref{formal-spaces-remark-composition-variant-Noetherian}
により，これは Lemma \ref{restricted-lemma-composition-rig-etale-continuous} から従う。
\end{proof}

\begin{lemma}
\label{restricted-lemma-rig-etale-permanence}
$S$ をスキームとする。$f : X \to Y$ と $g : Y \to Z$ を，
% P10-E0697
$S$ 上の局所 Noether 形式代数空間の射とする。
$g \circ f$ と $g$ が rig-\'etale ならば，$f$ も rig-\'etale である。
\end{lemma}

\begin{proof}
Formal Spaces, Remark
\ref{formal-spaces-remark-permanence-variant-Noetherian}
により，これは Lemma \ref{restricted-lemma-permanence-rig-etale-continuous} から従う。
\end{proof}

\begin{lemma}
\label{restricted-lemma-rig-etale-alternative-permanence}
$S$ をスキームとする。$f : X \to Y$ と $g : Y \to Z$ を，
$S$ 上の局所 Noether 形式代数空間の射とする。
$g \circ f$ が rig-\'etale で $g$ が adic モノ射ならば，
$f$ は rig-\'etale である。
\end{lemma}

\begin{proof}
Lemma \ref{restricted-lemma-base-change-rig-etale} と，$f$ が $g$ による
$g \circ f$ の基底変換であることを用いよ。
\end{proof}

\begin{lemma}
\label{restricted-lemma-closed-immersion-rig-smooth}
% P10-E0698
$S$ をスキームとする。$f : X \to Y$ を形式代数空間の射とする。
$X$ と $Y$ は局所 Noether で，$f$ は閉埋込みであると仮定する。
次の条件は同値である。
\begin{enumerate}
\item $f$ は rig-smooth である，
\item $f$ は rig-\'etale である，
\item 任意のアフィン形式代数空間 $V$ と，代数空間により表現可能かつ
\'etale である任意の射 $V \to Y$ に対して，射 $X \times_Y V \to V$ は
$\textit{WAdm}^{Noeth}$ の全射 $B \to A$ で，その核 $J$ が次の性質をもつものに対応する：
$B$ のある定義イデアル $I$ に対して $I(J/J^2) = 0$ である。
\end{enumerate}
\end{lemma}

\begin{proof}
% P10-E0699
(2) におけるような $V$ と $V \to Y$ が与えられたとき，$f$ にこれ以上の仮定を
課さなくても，Formal Spaces, Lemma
\ref{formal-spaces-lemma-closed-immersion-into-countably-indexed}
により，射 $X \times_Y V \to V$ は $\textit{WAdm}^{Noeth}$ の全射
$B \to A$ に対応することをまず観察しよう。

\medskip\noindent
Lemma \ref{restricted-lemma-rig-etale-rig-smooth-morphism} により
(2) $\Rightarrow$ (1) である。

\medskip\noindent
(3) $\Rightarrow$ (2) の証明。(3) を仮定する。Lemma
\ref{restricted-lemma-rig-etale-morphisms} により，(3) に現れる環写像
$B \to A$ が Definition \ref{restricted-definition-rig-etale-continuous-homomorphism}
の意味で rig-\'etale であることを示せば十分である。
$I$ を (3) のものとする。$(B, I)$ 上の $A$ の素朴な余接複体
$\NL_{A/B}^\wedge$ は，$J/J^2$ を次数 $-1$ に置いて得られる $A$-加群の複体である。
従って Definition \ref{restricted-definition-rig-etale-homomorphism} により，
$A$ は $(B, I)$ 上 rig-\'etale である。

\medskip\noindent
(1) を仮定し，$V$ と $B \to A$ を (3) のものとする。
Definition \ref{restricted-definition-rig-smooth} により，$B \to A$ は rig-smooth である。
任意の定義イデアル $I \subset B$ を選ぶ。このとき $A$ は
$(B, I)$ 上 rig-smooth である。上と同様，複体 $\NL_{A/B}^\wedge$ は
$J/J^2$ を次数 $-1$ に置いて得られる。従って Lemma
\ref{restricted-lemma-equivalent-with-artin-smooth} により，ある $n \geq 1$ に対して
$J/J^2$ は冪 $I^n$ により零化される。$B$ は adic なので，
$I^n$ は $B$ の定義イデアルであり，証明は完了する。
\end{proof}

\section{Rig-全射な射}
\label{restricted-section-rig-surjective}

\noindent
局所 Noether 形式代数空間の間の局所有限型射に対しては，
\cite{ArtinII} から借用した定義を用いることができる。
より一般の場合にどうすべきかについては Remark
\ref{restricted-remark-rig-surjective-more-general} を参照せよ。

\begin{definition}
\label{restricted-definition-rig-surjective}
$S$ をスキームとする。$f : X \to Y$ を $S$ 上の形式代数空間の射とする。
$X$ と $Y$ は局所 Noether で，$f$ は局所有限型であると仮定する。
実線部分が与えられた任意の図式
$$
\xymatrix{
\text{Spf}(R') \ar@{..>}[r] \ar@{..>}[d] & X \ar[d]^f \\
\text{Spf}(R) \ar[r]^-p & Y
}
$$
で，$R$ が完備離散付値環，$p$ が adic 射であるものに対し，
完備離散付値環の拡大 $R \subset R'$ と，表示された図式を可換にする射
$\text{Spf}(R') \to X$ が存在するとき，$f$ は {\it rig-全射} であるという。
\end{definition}

\noindent
以下の補題で，この概念が局所 Noether 形式代数空間と局所有限型射の文脈において
適切に振る舞うことを見る。次の注意では，この定義をより広い形式代数空間の射の
クラスへ修正する選択肢を論じる。

\begin{remark}
\label{restricted-remark-rig-surjective-more-general}
Definition \ref{restricted-definition-rig-surjective} で定式化した条件は，
局所 adic* 形式代数空間の有限型射に対してさえ適切ではない。
例えば，$A = (\bigcup_{n \geq 1} k[t^{1/n}])^\wedge$ とし，
完備化を $t$-進完備化とする。このとき，$R$ が完備離散付値環であるような
adic 射 $\text{Spf}(R) \to \text{Spf}(A)$ は存在しない。
従って任意の射 $X \to \text{Spf}(A)$ が rig-全射となる。しかし $A$ は整域で，
$t \in A$ は零でないので，$A$ は少なくとも一つの「rig-点」をもつと考えたいし，
$X = \emptyset$ を許したくはない。この特別な場合を扱うには，adic 射
$$
\text{Spf}(R) \longrightarrow Y
$$
で，$R$ が主イデアル $J$ に関して完備な付値環であり，その根基が
$\mathfrak m_R = \sqrt{J}$ であるものを考えればよい。この場合，$R$ の値群は
$(\mathbf{R}, +)$ に埋め込め，$Y$ に付随する解析空間を定義する際に Berkovich が
用いた観点が得られる。\cite{Berkovich} を参照せよ。別の方法は Huber が推進した
ものである。彼の理論では，$\Spec(R/J)$ が一点集合であるという仮定を外す。
\cite{Huber-continuous-valuations} を参照せよ。
\end{remark}

\begin{lemma}
\label{restricted-lemma-composition-rig-surjective}
\begin{slogan}
局所有限型射の rig-全射性は合成について保存される
\end{slogan}
$S$ をスキームとする。$f : X \to Y$ と $g : Y \to Z$ を，
$S$ 上の形式代数空間の射とする。$X$，$Y$，$Z$ は局所 Noether で，
% P10-E0701
$f$ と $g$ は局所有限型であると仮定する。$f$ と $g$ が rig-全射ならば，
$g \circ f$ も rig-全射である。
\end{lemma}

\begin{proof}
定義から直接従う（さらに Formal Spaces, Lemma
\ref{formal-spaces-lemma-composition-finite-type} を用いる）。
\end{proof}

\begin{lemma}
\label{restricted-lemma-base-change-rig-surjective}
% P10-E0700
$S$ をスキームとする。$f : X \to Y$ と $Z \to Y$ を，$S$ 上の
形式代数空間の射とする。$X$，$Y$，$Z$ は局所 Noether で，
% P10-E0701
$f$ と $g$ は局所有限型であると仮定する。$f$ が rig-全射ならば，
基底変換 $Z \times_Y X \to Z$ も rig-全射である。
\end{lemma}

\begin{proof}
定義から直接従う（さらに Formal Spaces, Lemmas
\ref{formal-spaces-lemma-fibre-product-Noetherian} と
\ref{formal-spaces-lemma-base-change-finite-type} を用いる）。
\end{proof}

\begin{lemma}
\label{restricted-lemma-rig-surjective-alternative-permanence}
$S$ をスキームとする。$f : X \to Y$ と $g : Y \to Z$ を，
$S$ 上の局所 Noether 形式代数空間の局所有限型射とする。
$g \circ f$ が rig-全射で $g$ がモノ射ならば，$f$ は rig-全射である。
\end{lemma}

\begin{proof}
Lemma \ref{restricted-lemma-base-change-rig-surjective} と，$f$ が $g$ による
$g \circ f$ の基底変換であることを用いよ。
\end{proof}

\begin{lemma}
\label{restricted-lemma-permanence-rig-surjective}
$S$ をスキームとする。$f : X \to Y$ と $g : Y \to Z$ を，$S$ 上の
形式代数空間の射とする。$X$，$Y$，$Z$ は局所 Noether で，
% P10-E0701
$f$ と $g$ は局所有限型であると仮定する。$g \circ f : X \to Z$ が
rig-全射ならば，$g : Y \to Z$ も rig-全射である。
\end{lemma}

\begin{proof}
定義から直ちに従う。
\end{proof}

\begin{lemma}
\label{restricted-lemma-etale-covering-rig-surjective}
% P10-E0712
$S$ をスキームとする。$f : X \to Y$ を，代数空間により表現可能，
\'etale，かつ全射である局所 Noether 形式代数空間の射とする。
このとき $f$ は rig-全射である。
\end{lemma}

\begin{proof}
$p : \text{Spf}(R) \to Y$ を adic 射とし，$R$ を完備離散付値環とする。
$Z = \text{Spf}(R) \times_Y X$ と置く。このとき
$Z \to \text{Spf}(R)$ は代数空間により表現可能，\'etale，かつ全射である。
従って $Z$ は空でない。空でないアフィン形式代数空間 $V$ と \'etale 射
$V \to Z$ を選ぶ（これは我々の定義により可能である）。このとき
$V \to \text{Spf}(R)$ は，$R \to A$ が \'etale 環写像であるような
$R \to A^\wedge$ に対応する。Formal Spaces, Lemma
\ref{formal-spaces-lemma-etale} を参照せよ。$A^\wedge \not = 0$
（$V \not = \emptyset$ だから）なので，$A$ の極大イデアル $\mathfrak m$ で
$\mathfrak m_R$ の上にあるものを見つけられる。このとき
$A_\mathfrak m$ は離散付値環である（More on Algebra, Lemma
\ref{more-algebra-lemma-Dedekind-etale-extension}）。従って
$R' = A_\mathfrak m^\wedge$ は完備離散付値環である（More on Algebra, Lemma
\ref{more-algebra-lemma-completion-dvr}）。Formal Spaces, Lemma
% P10-E0702
\ref{formal-spaces-lemma-morphism-between-formal-spectra} を適用すると，
所望の射 $\text{Spf}(R') \to V \to Z \to X$ が得られる。
\end{proof}

\noindent
上の補題の帰結として，$f : X \to Y$ が rig-全射であるかどうかを
$Y$ 上 \'etale 局所的に検査できる。

\begin{lemma}
\label{restricted-lemma-upshot}
% P10-E0712
$S$ をスキームとする。$f : X \to Y$ を，局所有限型な局所 Noether
形式代数空間の射とする。$\{g_i : Y_i \to Y\}$ を，代数空間により表現可能かつ
\'etale で，$\coprod g_i$ が全射であるような形式代数空間の射の族とする。
このとき，$f$ が rig-全射であるための必要十分条件は，各
$f_i : X \times_Y Y_i \to Y_i$ が rig-全射であることである。
\end{lemma}

\begin{proof}
実際，$f$ が rig-全射ならば，任意の基底変換も rig-全射である
（Lemma \ref{restricted-lemma-base-change-rig-surjective}）。逆に，すべての $f_i$ が
rig-全射ならば，$\coprod f_i : \coprod X \times_Y Y_i \to \coprod Y_i$
も rig-全射である。Lemma \ref{restricted-lemma-etale-covering-rig-surjective} により，射
$\coprod g_i : \coprod Y_i \to Y$ は rig-全射である。従って
$\coprod X \times_Y Y_i \to Y$ は rig-全射である
（Lemma \ref{restricted-lemma-composition-rig-surjective}）。この射は $X \to Y$ を経由するので，
Lemma \ref{restricted-lemma-permanence-rig-surjective} により $X \to Y$ は rig-全射である。
\end{proof}

\begin{lemma}
\label{restricted-lemma-faithfully-flat-rig-surjective}
$A$ を，イデアル $I$ に関して完備な Noether 環とする。
$B$ を $I$-進完備な $A$-代数とする。すべての $n$ に対して
$A/I^n \to B/I^nB$ が有限型かつ平坦であり，$n = 1$ に対して忠実平坦ならば，
$\text{Spf}(B) \to \text{Spf}(A)$ は rig-全射である。
\end{lemma}

\begin{proof}
形式スペクトル間の射が，対応する位相環間の連続写像によって与えられることを
以下では断りなく用いる。Formal Spaces, Lemma
\ref{formal-spaces-lemma-morphism-between-formal-spectra} を参照せよ。
% P10-E0703
$\varphi : A \to R$ を，完備離散付値環 $A$ への連続写像とする。
これは $\varphi(I) \subset \mathfrak m_R$ を含意する。他方，$\varphi$ が adic 射に
対応する場合にのみ持ち上げ
% P10-E0704
$\varphi' : B' \to R'$ を構成すればよいので，$\varphi(I) \not = 0$ と仮定してよい。
従って，例えば Remark \ref{restricted-remark-base-change} を用いて基底変換
$C = B \widehat{\otimes}_A R$ を考えられる。このとき $C$ は
$\mathfrak m_R$-進完備な $R$-代数であり，$C/\mathfrak m_R^n C$ は
$R/\mathfrak m_R^n$ 上有限型かつ平坦で，$C/\mathfrak m_R C$ は零でない。
$\mathfrak m_R$ の上にある任意の極大イデアル $\mathfrak m \subset C$ を選ぶ。
平坦性（これは going down を含意する）により，
$\Spec(C_\mathfrak m) \setminus V(\mathfrak m_R C_\mathfrak m)$
は空でない開集合である。従って，
% P10-E0705
$\mathfrak q$ が $\Spec(C_\mathfrak m) \setminus \{\mathfrak m\}$ の閉点を定め，かつ
$\mathfrak q \not \in V(IC_\mathfrak m)$ となる素イデアル
$\mathfrak q \subset \mathfrak m$ を選べる。Properties, Lemma
\ref{properties-lemma-complement-closed-point-Jacobson} を参照せよ。
このとき $C/\mathfrak q$ は $1$ 次元局所整域であり，離散付値環 $R'$ をもつ包含
$C/\mathfrak q \subset R'$ を見つけられる
（Algebra, Lemma \ref{algebra-lemma-exists-dvr}）。構成により
$\mathfrak m_R R' \subset \mathfrak m_{R'}$ であり，$C \to R'$ は連続写像
$C \to (R')^\wedge$ に拡張される（実際，現在の状況では
$R' = (R')^\wedge$ となるよう $R'$ を選べるが，これは必要ない）。
離散付値環の完備化は離散付値環なので，仮定から環の可換図式
$$
\xymatrix{
(R')^\wedge & C \ar[l] & B \ar[l] \\
R \ar[u] & R \ar[l] \ar[u] & A \ar[l] \ar[u]
}
$$
が得られ，これが所望の持ち上げを与える。
\end{proof}

\begin{lemma}
\label{restricted-lemma-flat-rig-surjective}
$A$ を，イデアル $I$ に関して完備な Noether 環とする。
$B$ を $I$-進完備な $A$-代数とする。次を仮定する。
\begin{enumerate}
\item $A$ の $I$-torsion は $0$ である，
\item すべての $n$ に対して $A/I^n \to B/I^nB$ は平坦かつ有限型である。
\end{enumerate}
このとき，$\text{Spf}(B) \to \text{Spf}(A)$ が rig-全射であるための必要十分条件は，
$A/I \to B/IB$ が忠実平坦であることである。
\end{lemma}

\begin{proof}
忠実平坦性は Lemma \ref{restricted-lemma-faithfully-flat-rig-surjective} により
rig-全射性を含意する。逆を示すため，$I$-torsion の消滅と $I$-冪 torsion の
消滅が同値であることを以下では断りなく用いる（More on Algebra, Lemma
\ref{more-algebra-lemma-torsion-free}）。また，形式スペクトル間の射が対応する
位相環間の連続写像によって与えられることも断りなく用いる。Formal Spaces, Lemma
\ref{formal-spaces-lemma-morphism-between-formal-spectra} を参照せよ。

\medskip\noindent
$\text{Spf}(B) \to \text{Spf}(A)$ が rig-全射であると仮定する。
極大イデアル $I \subset \mathfrak m \subset A$ を選ぶ。
$\Spec(A_\mathfrak m)$ の開集合
$U = \Spec(A_\mathfrak m) \setminus V(I_\mathfrak m)$ は空でない。実際，
$A_\mathfrak m$ の $I_\mathfrak m$-torsion は零である
（Algebra, Lemma \ref{algebra-lemma-Noetherian-power-ideal-kills-module} を用いよ）。
従って，$U$ の点を定め（すなわち $IA_\mathfrak m \not \subset \mathfrak q$），
$\Spec(A_\mathfrak m) \setminus \{\mathfrak m\}$ の閉点に対応する素イデアル
$\mathfrak q \subset A_\mathfrak m$ を見つけられる。Properties, Lemma
\ref{properties-lemma-complement-closed-point-Jacobson} を参照せよ。
このとき $A_\mathfrak m/\mathfrak q$ は $1$ 次元局所整域である。従って，
$R$ が離散付値環であるような局所環の単射局所準同型
$A_\mathfrak m/\mathfrak q \subset R$ を見つけられる
（Algebra, Lemma \ref{algebra-lemma-exists-dvr}）。構成により
$IR \subset \mathfrak m_R$ であり，$A \to R$ は連続写像
$A \to R^\wedge$ に拡張される。離散付値環の完備化は離散付値環なので，
仮定から環の可換図式
$$
\xymatrix{
R' & B \ar[l] \\
R^\wedge \ar[u] & A \ar[l] \ar[u]
}
$$
が得られる。従って $\mathfrak m$ の上にある $B$ の素イデアルが見つかる。
よって $\Spec(B/IB) \to \Spec(A/I)$ は全射であり，従って
$A/I \to B/IB$ は忠実平坦である
（Algebra, Lemma \ref{algebra-lemma-ff-rings}）。
\end{proof}

\begin{lemma}
\label{restricted-lemma-monomorphism-rig-surjective}
% P10-E0712
$S$ をスキームとする。$f : X \to Y$ を形式代数空間の射とする。
$X$ と $Y$ は局所 Noether，$f$ は局所有限型，かつ $f$ はモノ射であると仮定する。
% P10-E0706
このとき $f$ が rig 全射であるための必要十分条件は，$R$ が完備離散付値環である
任意の adic 射 $\text{Spf}(R) \to Y$ が $X$ を経由することである。
\end{lemma}

\begin{proof}
一方の向きは自明である。他方について，$\text{Spf}(R) \to Y$ を adic 射とし，
完備離散付値環の拡大 $R \subset R'$ が存在して，
% P10-E0707
$\text{Spf}(R') \to \text{Spf}(R) \to X$ が $Y$ を経由すると仮定する。
$\Spec(R'/\mathfrak m_R^n R') \to \Spec(R/\mathfrak m_R^n)$ は全射かつ平坦である。
従って，$X$ は fpqc 被覆に対する層条件を満たすので，射
% P10-E0708
$\Spec(R/\mathfrak m_R^n) \to X$ は $X$ を経由する。Formal Spaces, Lemma
\ref{formal-spaces-lemma-sheaf-fpqc} を参照せよ。言い換えれば，
$\text{Spf}(R) \to Y$ は $X$ を経由する。
\end{proof}

\begin{lemma}
\label{restricted-lemma-closed-immersion-rig-surjective}
% P10-E0712
$S$ をスキームとする。$f : X \to Y$ を形式代数空間の射とする。
$X$ と $Y$ は局所 Noether で，$f$ は閉埋込みであると仮定する。
次の条件は同値である。
\begin{enumerate}
\item $f$ は rig-全射である，および
\item 任意のアフィン形式代数空間 $V$ と，代数空間により表現可能かつ
\'etale である任意の射 $V \to Y$ に対して，射 $X \times_Y V \to V$ は
$\textit{WAdm}^{Noeth}$ の全射 $B \to A$ で，その核 $J$ が次の性質をもつものに
対応する：$B$ のある定義イデアル $I$ とある $n \geq 1$ に対して
$IJ^n = 0$ である。
\end{enumerate}
\end{lemma}

\begin{proof}
(2) におけるような $V$ と $V \to Y$ が与えられたとき，$f$ にこれ以上の仮定を
課さなくても，Formal Spaces, Lemma
\ref{formal-spaces-lemma-closed-immersion-into-countably-indexed}
により，射 $X \times_Y V \to V$ は $\textit{WAdm}^{Noeth}$ の全射
$B \to A$ に対応することをまず観察しよう。

\medskip\noindent
(1) を仮定する。Lemma \ref{restricted-lemma-base-change-rig-surjective} により，
$\text{Spf}(A) \to \text{Spf}(B)$ は rig-全射である。
$I \subset B$ を定義イデアルとする。$B$ は adic なので，すべての
$m \geq 1$ に対して $I^m \subset B$ は定義イデアルである。
すべての $n, m \geq 1$ に対して $I^m J^n \not = 0$ ならば，
$IJ$ は冪零でなく，従って $V(IJ) \not = \Spec(B)$ である。
よって $\mathfrak p \not \in V(I) \cup V(J)$ となる素イデアル
$\mathfrak p \subset B$ を見つけられる。
$I(B/\mathfrak p) \not = B/\mathfrak p$ であることに注意すると，
極大イデアル $\mathfrak p + I \subset \mathfrak m \subset B$ を見つけられる。
Algebra, Lemma \ref{algebra-lemma-exists-dvr} により，離散付値環 $R$ と
単射局所環準同型 $(B/\mathfrak p)_\mathfrak m \to R$ を見つけられる。
明らかに，環写像 $B \to R$ は $A = B/J$ を経由できない。Lemma
\ref{restricted-lemma-monomorphism-rig-surjective} によれば，これは
$\text{Spf}(A) \to \text{Spf}(B)$ が rig-全射であることに矛盾する。
従って，ある $n, m$ に対して $I^n J^m = 0$ であり，(2) が成り立つ。

\medskip\noindent
(2) を仮定する。Lemma \ref{restricted-lemma-upshot} により，
$\text{Spf}(A) \to \text{Spf}(B)$ が rig-全射であることを示せば十分である。
定義イデアル $I \subset B$ と，$I J^n = 0$ となる整数 $n$ を選ぶ。
$R$ を離散付値環とし，$I$ の像が零でない環写像 $B \to R$ を考える。
$R$ は整域なので，$R$ における $J$ の像は零である。従って $B \to R$ は
全射 $B \to A$ を経由し，rig-全射な射の定義により証明が終わる。
\end{proof}

\begin{lemma}
\label{restricted-lemma-closed-immersion-rig-smooth-rig-surjective}
% P10-E0712
$S$ をスキームとする。$f : X \to Y$ を形式代数空間の射とする。
$X$ と $Y$ は局所 Noether で，$f$ は閉埋込みであると仮定する。
次の条件は同値である。
\begin{enumerate}
\item $f$ は rig-smooth かつ rig-全射である，
\item $f$ は rig-\'etale かつ rig-全射である，および
\item 任意のアフィン形式代数空間 $V$ と，代数空間により表現可能かつ
\'etale である任意の射 $V \to Y$ に対して，射 $X \times_Y V \to V$ は
$\textit{WAdm}^{Noeth}$ の全射 $B \to A$ で，その核 $J$ が次の性質をもつものに
対応する：$B$ のある定義イデアル $I$ に対して $IJ = 0$ である。
\end{enumerate}
\end{lemma}

\begin{proof}
$I$ と $J$ を，$IJ^n = 0$ かつ $I(J/J^2) = 0$ を満たす環 $B$ の
イデアルとする。このとき $I^nJ = 0$ である（証明は省略する）。
従って本補題は，Lemmas \ref{restricted-lemma-closed-immersion-rig-smooth} と
\ref{restricted-lemma-closed-immersion-rig-surjective} の自明な組合せから従う。
\end{proof}

\begin{lemma}
\label{restricted-lemma-rig-etale-descent}
$S$ をスキームとする。$f : X \to Y$ と $g : Y \to Z$ を，
$S$ 上の局所 Noether 形式代数空間の射とする。次を仮定する。
\begin{enumerate}
\item $g$ は局所有限型である，
\item $f$ は rig-smooth（resp.\ rig-\'etale）かつ rig-全射である，
\item $g \circ f$ は rig-smooth（resp.\ rig-\'etale）である。
\end{enumerate}
% P10-E0713
このとき $g$ は rig-smooth（resp.\ rig-\'etale）である。
\end{lemma}

\begin{proof}
rig-smooth の場合を証明し，rig-\'etale の場合に必要な変更を証明の最後に示す。
可換図式
$$
\xymatrix{
X \times_Y V \ar[r] \ar[d] &
V \ar[d] \ar[r] &
W \ar[d] \\
X \ar[r] &
Y \ar[r] &
Z
}
$$
を考える。ここで $V$ と $W$ はアフィン形式代数空間，$V \to Y$ と $W \to Z$ は
代数空間により表現可能かつ \'etale である。Definition
\ref{restricted-definition-rig-smooth} を参照すると，$V \to W$ が adic Noether 位相環の
rig-smooth 写像に対応することを示さねばならない。
% P10-E0709
$V = \text{Spf}(B)$，$W = \text{Spf}(C)$ と書け，さらに $V \to W$ は
位相的有限型な adic 環写像 $C \to B$ に対応する。Lemma
\ref{restricted-lemma-finite-type-morphisms} を参照せよ。

\medskip\noindent
以下では，$X \times_Y V \to V$ が rig-smooth かつ rig-全射であることを
断りなく用いる。Lemmas \ref{restricted-lemma-base-change-rig-smooth} と
\ref{restricted-lemma-base-change-rig-surjective} を参照せよ。また，$g \circ f$ が
rig-smooth なので，合成 $X \times_Y V \to V \to W$ も rig-smooth である。

\medskip\noindent
$I \subset C$ を定義イデアルとする。加群
% P10-E0710
$C \to B$ が rig-smooth でないと仮定して矛盾を導く。これは，$IB$ を含まない
素イデアル $\mathfrak q \subset B$ が存在して，
% P10-E0711
$H^{-1}(\NL_{B/C}^\wedge)_\mathfrak p$ が零でないか，または
$H^0(\NL_{B/C}^\wedge)_\mathfrak p$ が有限自由
$B_\mathfrak q$-加群でないことを意味する。Lemma
\ref{restricted-lemma-equivalent-with-artin-smooth} を参照せよ。いくつかの詳細は省略する。
極大イデアル $IB + \mathfrak q \subset \mathfrak m$ を選べる。Algebra, Lemma
\ref{algebra-lemma-exists-dvr} により，完備離散付値環 $R$ と単射局所環準同型
$(B/\mathfrak q)_\mathfrak m \to R$ を見つけられる。

\medskip\noindent
$R$ を拡大で置き換えた後，adic 射
$\text{Spf}(R) \to V = \text{Spf}(B)$ の持ち上げ
$\text{Spf}(R) \to X \times_Y V$ が与えられていると仮定してよい。
Formal Spaces, Definition
\ref{formal-spaces-definition-formal-algebraic-space}
の意味での \'etale 被覆
$\{\text{Spf}(A_i) \to X \times_Y V\}$ を選ぶ。Lemma
\ref{restricted-lemma-etale-covering-rig-surjective} により，ある $i$ に対して
$\text{Spf}(R) \to X \times_Y V$ が射
$\text{Spf}(R) \to \text{Spf}(A_i)$ に持ち上がると仮定してよい
（このためには $R$ を別の拡大で置き換える必要があるかもしれない）。
$A = A_i$ と置く。環写像
$$
C \to B \to A \to R
$$
を考える。$\mathfrak p \subset A$ を写像 $A \to R$ の核とし，
$\mathfrak p$ は $\mathfrak q$ の上にあることに注意せよ。
$C \to A$ と $B \to A$ が rig-smooth であることは分かっている。
特に，Lemma \ref{restricted-lemma-rig-smooth-rig-flat} により環写像
$B_\mathfrak q \to A_\mathfrak p$ は平坦である。Lemmas
\ref{restricted-lemma-exact-sequence-NL} と
\ref{restricted-lemma-exact-sequence-NL-rig-smooth} の付随する完全列
$$
\xymatrix{
&
H^0(\NL_{B/C}^\wedge) \otimes_B A_\mathfrak p \ar[r] &
H^0(\NL_{A/C}^\wedge)_\mathfrak p \ar[r] &
H^0(\NL_{A/B}^\wedge)_\mathfrak p \ar[r] & 0 \\
0 \ar[r] &
H^{-1}(\NL_{B/C}^\wedge \otimes_B A)_\mathfrak p \ar[r] &
H^{-1}(\NL_{A/C}^\wedge)_\mathfrak p \ar[r] &
H^{-1}(\NL_{A/B}^\wedge)_\mathfrak p \ar[llu]
}
$$
を考える。$C \to A$ と $B \to A$ の rig-smooth 性から，
$H^{-1}(\NL_{B/C}^\wedge \otimes_B A)_\mathfrak p = 0$ であり，
$H^0(\NL_{B/C}^\wedge) \otimes_B A_\mathfrak p$ は有限自由
$A_\mathfrak p$-加群の全射の核として有限自由であると分かる。
$B_\mathfrak q \to A_\mathfrak p$ は平坦であり，従って忠実平坦なので，これは
$H^{-1}(\NL_{B/C}^\wedge)_\mathfrak q = 0$ であり，
$H^0(\NL_{B/C}^\wedge)_\mathfrak q$ が有限自由であることを含意する。
これは求めていた矛盾である。

\medskip\noindent
rig-\'etale の場合もまったく同様に論じるが，得られる結論は
$H^{-1}(\NL_{B/C}^\wedge)_\mathfrak q$ と
$H^0(\NL_{B/C}^\wedge)_\mathfrak q$ がともに零であることである。
\end{proof}

\section{完備離散付値環上の形式代数空間}
\label{restricted-section-over-cdrv}

\noindent
本節では次の用語法を用いる。$A$ が弱許容位相環であるとき，
「{it $X$ は $A$ 上の形式代数空間である}」とは，$X$ が形式代数空間であって，
形式代数空間の射 $p : X \to \text{Spf}(A)$ を備えていることをいう。
この状況では $p$ を {\it 構造射}と呼ぶ。

\begin{lemma}
\label{restricted-lemma-flat-locus}
$X$ を完備離散付値環 $A$ 上の局所 Noether 形式代数空間とする。
このとき，形式代数空間の閉埋込み $X' \to X$ で，$X'$ は $A$ 上平坦であり，
局所 Noether 形式代数空間の任意の射 $Y \to X$ で $Y$ が $A$ 上平坦なものが
$X'$ を経由するようなものが存在する。
\end{lemma}

\begin{proof}
$\pi \in A$ を一様化元とする。$A$-加群が平坦であるための必要十分条件は，
$\pi$-冪ねじれが $0$ であることを思い出そう。

\medskip\noindent
まず $X$ がアフィン形式代数空間であると仮定する。このとき
$X = \text{Spf}(B)$ であり，$B$ は adic Noether $A$-代数である。
この場合，$B' = B/\pi\text{-power torsion}$ として $X' = \text{Spf}(B')$ と置く。
$X'$ が $A$ 上平坦で，$X' \to X$ が閉埋込みであることは明らかである。
$g : Y \to X$ を局所 Noether 形式代数空間の射とし，$Y$ は $A$ 上平坦とする。
Formal Spaces, Definition
\ref{formal-spaces-definition-formal-algebraic-space}
の意味での被覆 $\{Y_j \to Y\}$ を選ぶ。このとき $Y_j = \text{Spf}(C_j)$ で，
$C_j$ は $A$ 上平坦である。従って，
% P10-E0714
$Y_j \to X$ は連続な
% P10-E0715
$R$-代数写像 $B \to C_j$ に対応し，$B \to C_j$ は明らかに
$\pi$-冪ねじれを零化するので $X'$ を経由する。
$\{Y_j \to Y\}$ は被覆で，$X' \to X$ はモノ射なので，
$g$ は $X'$ を経由すると結論できる。

\medskip\noindent
$X$ と $\{X_i \to X\}_{i \in I}$ を Formal Spaces, Definition
\ref{formal-spaces-definition-formal-algebraic-space}
におけるものとする。各 $i$ に対し，上で構成した平坦部分を
$X'_i \to X_i$ とする。$i, j \in I$ に対し，射影
$X'_i \times_X X_j \to X'_i$ は（仮定により）\'etale であり，
（Formal Spaces, Lemma
\ref{formal-spaces-lemma-presentation-representable} により）スキームの射である。
従って，代数空間により表現可能な \'etale 射は平坦なので（Lemma
\ref{restricted-lemma-representable-flat}），$X'_i \times_X X_j$ は $A$ 上平坦である。
よって射影 $X'_i \times_X X_j \to X_j$ は普遍性により $X'_j$ を経由する。
射 $X'_i \to X_i$ は層の単射なので，
$$
R_{ij} = X'_i \times_X X_j = X'_i \times_X X'_j = X_i \times_X X'_j
$$
と結論する。$U = \coprod X'_i$，$R = \coprod R_{ij}$ と置き，
$s, t : R \to U$ を二つの射影とする。層として $R = U \times_X U$ であり，
$s$ と $t$ は \'etale である。以上の観察により，
% P10-E0716
$(t, s) : R \to U$ は \'etale 同値関係を定める。従って
Spaces, Theorem \ref{spaces-theorem-presentation} により
$X' = U/R$ は代数空間である。構成により図式
$$
\xymatrix{
\coprod X'_i \ar[r] \ar[d] & \coprod X_i \ar[d] \\
X' \ar[r] & X
}
$$
は Cartesian である。右の垂直射は \'etale 全射で，上の水平射は表現可能な
閉埋込みなので，Bootstrap, Lemma
\ref{bootstrap-lemma-after-fppf-sep-lqf} により $X' \to X$ は表現可能である。
次に Spaces, Lemma
\ref{spaces-lemma-descent-representable-transformations-property}
を用いて，$X' \to X$ は閉埋込みであると結論できる。

\medskip\noindent
最後に，$Y \to X$ を射とし，$Y$ は $A$ 上平坦な局所 Noether
形式代数空間であると仮定する。このとき各 $X_i \times_X Y$ は $Y$ 上 \'etale，
従って $A$ 上平坦である（上を参照せよ）。よって
$X_i \times_X Y \to X_i$ は $X'_i$ を経由する。
$\{X_i \times_X Y \to Y\}$ は \'etale 被覆なので，$Y \to X$ は $X'$ を経由する。
\end{proof}

\begin{lemma}
\label{restricted-lemma-flat-and-diagonal-rig-surjective}
$X$ を，完備離散付値環 $A$ 上局所有限型な局所 Noether 形式代数空間とする。
$X' \subset X$ を Lemma \ref{restricted-lemma-flat-locus} におけるものとする。
$X \to X \times_{\text{Spf}(A)} X$ が rig-\'etale かつ rig-全射ならば，
$X' = \text{Spf}(A)$ または $X' = \emptyset$ である。
\end{lemma}

\begin{proof}
（補足：対角射は Formal Spaces, Lemma
\ref{formal-spaces-lemma-diagonal-morphism-formal-algebraic-spaces}
により常に局所有限型であり，$X \times_{\text{Spf}(A)} X$ は Formal Spaces, Lemmas
\ref{formal-spaces-lemma-base-change-finite-type} と
\ref{formal-spaces-lemma-locally-finite-type-locally-noetherian}
により局所 Noether である。従って，対角射にこれらの条件を課すことには意味がある。）
図式
$$
\xymatrix{
X' \ar[r] \ar[d] & X' \times_{\text{Spf}(A)} X' \ar[d] \\
X \ar[r] & X \times_{\text{Spf}(A)} X
}
$$
は Cartesian である。従って $X' \to X' \times_{\text{Spf}(A)} X'$ は
% P10-E0717
Lemma \ref{restricted-lemma-base-change-rig-surjective} により rig-\'etale かつ rig-全射である。
アフィン形式代数空間 $U$ と，代数空間により表現可能かつ \'etale である射
$U \to X'$ を選ぶ。このとき $U = \text{Spf}(B)$ である。ここで $B$ は
adic Noether 位相環であって，平坦な $A$-代数であり，その位相は
一様化元 $\pi \in A$ による $\pi$-進位相であり，さらに各 $n$ に対して
$A/\pi^n A \to B/\pi^n B$ は有限型である。後で用いるため，これは特に
次を含意することに注意する：$B \not = 0$ ならば，写像
$\text{Spf}(B) \to \text{Spf}(A)$ は層の全射である（いつものように
fppf 位相を用いていることを思い出されたい）。上の議論を繰り返すと，
$$
W = U \times_{X'} U =
X' \times_{X' \times_{\text{Spf}(A)} X'} (U \times_{\text{Spf}(A)} U)
\longrightarrow
U \times_{\text{Spf}(A)} U
$$
は閉埋込みであり，rig-\'etale かつ rig-全射である。Formal Spaces, Lemma
\ref{formal-spaces-lemma-fibre-product-affines-over-separated} により
$U \times_{\text{Spf}(A)} U = \text{Spf}(B \widehat{\otimes}_A B)$ である。
$B \widehat{\otimes}_A B$ は，平坦な $A$-代数 $B \otimes_A B$ の
$\pi$-進完備化として平坦な $A$-代数である。従って Lemma
\ref{restricted-lemma-closed-immersion-rig-smooth-rig-surjective} により
$W = U \times_{\text{Spf}(A)} U$ である。言い換えると，
$U \times_{X'} U = U \times_{\text{Spf}(A)} U$ であり，これはさらに
$U \to X'$ の像（層の写像として）が $\text{Spf}(A)$ へ単射的に写ることを意味する。
Formal Spaces, Definition
\ref{formal-spaces-definition-formal-algebraic-space}
の意味での被覆 $\{U_i \to X'\}$ を選ぶ。特に
$\coprod U_i \to X'$ は層の全射である。上の議論を
% P10-E0718
$U_i \coprod U_j \to X'$ に適用すると（$U_i \amalg U_j$ もアフィン形式代数空間
であることを用いる），$X' \to \text{Spf}(A)$ は fppf 層の単射であると分かる。
$X'$ は $A$ 上平坦なので，$X'$ は空である（すべての $i$ に対して $U_i$ が
空の場合）か，この写像は同型である（ある $i$ に対して $U_i$ が空でない場合；
このとき $U_i \to \text{Spf}(A)$ は層の全射であることを既に見た）。
これで証明は完了する。
\end{proof}

\begin{lemma}
\label{restricted-lemma-rig-monomorphism-rig-surjective}
% P10-E0721
$S$ をスキームとする。$f : X \to Y$ を形式代数空間の射とする。次を仮定する。
\begin{enumerate}
\item $X$ と $Y$ は局所 Noether である，
% P10-E0719
\item $f$ は局所有限型である，
\item $\Delta_f : X \to X \times_Y X$ は rig-\'etale かつ rig-全射である。
\end{enumerate}
% P10-E0720
このとき $f$ が rig 全射であるための必要十分条件は，$R$ が完備離散付値環である
任意の adic 射 $\text{Spf}(R) \to Y$ が射 $\text{Spf}(R) \to X$ に持ち上がることである。
\end{lemma}

\begin{proof}
一方の向きは自明である。他方について，$\text{Spf}(R) \to Y$ を adic 射とし，
完備離散付値環の拡大 $R \subset R'$ が存在して，
% P10-E0722
$\text{Spf}(R') \to \text{Spf}(R) \to X$ が $Y$ を経由すると仮定する。
ファイバー積図式
$$
\xymatrix{
\text{Spf}(R') \ar[r] \ar[rd] &
\text{Spf}(R) \times_Y X \ar[r] \ar[d]^p &
X \ar[d]^f \\
&
\text{Spf}(R) \ar[r] &
Y
}
$$
を考える。射 $p$ は $f$ の基底変換として局所有限型である。Formal Spaces, Lemma
\ref{formal-spaces-lemma-base-change-finite-type} を参照せよ。
対角射 $\Delta_p$ は $\Delta_f$ の基底変換なので，rig-\'etale かつ rig-全射である。
Lemma \ref{restricted-lemma-flat-and-diagonal-rig-surjective} により，
$\text{Spf}(R) \times_Y X$ の $R$ 上の平坦部分は，$\emptyset$ または
$\text{Spf}(R)$ に等しい。しかし $\text{Spf}(R')$ はそれを経由するので，
それは空でなく，従って所望の射
$\text{Spf}(R) \to \text{Spf}(R) \times_Y X \to X$ を得る。
\end{proof}

\section{完備化関手}
\label{restricted-section-completion-functor}

\noindent
この節では次の状況を考える。
まず基礎スキーム $S$ を固定する。すべての環，位相環，
スキーム，代数空間，形式代数空間，およびこれらの間の射は
$S$ 上にあるものとする。次に，代数空間 $X$ と閉部分集合
$T \subset |X|$ を固定する。
% P10-E0723
$|U| = |X| \setminus T$ を満たす開部分空間 $U \subset X$ を考える。
図式は
$$
U \to X \quad |X| = |U| \amalg T
$$
である。この状況で，$X$ 上の代数空間 $X'$，すなわち射
$f : X' \to X$ を備えた代数空間 $X'$ が与えられたとする。
% P10-E0724
$T$ の逆像を $T' \subset |X'|$ と書き，
$|U'| = |X'| \setminus T'$ を満たす開部分空間 $U' \subset X'$ を取る。
図式は
$$
U' = f^{-1}U
\quad\quad
\vcenter{
\xymatrix{
U' \ar[d] \ar[r] & X' \ar[d]_f \\
U \ar[r] & X
}
}
\quad\quad
\vcenter{
\xymatrix{
|U'| \ar[r] \ar[d] & |X'| \ar[d]^{|f|} & T' \ar[l] \ar[d] \\
|U| \ar[r] & |X| & T \ar[l]
}
}
\quad\quad
T' = |f|^{-1}T
$$
である。以下では $f$ の諸性質を，形式完備化の間に誘導される射
$$
f_{/T} : X'_{/T'} \longrightarrow X_{/T}
$$
の諸性質と関係づける。表示式が示すように，この射を
$f_{/T}$ と書く。Formal Spaces, Lemma
\ref{formal-spaces-lemma-map-completions-representable}
により，$f_{/T}$ が代数空間によって表現可能であることは既に見た。
実際，その補題の証明が示すように，図式
$$
\xymatrix{
X'_{/T'} \ar[d]_{f_{/T}} \ar[r] & X' \ar[d]^f \\
X_{/T} \ar[r] & X
}
$$
は Cartesian である。以下の補題の主張と証明を読む際には，
この事実を念頭に置かれたい。

\begin{lemma}
\label{restricted-lemma-map-completions-finite-type}
上の状況において，$f$ が局所有限型ならば，
$f_{/T}$ も局所有限型である。
\end{lemma}

\begin{proof}
（形式代数空間の有限型射については Formal Spaces, Section
\ref{formal-spaces-section-finite-type} で論じる。）
実際，$Z \to X$ を，$|Z|$ の像が $T$ に含まれるような
スキームから $X$ への射とする。上の Cartesian 正方形から，
$Z \times_X X'$ は $Z \times_{X_{/T}} X'_{/T'}$ を表現する
代数空間である。Morphisms of Spaces, Lemma
\ref{spaces-morphisms-lemma-base-change-finite-type} により
$Z \times_X X' \to Z$ は局所有限型なので，結論を得る。
\end{proof}

\begin{lemma}
\label{restricted-lemma-map-completions-etale}
上の状況において，$f$ が \'etale ならば，$f_{/T}$ も \'etale である。
\end{lemma}

\begin{proof}
Lemma \ref{restricted-lemma-map-completions-finite-type} の証明と同じ議論により，
これは Morphisms of Spaces, Lemma
\ref{spaces-morphisms-lemma-base-change-etale} から従う。
\end{proof}

\begin{lemma}
\label{restricted-lemma-closed-immersion-gives-closed-immersion}
上の状況において，$f$ が閉埋込みならば，$f_{/T}$ も閉埋込みである。
\end{lemma}

\begin{proof}
（形式代数空間の閉埋込みについては Formal Spaces, Section
\ref{formal-spaces-section-closed-immersions} で論じる。）
Lemma \ref{restricted-lemma-map-completions-finite-type} の証明と同じ議論により，
これは Spaces, Lemma
\ref{spaces-lemma-base-change-immersions} から従う。
\end{proof}

\begin{lemma}
\label{restricted-lemma-proper-gives-proper}
上の状況において，$f$ が proper ならば，$f_{/T}$ も proper である。
\end{lemma}

\begin{proof}
（形式代数空間の proper 射については Formal Spaces, Section
\ref{formal-spaces-section-proper} で論じる。）
Lemma \ref{restricted-lemma-map-completions-finite-type} の証明と同じ議論により，
これは Morphisms of Spaces, Lemma
\ref{spaces-morphisms-lemma-base-change-proper} から従う。
\end{proof}

\begin{lemma}
\label{restricted-lemma-quasi-compact-gives-quasi-compact}
上の状況において，$f$ が準コンパクトならば，$f_{/T}$ も準コンパクトである。
\end{lemma}

\begin{proof}
（形式代数空間の準コンパクト射については Formal Spaces, Section
\ref{formal-spaces-section-quasi-compact} で論じる。）
$(X'_{/T'})_{red} \to (X_{/T})_{red}$ が代数空間の準コンパクト射であることを
示せばよい。Formal Spaces, Lemma
\ref{formal-spaces-lemma-reduction-completion} により，これは射
$Z' \to Z$ である。ここで $Z' \subset X'$ および $Z \subset X$ は，
それぞれ $T'$ および $T$ 上に誘導される被約代数空間構造である。
従って $Z' \to f^{-1}Z = Z \times_X X'$ は厚化（すなわち基礎位相空間上の
同型を定める閉埋込み）である。$Z \times_X X' \to Z$ は $f$ の基底変換として
準コンパクトであるから（Morphisms of Spaces, Lemma
\ref{spaces-morphisms-lemma-base-change-quasi-compact}），
More on Morphisms of Spaces, Lemma
\ref{spaces-more-morphisms-lemma-thicken-property-morphisms}
により $Z' \to Z$ も準コンパクトである。
\end{proof}

\begin{remark}
\label{restricted-remark-diagonal-gives-diagonal}
上の状況において，対角射
$\Delta_f : X' \to X' \times_X X'$ および
$\Delta_{f_{/T}} : X'_{/T'} \to X'_{/T'} \times_{X_{/T}} X'_{/T'}$
を考える。$X' \times_X X'$ の部分関手として
$$
X'_{/T'} \times_{X_{/T}} X'_{/T'} = (X' \times_X X')_{/T''}
$$
が成り立つことは容易に分かる。ここで
$T'' \subset |X' \times_X X'|$ は $T$ の逆像である。
従って $\Delta_{f_{/T}} = (\Delta_f)_{/T''}$ である。
以下ではこれを用いて，$\Delta_f$ の諸性質が
$\Delta_{f_{/T}}$ に継承されることを示す。
\end{remark}

\begin{lemma}
\label{restricted-lemma-quasi-separated-gives-quasi-separated}
上の状況において，$f$ が（準）分離的ならば，$f_{/T}$ もそうである。
\end{lemma}

\begin{proof}
（形式代数空間の射の分離条件については Formal Spaces, Section
\ref{formal-spaces-section-separation-axioms} で論じる。）
$\Delta_f$ が準コンパクト，それぞれ閉埋込みならば，
$\Delta_{f_{/T}}$ についても同じことが成り立つと示せばよい。
これは Remark \ref{restricted-remark-diagonal-gives-diagonal} の議論と
Lemmas \ref{restricted-lemma-quasi-compact-gives-quasi-compact} および
\ref{restricted-lemma-closed-immersion-gives-closed-immersion} から従う。
\end{proof}

\begin{lemma}
\label{restricted-lemma-smooth-gives-rig-smooth}
上の状況において，$X$ が局所 Noether，$f$ が局所有限型，かつ
$U' \to U$ が smooth ならば，$f_{/T}$ は rig-smooth である。
\end{lemma}

\begin{proof}
証明の方針は次の通りである。$X$ と $X'$ がアフィンである場合に還元し，
アフィンの場合を代数に翻訳し，最後に Lemma \ref{restricted-lemma-rig-smooth}
を適用する。読者には詳細を飛ばすことを勧める。

\medskip\noindent
Properties of Spaces, Lemma
\ref{spaces-properties-lemma-cover-by-union-affines} により，
アフィンスキームの非交和 $W = \coprod W_i$ をもつ全射 \'etale 射
$W \to X$ を選ぶ。各 $i$ に対して，全射 \'etale 射
$W'_i \to W_i \times_X X'$ を選ぶ。ここで
$W'_i = \coprod W'_{ij}$ はアフィンの非交和である。
特に $\coprod W'_{ij} \to X'$ は全射かつ \'etale である。
% P10-E0725
与えられた射を $f_{ij} : W_{ij} \to W_i$ と書く。
$T$ の逆像を $T_i \subset W_i$ および $T'_{ij} \subset W_{ij}$ と書く。
$T$ の逆像に沿う完備化を取ると Cartesian 図式が得られるので（上を参照），
$(W_i)_{/T_i} = W_i \times_X X_{/T}$ であり，同様に
$(W'_{ij})_{/T'_{ij}} = W'_{ij} \times_{X'} X'_{/T'}$ である。
さらに，$(W_i)_{/T_i}$ と $(W'_{ij})_{/T'_{ij}}$ は
アフィン形式代数空間であることを思い出そう。
% P10-E0726
従って $\{W'_{ij})_{/T'_{ij}} \to X'_{/T'}\}$ は
Formal Spaces, Definition
\ref{formal-spaces-definition-formal-algebraic-space}
の意味での被覆である。Lemma \ref{restricted-lemma-rig-smooth-morphisms}
により，
$$
(W'_{ij})_{/T'_{ij}} \longrightarrow (W_i)_{/T_i}
$$
が rig-smooth であることを示せば十分である。
$W'_{ij} \to W_i$ は局所有限型であり，smooth 射
$W'_{ij} \setminus T'_{ij} \to W_i \setminus T_i$ を誘導することに注意する
（これは $f$ について成り立ち，射のこれらの性質は始域と終域について
\'etale 局所的だからである）。また $W_i$ は局所 Noether である
（$X$ が局所 Noether であり，この性質は代数空間上 \'etale 局所的だからである）。
従って $X$ と $X'$ がアフィンスキームである場合を証明すればよい。

\medskip\noindent
$X = \Spec(A)$ および $X' = \Spec(A')$ がアフィンスキームであると仮定する。
$X$ は Noether なので，$A$ は Noether である。射 $f$ は有限型環準同型
$A \to A'$ により与えられる。$T$ を切り出すイデアルを $I \subset A$
とする。このとき $IA'$ は $T'$ を切り出す。また
$\Spec(A') \to \Spec(A)$ は $\Spec(A) \setminus T$ 上 smooth である。
$A^\wedge$ と $(A')^\wedge$ を $I$-進完備化とする。Formal Spaces, Lemma
\ref{formal-spaces-lemma-formal-completion-types} の証明により，
$X_{/T} = \text{Spf}(A^\wedge)$ および
$X'_{/T'} = \text{Spf}((A')^\wedge)$ である。
% P10-E0727
Lemma \ref{restricted-lemma-rig-smooth} により，$(A')^\wedge$ は $(A. I)$ 上
rig-smooth である。これはさらに $A^\wedge \to (A')^\wedge$ が
rig-smooth であることを意味し，最終的に Lemma
% P10-E0728
\ref{restricted-lemma-rig-smooth-morphisms} により
$X'_{/T'} \to X_{/T}$ が rig smooth であることを含意する。
\end{proof}

\begin{lemma}
\label{restricted-lemma-etale-gives-rig-etale}
上の状況において，$X$ が局所 Noether，$f$ が局所有限型，かつ
$U' \to U$ が \'etale ならば，$f_{/T}$ は rig-\'etale である。
\end{lemma}

\begin{proof}
証明は Lemma \ref{restricted-lemma-smooth-gives-rig-smooth} の証明とまったく同じである。
ただし Lemmas \ref{restricted-lemma-rig-smooth} および
\ref{restricted-lemma-rig-smooth-morphisms} を，それぞれ Lemmas
\ref{restricted-lemma-rig-etale} および \ref{restricted-lemma-rig-etale-morphisms} に置き換える。
\end{proof}

\begin{lemma}
\label{restricted-lemma-completion-proper-surjective-rig-surjective}
上の状況において，$X$ が局所 Noether，$f$ が proper，かつ
$U' \to U$ が全射ならば，$f_{/T}$ は rig-全射である。
\end{lemma}

\begin{proof}
（この主張が意味をもつことは Lemma
\ref{restricted-lemma-map-completions-finite-type} および Formal Spaces, Lemma
\ref{formal-spaces-lemma-formal-completion-types} による。）
$R$ を分数体 $K$ をもつ完備離散付値環とする。
$p : \text{Spf}(R) \to X_{/T}$ を形式代数空間の adic 射とする。
Formal Spaces, Lemma
\ref{formal-spaces-lemma-adic-into-completion} により，合成
$\text{Spf}(R) \to X_{/T} \to X$ は射 $q : \Spec(R) \to X$
に対応し，この射は $\Spec(K)$ を $U$ に写す。
$U' \to U$ は proper かつ全射なので，
$\Spec(K) \times_U U'$ は空でなく $K$ 上 proper である。
従って体拡大 $K'/K$ と可換図式
$$
\xymatrix{
\Spec(K') \ar[r] \ar[d] & U' \ar[r] \ar[d] & X' \ar[d] \\
\Spec(K) \ar[r] & U \ar[r] & X
}
$$
を選べる。Algebra, Lemma \ref{algebra-lemma-exists-dvr} により，
$R$ を支配し分数体 $K'$ をもつ離散付値環 $R' \subset K'$ を取る。
$\Spec(K) \to X$ は $\Spec(R) \to X$ に延長されるので，proper 性の付値判定法
（Morphisms of Spaces, Lemma
\ref{spaces-morphisms-lemma-characterize-proper}）により，
$U'$ のこの $K'$-点を，$\Spec(R) \to X$ 上の射
$\Spec(R') \to X'$ に延長できる。従って $\Spec(R')$ における
$T'$ の逆像は閉点であり，所望の通り $p$ を持ち上げる adic 射
$\text{Spf}((R')^\wedge) \to X'_{/T'}$ を得る
（More on Algebra, Lemma
\ref{more-algebra-lemma-completion-dvr} により $(R')^\wedge$ は
完備離散付値環であることに注意せよ）。
\end{proof}

\begin{lemma}
\label{restricted-lemma-separated-mono-open-diagonal-rig-surjective}
上の状況において，$X$ が局所 Noether，$f$ が分離的かつ局所有限型，
さらに $U' \to U$ がモノ射ならば，$\Delta_{f_{/T}}$ は rig-全射である。
\end{lemma}

\begin{proof}
対角射 $\Delta_f : X' \to X' \times_X X'$ は閉埋込みであり，
$\Delta_f$ の制限 $U' \to U' \times_U U'$ は全射である。
従ってこの補題は Remark \ref{restricted-remark-diagonal-gives-diagonal} の議論と
Lemma \ref{restricted-lemma-completion-proper-surjective-rig-surjective} から従う。
\end{proof}

\section{形式的修正}
\label{restricted-section-formal-modifications}

\noindent
この節では，局所 Noether 形式代数空間の形式的修正という Artin の概念を
定義し，調べる。まず定義を与える。

\begin{definition}
\label{restricted-definition-formal-modification}
$S$ をスキームとする。$f : X \to Y$ を，$S$ 上の局所 Noether
形式代数空間の射とする。次を満たすとき，$f$ を
{\it 形式的修正}という。
\begin{enumerate}
\item $f$ は proper 射である（Formal Spaces, Definition
\ref{formal-spaces-definition-proper}），
\item $f$ は rig-\'etale である，
\item $f$ は rig-全射である，
\item $\Delta_f : X \to X \times_Y X$ は rig-全射である。
\end{enumerate}
\end{definition}

\noindent
典型例は Lemma \ref{restricted-lemma-modification-gives-formal-modification}
で与えられる。実際，後にすべての形式的修正がこの型で
「形式的に局所的」であることを示す。Lemma
\ref{restricted-lemma-formal-modifications-locally-algebraic} を参照せよ。
これらの条件を Artin の論文の条件と比較しよう。

\begin{remark}
\label{restricted-remark-compare-formal-modification-artin}
\cite[Definition 1.7]{ArtinII} では，形式的修正は，次の三条件を満たす
局所 Noether 形式代数空間の proper 射 $f : X \to Y$ として定義される
\footnote{これらの条件を Stacks project で展開した言語へ完全に翻訳することは
しない。それでも，ここの議論が読者の役に立つことを期待する。}。
\begin{enumerate}
% P10-E0729
\item[(\romannumeral1)] $f$ の Cramer イデアルと Jacobian イデアルは，
それぞれ $X$ の定義イデアルを含む，
\item[(\romannumeral2)] 対角写像
$\Delta : X \to X \times_Y X$ を定義するイデアルは，
$X \times_Y X$ の定義イデアルにより零化される，
\item[(\romannumeral3)] $R$ が完備離散付値環であるとき，
任意の adic 射 $\text{Spf}(R) \to Y$ は
$\text{Spf}(R) \to X$ に持ち上がる。
\end{enumerate}
これらを上の条件一覧と比較しよう。

\medskip\noindent
(\romannumeral1) について。性質 (\romannumeral1) は，$f$ が
rig-\'etale 射であるという我々の条件に一致する。これは Lemma
\ref{restricted-lemma-equivalent-with-artin} の
(\ref{restricted-item-condition-artin}) から従う。

\medskip\noindent
(\romannumeral2) について。$f$ が rig-\'etale であると仮定する。
このとき $\Delta_f : X \to X \times_Y X$ は，$X$ 上 rig-\'etale である
局所 Noether 形式代数空間の間の射として rig-\'etale である
（前者には $\text{id}_X$ を，後者には
$\text{pr}_1$ を用いる）。Lemmas \ref{restricted-lemma-base-change-rig-etale}
および \ref{restricted-lemma-rig-etale-permanence} を参照せよ。
従って Lemma \ref{restricted-lemma-closed-immersion-rig-smooth-rig-surjective}
により，性質 (\romannumeral2) は $\Delta_f$ が rig-全射であるという
我々の条件に一致する。

\medskip\noindent
(\romannumeral3) について。性質 (\romannumeral3) は，
我々の rig-全射の概念とは完全には一致しない。Artin はすべての adic 射
$\text{Spf}(R) \to Y$ が $X$ への射に持ち上がることを要求する一方，
我々の rig-全射の概念は，$R$ をある拡大で置き換えた後に持上げが存在すること
だけを主張するからである。しかし (\romannumeral1) と
(\romannumeral2) により $\Delta_f$ が既に rig-\'etale かつ
rig-全射であるから，Lemma \ref{restricted-lemma-rig-monomorphism-rig-surjective}
によりこれらの条件は同値である。
\end{remark}

\begin{lemma}
\label{restricted-lemma-modification-gives-formal-modification}
$S$, $f : X' \to X$, $T \subset |X|$, $U \subset X$,
$T' \subset |X'|$, および $U' \subset X'$ を Section
\ref{restricted-section-completion-functor} の通りとする。
$X$ が局所 Noether，$f$ が proper，かつ $U' \to U$ が同型ならば，
$f_{/T} : X'_{/T'} \to X_{/T}$ は形式的修正である。
\end{lemma}

\begin{proof}
% P10-E0730
Formal Spaces, Lemmas \ref{formal-spaces-lemma-formal-completion-types}
により，この矢印の始域と終域は局所 Noether 形式代数空間である。
その他の条件は Lemmas \ref{restricted-lemma-proper-gives-proper},
\ref{restricted-lemma-etale-gives-rig-etale},
\ref{restricted-lemma-completion-proper-surjective-rig-surjective}, および
\ref{restricted-lemma-separated-mono-open-diagonal-rig-surjective} から従う。
\end{proof}

\begin{lemma}
\label{restricted-lemma-base-change-formal-modification}
$S$ をスキームとする。$f : X \to Y$ を，$S$ 上の局所 Noether
形式代数空間の射であって形式的修正であるものとする。
局所 Noether 形式代数空間の任意の adic 射 $Y' \to Y$ に対し，
基底変換 $f' : X \times_Y Y' \to Y'$ は形式的修正である。
\end{lemma}

\begin{proof}
Formal Spaces, Lemma \ref{formal-spaces-lemma-base-change-proper}
により $f'$ は proper である。
% P10-E0731
Lemma \ref{restricted-lemma-base-change-rig-etale} により $f'$ は rig-etale である。
% P10-E0732
次に射 $f'$ は Lemma \ref{restricted-lemma-base-change-rig-surjective}
により rig-全射である。
% P10-E0733
$X' = X \times_ Y'$ と置く。射 $\Delta_{f'}$ は，adic 射
$X' \times_{Y'} X' \to X \times_Y X$ による $\Delta_f$ の基底変換である。
従って Lemma \ref{restricted-lemma-base-change-rig-surjective}
により $\Delta_{f'}$ は rig-全射である。
\end{proof}

\section{完備化と射，I}
\label{restricted-section-completion-and-morphisms}

\noindent
この節では，Theorem \ref{restricted-theorem-dilatations-general} の証明で用いる
完備化に関する準備的結果をいくつか与える。ここで述べ証明する補題は
自明ではないが（いくつかは rig-\'etale 代数の代数化に関する成果に基づく），
それでも初読時にはこの節を飛ばすことを勧める。

\begin{lemma}
\label{restricted-lemma-algebraize-rig-etale-affine}
$T \subset X$ を Noether アフィンスキーム $X$ の閉部分集合とする。
$W$ を Noether アフィン形式代数空間とする。
$g : W \to X_{/T}$ を rig-\'etale 射とする。このときアフィンスキーム
$X'$ と，$X \setminus T$ 上 \'etale である有限型射
$f : X' \to X$ が存在し，$f_{/T}$ および $g$ と整合する同型
$X'_{/f^{-1}T} \cong W$ が存在する。さらに，
$W \to X_{/T}$ が \'etale ならば，$X' \to X$ も \'etale である。
\end{lemma}

\begin{proof}
$X'$ の存在は Lemma
\ref{restricted-lemma-approximate-by-etale-over-complement} の言い換えである。
最後の主張は More on Morphisms, Lemma
\ref{more-morphisms-lemma-check-smoothness-on-infinitesimal-nbhds}
から従う。
\end{proof}

\begin{lemma}
\label{restricted-lemma-algebraize-morphism-rig-etale}
次を仮定する。
\begin{enumerate}
\item $X$, $X'$, $Y$ は Noether アフィンスキームである，
\item $T \subset |X|$ は閉部分集合である，
\item $f : X' \to X$ は局所有限型の射であり，
$X \setminus T$ 上 \'etale である，
\item $h : Y \to X$ は射である，
\item $\alpha : Y_{/T} \to X'_{/T}$ は $X_{/T}$ 上の射である
（記法については証明を参照せよ）。
\end{enumerate}
このとき，同型 $b_{/T} : Y'_{/T} \to Y_{/T}$ を誘導する
アフィンスキームの \'etale 射 $b : Y' \to Y$ と，$X$ 上の射
$a : Y' \to X'$ が存在して，
$\alpha = a_{/T} \circ b_{/T}^{-1}$ を満たす。
\end{lemma}

\begin{proof}
主張中の添字 ${}_{/T}$ を用いる記法は，スキームの射
$g : V \to X$ に形式代数空間の射
$g_{/T} : V_{/g^{-1}T} \to X_{/T}$ を対応させる構成を指す。
これは $X$ 上のスキームの圏から $X_{/T}$ 上の形式代数空間の圏への関手である。
Section \ref{restricted-section-completion-functor} を参照せよ。
この説明の下で，補題は Lemma
\ref{restricted-lemma-fully-faithful-etale-over-complement} の単なる言い換えである。
\end{proof}

\begin{lemma}
\label{restricted-lemma-factor}
% P10-E0734
$S$ をスキームとする。$f : X \to Y$ および $g : Z \to Y$ を
代数空間の射とする。$T \subset |X|$ を閉部分集合とする。次を仮定する。
\begin{enumerate}
\item $X$ は局所 Noether である，
\item $g$ は局所有限型のモノ射である，
\item $f|_{X \setminus T} : X \setminus T \to Y$ は $g$ を経由する，
\item $f_{/T} : X_{/T} \to Y$ は $g$ を経由する。
\end{enumerate}
このとき $f$ は $g$ を経由する。
\end{lemma}

\begin{proof}
ファイバー積 $E = X \times_Y Z \to X$ を考える。
仮定により開埋込み $X \setminus T \to X$ は $E$ を経由し，
$|\varphi|(|X'|) \subset T$ を満たす任意の射
$\varphi : X' \to X$ も $E$ を経由する。Formal Spaces, Section
\ref{formal-spaces-section-completion} を参照せよ。
More on Morphisms of Spaces, Lemma
\ref{spaces-more-morphisms-lemma-check-smoothness-on-infinitesimal-nbhds}
により，これは $E \to X$ が $T$ の点へ写る $E$ のすべての点で
\'etale であることを含意する。従って $E \to X$ は \'etale モノ射，
ゆえに開埋込みである（Morphisms of Spaces, Lemma
\ref{spaces-morphisms-lemma-etale-universally-injective-open}）。
このとき仮定から $|X| = |E|$ なので，$E = X$ が従う。
\end{proof}

\begin{lemma}
\label{restricted-lemma-faithful-general}
$S$ をスキームとする。$X$, $W$ を $S$ 上の代数空間とし，
$X$ は局所 Noether であるとする。$T \subset |X|$ を閉部分集合とする。
$a, b : X \to W$ を $S$ 上の代数空間の射であって，
$a|_{X \setminus T} = b|_{X \setminus T}$ かつ，
$X_{/T} \to W$ の射として $a_{/T} = b_{/T}$ を満たすものとする。
このとき $a = b$ である。
\end{lemma}

\begin{proof}
$E$ を $a$ と $b$ の等化子とする。このとき $E$ は代数空間であり，
$E \to X$ は局所有限型のモノ射である。Morphisms of Spaces, Lemma
\ref{spaces-morphisms-lemma-properties-diagonal} を参照せよ。
仮定により，二つの射 $f = \text{id} : X \to X$ と
$g : E \to X$，および $|X|$ の閉部分集合 $T$ に
Lemma \ref{restricted-lemma-factor} を適用できる。
\end{proof}

\begin{lemma}
\label{restricted-lemma-faithful}
$S$ をスキームとする。$X$, $Y$ を $S$ 上の局所 Noether 代数空間とする。
$T \subset |X|$ および $T' \subset |Y|$ を閉部分集合とする。
$a, b : X \to Y$ を $S$ 上の代数空間の射であって，
$a|_{X \setminus T} = b|_{X \setminus T}$,
$|a|(T) \subset T'$, $|b|(T) \subset T'$，かつ
$X_{/T} \to Y_{/T'}$ の射として $a_{/T} = b_{/T}$ を満たすものとする。
このとき $a = b$ である。
\end{lemma}

\begin{proof}
より一般的な Lemma \ref{restricted-lemma-faithful-general} の帰結である。
\end{proof}

\begin{lemma}
\label{restricted-lemma-equivalence-relation}
$S$ をスキームとする。$X$ を $S$ 上の局所 Noether 代数空間とする。
$T \subset |X|$ を閉部分集合とする。
$s, t : R \to U$ を $X$ 上の代数空間の二つの射とする。次を仮定する。
\begin{enumerate}
\item $R$, $U$ は $X$ 上局所有限型である，
\item $s$ と $t$ の $X \setminus T$ への基底変換は \'etale 同値関係である，
\item 形式完備化
$(t_{/T}, s_{/T}) : R_{/T} \to U_{/T} \times_{X_{/T}} U_{/T}$
も同値関係である（記法については証明を参照せよ）。
\end{enumerate}
このとき $(t, s) : R \to U \times_X U$ は \'etale 同値関係である。
\end{lemma}

\begin{proof}
主張中の添字 ${}_{/T}$ を用いる記法は，代数空間の射
$f : X' \to X$ に形式代数空間の射
$f_{/T} : X'_{/f^{-1}T} \to X_{/T}$ を対応させる構成を指す。
Section \ref{restricted-section-completion-functor} を参照せよ。
仮定により射 $s, t : R \to U$ は $X \setminus T$ 上 \'etale である。
写像 $s, t : R \to U$ の形式完備化も \'etale なので，例えば
More on Morphisms, Lemma
\ref{more-morphisms-lemma-check-smoothness-on-infinitesimal-nbhds}
により $s$ と $t$ は \'etale である。射
$\text{id} : R \times_{U \times_X U} R \to R \times_{U \times_X U} R$
および $\Delta : R \to R \times_{U \times_X U} R$ に
Lemma \ref{restricted-lemma-factor} を適用すると，$(t, s)$ がモノ射であると分かる。
さらに
$(t \circ \text{pr}_0, s \circ \text{pr}_1) :
R \times_{s, U, t} R \to U \times_X U$ および
$(t, s) : R \to U \times_X U$ に再び適用すると，
「推移律」が成り立つと分かる。同値関係の他の二公理の証明は省略する。
\end{proof}

\begin{lemma}
\label{restricted-lemma-smash-away-from-T}
$S$ をスキームとする。$X$ を $S$ 上の局所 Noether 代数空間とし，
$T \subset |X|$ を閉部分集合とする。$f : X' \to X$ を，
局所有限型で $T$ の外で \'etale である代数空間の射とする。
$f$ の分解
$$
X' \longrightarrow X'' \longrightarrow X
$$
であって，次の性質をもつものが存在する：
$X'' \to X$ は局所有限型であり，
$X'' \to X$ は $X \setminus T$ 上の同型であり，
$X'_{/T} \to X''_{/T}$ は同型である（記法については証明を参照せよ）。
\end{lemma}

\begin{proof}
主張中の添字 ${}_{/T}$ を用いる記法は，代数空間の射
$f : X' \to X$ に形式代数空間の射
$f_{/T} : X'_{/f^{-1}T} \to X_{/T}$ を対応させる構成を指す。
Section \ref{restricted-section-completion-functor} を参照せよ。
また Section \ref{restricted-section-completion-functor} で導入した
$|U| = |X| \setminus T$ および
% P10-E0735
$U' = |X'| \setminus f^{-1}T$ を満たす開部分空間を表すため，
$U \subset X$ および $U' \subset X'$ という概念も用いる。

\medskip\noindent
$X'$ を $X' \amalg U$ で置き換えることにより，
$X' \to X$ の像が $U$ を含むと仮定してよく，そう仮定する。
$U' \to X'$ と対角射
$U' \to U' \times_U U' = U' \times_X U'$ の押出しを
$$
R = X' \amalg_{U'} (U' \times_U U')
$$
とする。$U' \to X$ は \'etale なので，この対角射は開埋込みであり，
$R$ は代数空間である（これは例えば Spaces, Lemma
\ref{spaces-lemma-glueing-algebraic-spaces} から従う）。
二つの射影 $U' \times_U U' \to U'$ は $R$ へ延長され，
二つの \'etale 射 $s, t : R \to X'$ を得る。
各部分上で別々に確認すると，$R$ は $X'$ 上の \'etale 同値関係である。
$X'' = X'/R$ と置く。Bootstrap, Theorem
\ref{bootstrap-theorem-final-bootstrap} により，これは代数空間である。
% P10-E0736
構成により補題の分解をもち，射 $X'' \to X$ は局所有限型である
（これは \'etale 局所的，すなわち $X'$ 上で確認できる）。
$U' \to U$ は全射 \'etale 射であり，
$s^{-1}(U') = t^{-1}(U') = U' \times_U U'$ なので，
$U'' = U \times_X X'' \to U$ は同型である。
最後に，射 $X' \to X''$ が同型
$X'_{/T} \to X''_{/T}$ を誘導することを示さなければならない。
これを見るため，$R$ の $T$ の逆像に沿う形式完備化が，
$R$ の選び方により $X'$ の $T$ の逆像に沿う形式完備化に等しいことに注意する！
Formal Spaces, Section \ref{formal-spaces-section-completion}
における形式完備化の構成により，層として
$X''_{/T} = (X'_{/T}) / (R_{/T})$ である。
$X'_{/T} = R_{/T}$ なので $X'_{/T} = X''_{/T}$ と結論でき，
証明が完了する。
\end{proof}

\section{射の rig 貼合せ}
\label{restricted-section-glue-formal-to-actual}

\noindent
$X$, $W$ を代数空間とし，$X$ は Noether であるとする。
$Z \subset X$ を，開補集合 $U$ をもつ閉部分空間とする。
下の命題は大まかに言えば
$$
\{\text{射 }X \to W\} =
\{\text{両立する射 }U \to W\text{ と }X_{/Z} \to W\}
$$
を主張する。ここで $a : U \to W$ と $b : X_{/Z} \to W$ が
両立するとは，$a$ と $b$ が同じ「rig 空間の射」を定めることを意味する。
「rig 空間」の圏を導入するには多くの作業が必要だが，この場合に条件が
何を意味するかを正確に述べるために，それを行う必要はない。

\begin{proposition}
\label{restricted-proposition-glue-modification}
$S$ をスキームとする。$X$ を $S$ 上の局所 Noether 代数空間とする。
$T \subset |X|$ を，補開部分空間 $U \subset X$ をもつ閉部分集合とする。
$f : X' \to X$ を，$f^{-1}(U) \to U$ が同型であるような
代数空間の proper 射とする。$S$ 上の任意の代数空間 $W$ に対して，写像
$$
\Mor_S(X, W) \longrightarrow
\Mor_S(X', W) \times_{\Mor_S(X'_{/T}, W)} \Mor_S(X_{/T}, W)
$$
は全単射である。
\end{proposition}

\begin{proof}
$w' : X' \to W$ と $\hat w : X_{/T} \to W$ を，
% P10-E0737
$X'_{/T} \to X$ および $X'_{/T} \to X_{/T}$ との合成により
同じ射 $X'_{/T} \to W$ を定める射とする。
$X' \to X$ および $X_{/T} \to X$ との合成が $w'$ と $\hat w$ を
それぞれ再現する一意な射 $w : X \to W$ が存在することを示さなければならない。
一意性は Lemma \ref{restricted-lemma-faithful-general} から直ちに従う。

\medskip\noindent
$T$ と $f$ に関する仮定は，代数空間の任意の \'etale 射
$X_1 \to X$ による基底変換で保たれる。
形式代数空間は \'etale 位相に関する層であり，しかも一意性は既に得られているので，
$X$ をある \'etale 被覆の各要素で置き換えた後に存在を証明すれば十分である。
従って $X$ はアフィン Noether スキームであると仮定してよい。

\medskip\noindent
$X$ はアフィン Noether スキームであると仮定する。
Pushouts of Spaces, Section
\ref{spaces-pushouts-section-coequalizer-glue} の内容を用いて
射 $w : X \to W$ を構成する。
% P10-E0738
証明を読み進める前に，その節の内容を少し読むのが有益である。

\medskip\noindent
$X'' = X' \times_X X'$ と置き，二つの射
$a = w' \circ \text{pr}_1 : X'' \to W$ および
$b = w' \circ \text{pr}_2 : X'' \to W$ を考える。
$a$ と $b$ は開部分 $U$ 上で一致し，さらに
% P10-E0739
$a_{/T}$ と $b_{a/T}$ も一致することが分かる
（いずれも合成 $X''_{/T} \to X_{/T} \to W$ に等しく，
第二の矢印は $\hat w$ である）。従って Lemma
\ref{restricted-lemma-faithful-general} により $a = b$ である。

\medskip\noindent
$T$ 上に誘導される被約閉部分スキーム構造を $Z \subset X$ と書く。
$n \geq 1$ に対して，$Z$ の第 $n$ 無限小近傍を $Z_n \subset X$ と書く。
$w_n = \hat w|_{Z_n} : Z_n \to W$ と書けば，
$X_{/T} = \colim Z_n$ 上で $\hat w = \colim w_n$ である。
$Y_n = X' \amalg Z_n$ と置く。二つの射影
$$
s_n, t_n : R_n = Y_n \times_X Y_n \longrightarrow Y_n
$$
を考える。$s_n$ と $t_n$ の余等化子を Pushouts of Spaces, Section
\ref{spaces-pushouts-section-coequalizer-glue} の通りに
$Y_n \to X_n \to X$ とする（特にこの余等化子は存在し，良い性質をもつ等々。
Pushouts of Spaces, Lemma
\ref{spaces-pushouts-lemma-coequalizer} を参照せよ）。
% P10-E0740
前段落で得た $a = b$ と，$X'_{/T}$ 上で $w'$ と $\hat w$ が一致することから，射
$$
w' \amalg w_n : Y_n \longrightarrow W
$$
は射 $s_n$ と $t_n$ を等化する。従ってすべての $n \geq 1$ に対して，
$w'$ および $w_n$ と両立する射 $w^n : X_n \to W$ が存在する。
さらに $m \geq 1$ に対し，合成
$$
X_n \to X_{n + m} \xrightarrow{w^{n + m}} W
$$
は構成により $w^n$ に等しい（対応する主張が
$w' \amalg w_{n + m}$ と $w' \amalg w_n$ について成り立つからである）。
Pushouts of Spaces, Lemma
\ref{spaces-pushouts-lemma-essentially-constant} および Remark
\ref{spaces-pushouts-remark-essentially-constant} により，
代数空間の系 $X_n$ は値 $X$ をもって本質的に定常である。これで結論を得る。
\end{proof}

\section{rig-\'etale 射の代数化}
\label{restricted-section-algebraization}

\noindent
この節では，Artin の論文 \cite{ArtinII} にある dilatation に関する結果の
一般化を証明する。

\medskip\noindent
この節の記法は Section \ref{restricted-section-completion-functor} の記法と一致させるが，
代数空間と形式代数空間は局所 Noether であるとする。

\medskip\noindent
従ってまず基礎スキーム $S$ を固定する。すべての環，位相環，スキーム，
代数空間，形式代数空間，およびこれらの間の射は $S$ 上にあるものとする。
次に，局所 Noether 代数空間 $X$ と閉部分集合 $T \subset |X|$ を固定する。
% P10-E0741
$|U| = |X| \setminus T$ を満たす開部分空間 $U \subset X$ を考える。
図式は
$$
U \to X \quad |X| = |U| \amalg T
$$
である。代数空間の射 $f : X' \to X$ が与えられたとき，Section
\ref{restricted-section-completion-functor} と同じく
$U' = f^{-1}U$, $T' = |f|^{-1}(T)$, および
$f_{/T} : X'_{/T'} \to X_{/T}$ という記法を用いる。
% P10-E0742
ときには $X'_{/T'}$ の代わりに $X'_{/T}$ と書く。さらに一般に，
$X$ 上の代数空間の射 $a : X' \to X''$ に対し，$X'$ と $X''$ における
$T$ の逆像に沿って射 $a$ を完備化して得られる形式代数空間の射を
$a_{/T} : X'_{/T} \to X''_{/T}$ と書く。

\medskip\noindent
この設定の下で，関手
\begin{equation}
\label{restricted-equation-completion-functor}
\left\{
\begin{matrix}
\text{代数空間の射}\\
f : X' \to X\text{ であって局所的に}\\
\text{有限型で，かつ}\\
U' \to U\text{ は同型である}
\end{matrix}
\right\}
\longrightarrow
\left\{
\begin{matrix}
\text{射 }g : W \to X_{/T}\\
\text{（形式代数空間の射）}\\
\text{ここで }W\text{ は局所 Noether である}\\
\text{かつ }g\text{ は rig-\'etale である}
\end{matrix}
\right\}
\end{equation}
を考える。この関手は $f : X' \to X$ を
$f_{/T} : X'_{/T'} \to X_{/T}$ に送る。
Lemma \ref{restricted-lemma-etale-gives-rig-etale} により $f_{/T}$ は
rig-\'etale なので，これは意味をもつ。

\begin{lemma}
\label{restricted-lemma-functoriality-completion-functor}
上の状況で，$X_1$ が局所 Noether であるような代数空間の射
$X_1 \to X$ を取る。$T$ の逆像を $T_1 \subset |X_1|$，
$U$ の逆像を $U_1 \subset X_1$ と書く。次の記法を用いる。
\begin{enumerate}
\item $\mathcal{C}_{X, T}$ は，局所有限型であって
$U' = f^{-1}U \to U$ が同型であるような代数空間の射
$f : X' \to X$ を対象とする圏である，
\item $\mathcal{C}_{X_1, T_1}$ は，局所有限型であって
$f_1^{-1}U_1 \to U_1$ が同型であるような代数空間の射
$f_1 : X_1' \to X_1$ を対象とする圏である，
\item $\mathcal{C}_{X_{/T}}$ は，$W$ が局所 Noether で $g$ が
rig-\'etale であるような形式代数空間の射
$g : W \to X_{/T}$ を対象とする圏である，
% P10-E0743
\item $\mathcal{C}_{X_{1, /T_1}}$ は，$W_1$ が局所 Noether で
$g_1$ が rig-\'etale であるような形式代数空間の射
$g_1 : W_1 \to X_{1, /T_1}$ を対象とする圏である。
\end{enumerate}
このとき図式
$$
\xymatrix{
\mathcal{C}_{X, T} \ar[d] \ar[r] &
\mathcal{C}_{X_{/T}} \ar[d] \\
\mathcal{C}_{X_1, T_1} \ar[r] &
\mathcal{C}_{X_{1, /T_1}}
}
$$
は可換である。ここで横の矢印は
(\ref{restricted-equation-completion-functor}) により，縦の矢印は
$X_1 \to X$ および $X_{1, /T_1} \to X_{/T}$ に沿う基底変換により与えられる。
\end{lemma}

\begin{proof}
これは，$X$ 上の代数空間の圏上の完備化関手
$(h : Y \to X) \mapsto Y_{/T} = Y_{/|h|^{-1}T}$ が
ファイバー積と可換であることから直ちに従う。
\end{proof}

\begin{lemma}
\label{restricted-lemma-completion-functor-fully-faithful}
上の状況で，$f : X' \to X$ を，局所有限型で $U$ 上の同型である
代数空間の射とする。$g : Y \to X$ を，$Y$ が局所 Noether である射とする。
このとき完備化は全単射
$$
\Mor_X(Y, X') \longrightarrow \Mor_{X_{/T}}(Y_{/T}, X'_{/T})
$$
を定める。特に，関手 (\ref{restricted-equation-completion-functor}) は充満忠実である。
\end{lemma}

\begin{proof}
$a, b : Y \to X'$ を，$a_{/T} = b_{/T}$ を満たす $X$ 上の射とする。
このとき $a$ と $b$ は開部分空間 $g^{-1}U$ 上，かつ $g^{-1}T$ に沿う
完備化の後に一致する。従って Lemma \ref{restricted-lemma-faithful}
により $a = b$ である。言い換えると，完備化写像は常に単射である。

\medskip\noindent
$\alpha : Y_{/T} \to X'_{/T}$ を $X_{/T}$ 上の形式代数空間の射とする。
$\alpha = a_{/T}$ を満たす $X$ 上の射 $a : Y \to X'$ が存在することを
示さなければならない。証明は標準的だが煩雑なアフィンの場合への還元を行い，
その後 Lemma \ref{restricted-lemma-algebraize-morphism-rig-etale} を適用する。

\medskip\noindent
$\{h_i : Y_i \to Y\}$ を代数空間の \'etale 被覆とする。
各 $i$ に対して，その完備化
$(a_i)_{/T} : (Y_i)_{/T} \to X'_{/T}$ が
$\alpha \circ (h_i)_{/T}$ に等しいような $X$ 上の射
$a_i : Y_i \to X'$ を見いだせるならば，
$\alpha = a_{/T}$ を満たす射 $a : Y \to X'$ を得る。
実際，まず $\alpha$ との一致により，
$(Y_i \times_Y Y_j)_{/T} \to X'_{/T}$ の射として
$(a_i)_{/T} \circ \text{pr}_1 = (a_j)_{/T} \circ \text{pr}_2$
であることに注意する（ここでは完備化 ${}_{/T}$ がファイバー積と可換であることを用いる）。
既に証明した単射性により，これは $Y_i \times_Y Y_j \to X'$ の射として
$a_i \circ \text{pr}_1 = a_j \circ \text{pr}_2$ であることを示す。
$X'$ は fppf 層なので，射の族 $a_i$ は射 $a : Y \to X'$ に降下する。
% P10-E0744
$\{(a_i)_{/T} : (Y_i)_{/T} \to X'_{/T}\}$ は \'etale 被覆なので，
$\alpha = a_{/T}$ である。

\medskip\noindent
前段落の結果により，存在を証明するには $Y$ がアフィンで，
$g : Y \to X$ が $g_1 : Y \to X_1$ とアフィン $X_1$ をもつ
\'etale 射 $X_1 \to X$ との合成に分解すると仮定してよい。
このとき $T$ の逆像 $T_1 \subset |X_1|$ を考え，
$X'_1 = X' \times_X X_1$ と置き，射影を $f_1 : X'_1 \to X_1$ と書く。
さらに
$$
\alpha_1 = (\alpha, (g_1)_{/T_1}) :
Y_{/T_1} = Y_{/T}
\longrightarrow
X'_{/T} \times_{X_{/T}} (X_1)_{/T_1} = (X'_1)_{/T_1}
$$
と置く。$X_1$ 上の $\alpha_1$ について存在を証明すれば十分である。
% P10-E0745
言い換えると，$X, T, X', Y, f, g, \alpha$ を
$X_1, T_1, X'_1, Y, g_1, \alpha_1$ で置き換えてよい。
これにより次段落で述べる場合へ還元される。

\medskip\noindent
$Y$ と $X$ はアフィンであると仮定する。$(Y_{/T})_{red}$ は
アフィンスキームであることを思い出そう（$g^{-1}T \subset Y$ 上に誘導される
被約スキーム構造と同型である。Formal Spaces, Lemma
\ref{formal-spaces-lemma-reduction-completion} を参照せよ）。
従って $\alpha_{red} : (Y_{/T})_{red} \to (X'_{/T})_{red}$ は
$f^{-1}T$ 内に準コンパクトな像 $E$ をもつ（上と同じ補題により，これは
$(X'_{/T})_{red}$ の基礎位相空間である）。従ってアフィンスキーム $V$ と，
% P10-E0746
$h$ の像が $E$ を含む \'etale 射 $h : V \to X'$ を見いだせる。
実線部分が Cartesian である図式
$$
\xymatrix{
Y'_{/T} \ar@{..>}[rd] \ar@{..>}[r] &
W \ar[d] \ar[r] & V_{/T} \ar[d]^{h_{/T}} \\
& Y_{/T} \ar[r]^\alpha & X'_{/T}
}
$$
を選ぶ。構成により射 $W \to Y_{/T}$ は代数空間によって表現可能で，
\'etale かつ全射である（全射性は被約化を見れば分かる。Formal Spaces, Lemma
\ref{formal-spaces-lemma-reduction-surjective} を参照せよ）。
Lemma \ref{restricted-lemma-algebraize-rig-etale-affine} により，
$Y' \to Y$ が \'etale かつ $Y'$ がアフィンであるように
$W = Y'_{/T}$ と書ける。これにより図式中の点線矢印が得られる。
$W \to Y_{/T}$ は全射なので，$Y' \to Y$ の像は $g^{-1}T$ を含む。
従って $\{Y' \to Y, Y \setminus g^{-1}T \to Y\}$ は \'etale 被覆である。
$f$ は $U$ 上の同型なので，$X$ 上の（一意な）射
$Y \setminus g^{-1}T \to X'$ が存在し，完備化上で $\alpha$ と一致する
（$Y \setminus g^{-1}T$ の完備化は空だからである）。
従って $Y'$ について存在を証明すれば十分であり，次段落で調べる場合へ還元される。

\medskip\noindent
前段落の結果により，$Y$ がアフィンであり，$\alpha$ が
$Y_{/T} \to V_{/T} \to X'_{/T}$ と分解すると仮定してよい。
ここで $V$ は $X'$ 上 \'etale なアフィンスキームである。
さらに $Y$ をアフィン \'etale 被覆の各要素で置き換えてよい。
Lemma \ref{restricted-lemma-algebraize-morphism-rig-etale} により，
同型 $b_{/T} : Y'_{/T} \to Y_{/T}$ を誘導するアフィンスキームの
\'etale 射 $b : Y' \to Y$ と，$c_{/T} \circ b_{/T}^{-1}$ が与えられた射
$Y_{/T} \to V_{/T}$ に等しいような射 $c : Y' \to V$ を見いだせる。
$a' : Y' \to X'$ を $c$ と $V \to X'$ の合成とすると，
$a'_{/T} = \alpha \circ b_{/T}$ である。言い換えると，
$Y'$ と $\alpha \circ b_{/T}$ について存在が得られる。
% P10-E0747
そこで \'etale 被覆 $\{Y' \to Y, Y \setminus g^{-1}T \to Y\}$
をもう一度考えることにより証明が完了する。
\end{proof}

\begin{lemma}
\label{restricted-lemma-dilatations-affine}
上の状況で $X$ がアフィンであると仮定する。このとき関手
(\ref{restricted-equation-completion-functor}) は圏同値である。
\end{lemma}

\noindent
この補題を証明する前に，一つの例を論じよう。
$S = \Spec(k)$, $X = \mathbf{A}^1_k$, $T = \{0\}$ と仮定する。
このとき $X_{/T} = \text{Spf}(k[[x]])$ である。
$W = \text{Spf}(k[[x]] \times k[[x]])$ とする。
すると対応する $f : X' \to X$ は，原点を二重化したアフィン直線から
アフィン直線への写像である（Schemes, Example
\ref{schemes-example-affine-space-zero-doubled}）。
さらにこれは，$X$ 上の $X \amalg X$ から出発する Lemma
\ref{restricted-lemma-smash-away-from-T} の構成の出力である。

\begin{proof}
Lemma \ref{restricted-lemma-completion-functor-fully-faithful} により，
関手が充満忠実であることは既に分かっている。本質的全射性を示す。
$g : W \to X_{/T}$ を，$W$ が局所 Noether で $g$ が rig-\'etale である
形式代数空間の射とする。いくつかの段階に分けて，$W$ が本質像に属することを示す。

\medskip\noindent
Step 1: $W$ がアフィン形式代数空間である場合。
% P10-E0748
Lemma \ref{restricted-lemma-algebraize-rig-etale-affine} により，有限型で
$X \setminus T$ 上 \'etale であり，$U_{/T}$ が $W$ と同型であるような
$U \to X$ を見いだせる。従って Lemma
\ref{restricted-lemma-smash-away-from-T} により $W$ は本質像に属する。

\medskip\noindent
Step 2: $W$ が分離的である場合。Formal Spaces, Definition
\ref{formal-spaces-definition-formal-algebraic-space} の通りに
$\{W_i \to W\}$ を選ぶ。Step 1 により，形式代数空間
$W_i$ と $W_i \times_W W_j$ は本質像に属する。
$W_i = (X'_i)_{/T}$ および
$W_i \times_W W_j = (X'_{ij})_{/T}$ と書く。
% P10-E0749
充満忠実性により，射影 $W_i \times_W W_j \to W_i$ および
$W_i \times_W W_j \to W_j$ に対応する射
$t_{ij} : X'_{ij} \to X'_i$ および
$s_{ij} : X'_{ij} \to X'_j$ を得る。構造
$$
R = \coprod X'_{ij},\quad
V = \coprod X'_i,\quad
s = \coprod s_{ij},\quad
t = \coprod t_{ij}
$$
を考える（文字 $U$ は既に用いられているので使えない）。
Lemma \ref{restricted-lemma-equivalence-relation} を適用すると，
$(t, s) : R \to V \times_X V$ は $X$ 上の $V$ に \'etale 同値関係を定める。
従って商 $X' = V/R$ を取ることができ，Bootstrap, Theorem
\ref{bootstrap-theorem-final-bootstrap} により，これは代数空間である。
完備化はファイバー積および商層を取る操作と可換なので，
$X_{/T}$ 上の形式代数空間として $X'_{/T} \cong W$ を得る。

\medskip\noindent
Step 3: $W$ が一般の場合。Formal Spaces, Definition
\ref{formal-spaces-definition-formal-algebraic-space} の通りに
$\{W_i \to W\}$ を選ぶ。形式代数空間 $W_i$ と
$W_i \times_W W_j$ は分離的である。従って Step 2 により，
形式代数空間 $W_i$ と $W_i \times_W W_j$ は本質像に属する。
そこで前段落とまったく同じ議論を行えば，$W$ も本質像に属することが分かる。
これで証明は完了する。
\end{proof}

\begin{theorem}
\label{restricted-theorem-dilatations-general}
$S$ をスキームとする。$X$ を $S$ 上の局所 Noether 代数空間とする。
$T \subset |X|$ を閉部分集合とする。$U \subset X$ を
$|U| = |X| \setminus T$ を満たす開部分空間とする。完備化関手
(\ref{restricted-equation-completion-functor})
$$
\left\{
\begin{matrix}
\text{代数空間の射}\\
f : X' \to X\text{ であって局所的に}\\
\text{有限型で，かつ}\\
f^{-1}U \to U\text{ は同型である}
\end{matrix}
\right\}
\longrightarrow
\left\{
\begin{matrix}
\text{射 }g : W \to X_{/T}\\
\text{（形式代数空間の射）}\\
\text{ここで }W\text{ は局所 Noether である}\\
\text{かつ }g\text{ は rig-\'etale である}
\end{matrix}
\right\}
$$
は，$f : X' \to X$ を $f_{/T} : X'_{/T'} \to X_{/T}$ に送り，圏同値である。
\end{theorem}

\begin{proof}
Lemma \ref{restricted-lemma-completion-functor-fully-faithful} により，この関手は充満忠実である。
$g : W \to X_{/T}$ を，$W$ が局所 Noether で $g$ が rig-\'etale である
形式代数空間の射とする。証明を終えるため，$W$ が本質像に属することを示す。

\medskip\noindent
各 $i$ に対して $X_i$ がアフィンであるような \'etale 被覆
$\{X_i \to X\}$ を選ぶ。$U$ の逆像を $U_i \subset X_i$，
$T$ の逆像を $T_i \subset X_i$ と書く。
$(X_i)_{/T_i} = (X_i)_{/T} = (X_i \times_X X)_{/T}$ および
$W_i = X_i \times_X W = (X_i)_{/T} \times_{X_{/T}} W$ であることを
思い出そう。Lemma \ref{restricted-lemma-functoriality-completion-functor} を参照せよ。
適切な余サイクル条件を満たす同型
$$
\alpha_{ij} :
W_i \times_{X_{/T}} (X_j)_{/T}
\longrightarrow 
(X_i)_{/T} \times_{X_{/T}} W_j
$$
を得る。Lemma \ref{restricted-lemma-dilatations-affine} を
$X_i, T_i, U_i, W_i \to (X_i)_{/T}$ に適用すると，
局所有限型で $U_i$ 上の同型である代数空間の射
$X'_i \to X_i$ と，$(X_i)_{/T}$ 上の同型
$\beta_i : (X'_i)_{/T} \cong W_i$ が存在する。
% P10-E0750
充満忠実性により，$X_i \times_X X_j$ 上の同型
$$
a_{ij} : X'_i \times_X X_j \longrightarrow X_i \times_X X'_j
$$
であって
$\alpha_{ij} = \beta_j|_{X_i \times_X X_j}
\circ (a_{ij})_{/T} \circ \beta_i^{-1}|_{X_i \times_X X_j}$
を満たすものを得る。再び充満忠実性を用いると（今度は
$X_i \times_X X_j \times_X X_k$ 上で），これらの射 $a_{ij}$ は
$\alpha_{ij}$ が満たすものと同じ余サイクル条件を満たすことが分かる。
言い換えると，族 $\{X_i \to X\}$ に関する降下データ
（Descent on Spaces, Definition
\ref{spaces-descent-definition-descent-datum-for-family-of-morphisms} の意味で）
$(X'_i, a_{ij})$ を得る。Bootstrap, Lemma
\ref{bootstrap-lemma-descend-algebraic-space} により，この降下データは有効である。
従って代数空間の射 $f : X' \to X$ と，$X_i$ 上の同型
$h_i : X' \times_X X_i \to X'_i$ であって
$a_{ij} = h_j|_{X_i \times_X X_j} \circ h_i^{-1}|_{X_i \times_X X_j}$
を満たすものを得る。読者は，これにより得られる $X_i$ 上の同型
$$
(X' \times_X X_i)_{/T}
\xrightarrow{\beta_i \circ (h_i)_{/T}}
W_i
$$
が貼り合わさって $X_{/T}$ 上の同型 $X'_{/T} \to W$ を与えることを
確認できる。詳細の一部は省略する。
\end{proof}

\section{完備化と射，II}
\label{restricted-section-completion-and-morphisms-bis}

\noindent
Artin の dilatation 定理を得るには，Theorem
\ref{restricted-theorem-dilatations-general} が与える対応の下で，形式的修正と
実際の修正とを対応させる必要がある。読者にはこの節を飛ばすことを勧める。

\begin{lemma}
\label{restricted-lemma-output-quasi-compact}
Theorem \ref{restricted-theorem-dilatations-general} の仮定と記法の下で，
$f : X' \to X$ が $g : W \to X_{/T}$ に対応するとする。
このとき $f$ が準コンパクトであるための必要十分条件は，
$g$ が準コンパクトであることである。
\end{lemma}

\begin{proof}
$f$ が準コンパクトならば，Lemma
\ref{restricted-lemma-quasi-compact-gives-quasi-compact} により $g$ も準コンパクトである。
逆に $g$ が準コンパクトであると仮定する。
$X_i$ がアフィンであるような \'etale 被覆 $\{X_i \to X\}$ を選ぶ。
Morphisms of Spaces, Lemma
\ref{spaces-morphisms-lemma-quasi-compact-local} により，基底変換
$X' \times_X X_i \to X_i$ が準コンパクトであることを示せば十分である。
% P10-E0751
Formal Spaces, Lemma
\ref{formal-spaces-lemma-characterize-quasi-compact-morphism} により，基底変換
$W_i \times_{X_{/T}} (X_i)_{/T} \to (X_i)_{/T}$ は準コンパクトである。
Lemma \ref{restricted-lemma-functoriality-completion-functor} により，
次段落で述べる場合へ還元される。

\medskip\noindent
$X$ はアフィンで，$g : W \to X_{/T}$ は準コンパクトであると仮定する。
$X'$ が準コンパクトであることを示さなければならない。
$V = \coprod_{j \in J} V_j$ がアフィンの非交和であるような全射 \'etale 射
$V \to X'$ を取る。このとき $V_{/T} \to X'_{/T} = W$ は全射 \'etale 射である。
% P10-E0757
$W$ は準コンパクトなので，有限部分集合 $J' \subset J$ であって
$\coprod_{j \in J'} (V_j)_{/T} \to W$ が全射であるものを見いだせる。
すると
$$
U \amalg \coprod\nolimits_{j \in J'} V_j \longrightarrow X'
$$
は全射である（従って $X'$ は準コンパクトである）。
実際，$X'_{/T} = W$ なので $|X'| = |U| \amalg |W_{red}|$ である。
\end{proof}

\begin{lemma}
\label{restricted-lemma-output-quasi-separated}
Theorem \ref{restricted-theorem-dilatations-general} の仮定と記法の下で，
$f : X' \to X$ が $g : W \to X_{/T}$ に対応するとする。
このとき $f$ が準分離的であるための必要十分条件は，
$g$ が準分離的であることである。
\end{lemma}

\begin{proof}
$f$ が準分離的ならば，Lemma
\ref{restricted-lemma-quasi-separated-gives-quasi-separated} により $g$ も準分離的である。
逆に $g$ が準分離的であると仮定する。$f$ が準分離的であることを
示さなければならない。Lemma \ref{restricted-lemma-output-quasi-compact} の証明と
まったく同様に，Morphisms of Spaces, Lemma
\ref{spaces-morphisms-lemma-separated-local} および Formal Spaces, Lemma
\ref{formal-spaces-lemma-separated-local} を用い，アフィンスキームからなる
$X$ の \'etale 被覆の各要素上で確認してよい。
従って $X$ はアフィンであると仮定してよく，そう仮定する。

\medskip\noindent
$V = \coprod_{j \in J} V_j$ がアフィンの非交和であるような全射 \'etale 射
$V \to X'$ を取る。$X'$ が準分離的であることを示すには，
すべての $j, j' \in J$ に対し $V_j \times_{X'} V_{j'}$ が
準コンパクトであることを示せば十分である。$W$ は準分離的なので，
% P10-E0752
すべての $j, j' \in J$ に対しファイバー積
$(V_j \times_Y V_{j'})_{/T} =
(V_j)_{/T} \times_{X'_{/T}} (V_{j'})_{/T}$ は準コンパクトである。
$X$ は Noether アフィンであり，$U' \to U$ は同型なので，
$$
(V_j \times_{X'} V_{j'}) \times_X U =
(V_j \times_X V_{j'}) \times_X U
$$
は準コンパクトである。従って等式
$$
|V_j \times_{X'} V_{j'}| =
|(V_j \times_{X'} V_{j'}) \times_X U| \amalg
|(V_j \times_{X'} V_{j'})_{/T, red}|
$$
と，形式代数空間が準コンパクトであるための必要十分条件は，
その付随被約代数空間が準コンパクトであることだという事実から結論を得る。
\end{proof}

\begin{lemma}
\label{restricted-lemma-output-separated}
Theorem \ref{restricted-theorem-dilatations-general} の仮定と記法の下で，
$f : X' \to X$ が $g : W \to X_{/T}$ に対応するとする。このとき
$f$ が分離的であること $\Leftrightarrow$ $g$ が分離的であって
$\Delta_g : W \to W \times_{X_{/T}} W$ が rig-全射であることとは同値である。
\end{lemma}

\begin{proof}
$f$ が分離的ならば，Lemmas
\ref{restricted-lemma-quasi-separated-gives-quasi-separated} および
\ref{restricted-lemma-separated-mono-open-diagonal-rig-surjective} により，
$g$ は分離的で $\Delta_g$ は rig-全射である。
$g$ は分離的で $\Delta_g$ は rig-全射であると仮定する。
Lemma \ref{restricted-lemma-output-quasi-compact} の証明とまったく同様に，
Morphisms of Spaces, Lemma
\ref{spaces-morphisms-lemma-base-change-separated}
（代数空間の射が分離的であるという性質の基底上の局所性），
Formal Spaces, Lemma
\ref{formal-spaces-lemma-base-change-separated}
（形式代数空間の射が分離的であるという性質は基底変換で保たれる），および
Lemma \ref{restricted-lemma-base-change-rig-surjective}
（rig-全射性は基底変換で保たれる）を用い，アフィンスキームからなる
$X$ の \'etale 被覆の各要素上で確認してよい。
従って $X$ はアフィンであると仮定してよく，そう仮定する。
さらに Lemma \ref{restricted-lemma-output-quasi-separated} により，
$f : X' \to X$ が準分離的であることは既に分かっている。

\medskip\noindent
Cohomology of Spaces, Lemma
\ref{spaces-cohomology-lemma-check-separated-dvr} および Remark
\ref{spaces-cohomology-remark-variant} により，任意の可換図式
$$
\xymatrix{
\Spec(K) \ar[r] \ar[d] & X' \ar[d] \\
\Spec(R) \ar[r]^p \ar@{-->}[ru] & X' \times_X X'
}
$$
が与えられたとき，図式を可換にする点線矢印が存在することを示せば十分である。
ここで $R$ は分数体 $K$ をもつ完備離散付値環である
（これにより付値判定法の一意性部分が得られる）。
% P10-E0753
$h : \Spec(R) \to X$ を，$p$ と射 $Y \times_X Y \to X$ との合成とする。
三つの場合がある。
Case I: $h(\Spec(R)) \subset U$。この場合は
$U' = X' \times_X U \to U$ が同型なので自明である。
Case II: $h$ が $\Spec(R)$ を $T$ に写す場合。この場合は
$g : W \to X_{/T}$ が分離的であるという仮定から従う。
実際，$Z$ を $T$ 上に誘導される被約閉部分空間構造とすると，
$h$ は $Z$ を経由し，
$$
W \times_{X_{/T}} Z = X' \times_X Z \longrightarrow Z
$$
は仮定により分離的である（例えば Formal Spaces, Lemma
\ref{formal-spaces-lemma-separated-local} を参照せよ）。
従って表示した矢印に Cohomology of Spaces, Lemma
\ref{spaces-cohomology-lemma-check-separated-dvr} を適用すると，持上げ性が得られる。
Case III: $h(\Spec(K))$ は $T$ に含まれないが，$h$ が
$\Spec(R)$ の閉点を $T$ に写す場合。この場合，対応する射
$$
p_{/T} :
\text{Spf}(R)
\longrightarrow
(X' \times_X X')_{/T} =
W \times_{X_{/T}} W
$$
は adic 射である（Formal Spaces, Lemma
\ref{formal-spaces-lemma-map-completions-representable} および Definition
\ref{formal-spaces-definition-adic-morphism} による）。
従って $\Delta_g : W \to W \times_{X_{/T}} W$ が rig-全射であるという仮定から，
Lemma \ref{restricted-lemma-monomorphism-rig-surjective} により，$p_{/T}$ を射
$\text{Spf}(R) \to W = X'_{/T}$ に持ち上げられる。
Formal Spaces, Lemma
\ref{formal-spaces-lemma-map-into-algebraic-space} を用いて合成
$\text{Spf}(R) \to X'$ を代数化すると，所望の $p$ の持上げ
$\Spec(R) \to X'$ を得る。
\end{proof}

\begin{lemma}
\label{restricted-lemma-output-proper}
Theorem \ref{restricted-theorem-dilatations-general} の仮定と記法の下で，
$f : X' \to X$ が $g : W \to X_{/T}$ に対応するとする。
このとき $f$ が proper であるための必要十分条件は，$g$ が形式的修正
（Definition \ref{restricted-definition-formal-modification}）であることである。
\end{lemma}

\begin{proof}
$f$ が proper ならば，Lemma
\ref{restricted-lemma-modification-gives-formal-modification} により $g$ は形式的修正である。
$g$ が形式的修正であると仮定する。Lemmas
\ref{restricted-lemma-output-quasi-compact} および \ref{restricted-lemma-output-separated}
により，$f$ は準コンパクトかつ分離的である。

\medskip\noindent
Cohomology of Spaces, Lemma \ref{spaces-cohomology-lemma-check-proper-dvr}
および Remark \ref{spaces-cohomology-remark-variant} により，
任意の可換図式
$$
\xymatrix{
\Spec(K) \ar[r] \ar[d] & X' \ar[d]^f \\
\Spec(R) \ar[r]^p \ar@{-->}[ru] & X
}
$$
が与えられたとき，図式を可換にする点線矢印が存在することを示せば十分である。
ここで $R$ は分数体 $K$ をもつ完備離散付値環である。
三つの場合がある。
Case I: $p(\Spec(R)) \subset U$。この場合は $U' \to U$ が同型なので自明である。
Case II: $p$ が $\Spec(R)$ を $T$ に写す場合。この場合は
$g : W \to X_{/T}$ が proper であるという仮定から従う。
実際，$Z$ を $T$ 上に誘導される被約閉部分空間構造とすると，
$p$ は $Z$ を経由し，
$$
W \times_{X_{/T}} Z = X' \times_X Z \longrightarrow Z
$$
は仮定により proper である。従って表示した矢印に
Cohomology of Spaces, Lemma
\ref{spaces-cohomology-lemma-check-proper-dvr} を適用すると持上げ性が得られる。
Case III: $p(\Spec(K))$ は $T$ に含まれないが，$p$ が
$\Spec(R)$ の閉点を $T$ に写す場合。この場合，対応する射
$$
p_{/T} : \text{Spf}(R) \longrightarrow X'_{/T} = W
$$
は adic 射である（Formal Spaces, Lemma
\ref{formal-spaces-lemma-map-completions-representable} および Definition
\ref{formal-spaces-definition-adic-morphism} による）。
% P10-E0754
従って $g : W \to X_{/T}$ が rig-全射であるという仮定から，
% P10-E0755
$g_{/T}$ を射 $\text{Spf}(R') \to W = X'_{/T}$ に持ち上げられる。
ここで $R \subset R'$ は完備離散付値環のある拡大である。
% P10-E0756
Formal Spaces, Lemma
\ref{formal-spaces-lemma-map-into-algebraic-space} を用いて合成
$\text{Spf}(R') \to X'$ を代数化すると，所望の $p$ の持上げ
$\Spec(R') \to X'$ を得る。
\end{proof}

\begin{lemma}
\label{restricted-lemma-output-etale}
Theorem \ref{restricted-theorem-dilatations-general} の仮定と記法の下で，
$f : X' \to X$ が $g : W \to X_{/T}$ に対応するとする。
このとき $f$ が \'etale であるための必要十分条件は，$g$ が \'etale であることである。
\end{lemma}

\begin{proof}
$f$ が \'etale ならば，Lemma \ref{restricted-lemma-map-completions-etale}
により $g$ も \'etale である。逆に $g$ が \'etale であると仮定する。
$f$ は $U$ 上の同型なので，$f$ は $U$ 上 \'etale である。
従って $T$ の上にある $X'$ の任意の点で $f$ が \'etale であることを
示せば十分である。$Z \subset X$ を，基礎位相空間が
$|Z| = T \subset |X|$ である被約閉部分空間とする。Properties of Spaces,
Definition \ref{spaces-properties-definition-reduced-induced-space} を参照せよ。
$Z_n \subset X$ を第 $n$ 無限小近傍とすれば，
$X_{/T} = \colim Z_n$ である。$X'_{/T} = W \to X_{/T}$ なので，
$f^{-1}(Z_n) = X' \times_X Z_n \to Z_n$ は，$g$ が \'etale であるという
仮定により \'etale である。More on Morphisms of Spaces, Lemma
\ref{spaces-more-morphisms-lemma-check-smoothness-on-infinitesimal-nbhds}
により，$f$ は $T$ の上にある点で \'etale である。
\end{proof}

\section{dilatation に関する Artin の定理}
\label{restricted-section-dilatations}

\noindent
この節では，\cite{ArtinII} の対応する主張との類似性を強調するため，
形式代数空間に別のフォントを用いる。まず本章の第一の主定理を述べる。

\begin{theorem}
\label{restricted-theorem-dilatations}
\begin{reference}
\cite[Theorem 3.2]{ArtinII}
\end{reference}
$S$ をスキームとする。$X$ を $S$ 上の局所 Noether 代数空間とする。
$T \subset |X|$ を閉部分集合とする。$\mathfrak X = X_{/T}$ を
$X$ の $T$ に沿う形式完備化とする。
$$
\mathfrak f : \mathfrak X' \to \mathfrak X
$$
を形式的修正（Definition \ref{restricted-definition-formal-modification}）とする。
このとき，$X$ における $T$ の補集合上で同型であり，その完備化
$f_{/T}$ が $\mathfrak f$ を再現する一意な proper 射
$f : X' \to X$ が存在する。
\end{theorem}

\begin{proof}
Theorem \ref{restricted-theorem-dilatations-general} と Lemma
\ref{restricted-lemma-output-proper} から従う。
\end{proof}

\noindent
Section \ref{restricted-section-formal-modifications} で約束した形式的修正の特徴付けを述べる。

\begin{lemma}
\label{restricted-lemma-formal-modifications-locally-algebraic}
$S$ をスキームとする。$\mathfrak X' \to \mathfrak X$ を，$S$ 上の
局所 Noether 形式代数空間の形式的修正
（Definition \ref{restricted-definition-formal-modification}）とする。次が与えられたとする。
\begin{enumerate}
\item 任意の adic Noether 位相環 $A$，
\item 任意の adic 射 $\text{Spf}(A) \longrightarrow \mathfrak X$。
\end{enumerate}
このとき代数空間の proper 射 $X \to \Spec(A)$ と，同型
$$
\text{Spf}(A) \times_{\mathfrak X} \mathfrak X'
\longrightarrow
X_{/Z}
$$
が存在する。この同型は $\text{Spf}(A)$ 上で，
% P10-E0758
$\mathfrak X$ の基底変換を，$A$ 上の $X$ の「閉ファイバー」
$Z = X \times_{\Spec(A)} \text{Spf}(A)_{red}$ に沿う $X$ の形式完備化と同一視する。
\end{lemma}

\begin{proof}
射 $\text{Spf}(A) \times_{\mathfrak X} \mathfrak X'
\to \text{Spf}(A)$ は Lemma
\ref{restricted-lemma-base-change-formal-modification} により形式的修正である。
従って Theorem \ref{restricted-theorem-dilatations} から結論が従う。
\end{proof}

\section{修正への応用}
\label{restricted-section-modifications}

\noindent
$A$ を Noether 環，$I \subset A$ をイデアルとする。
$X = \Spec(A)$ および $U = X \setminus V(I)$ と置く。
この節では圏
\begin{equation}
\label{restricted-equation-modification}
\left\{
f : X' \longrightarrow X
\quad \middle| \quad
\begin{matrix}
X'\text{ は代数空間である}\\
f\text{ は局所有限型である}\\
f^{-1}(U) \to U\text{ は同型である}
\end{matrix}
\right\}
\end{equation}
を考える。$X'/X$ から $X''/X$ への射とは，$X$ 上の代数空間の射
$X' \to X''$ とする。

\medskip\noindent
$A \to B$ を Noether 環の準同型とし，$J = \sqrt{I B}$ を満たす
イデアル $J \subset B$ を取る。このとき射
$\Spec(B) \to \Spec(A)$ に沿う基底変換は，$A$ に対する圏
(\ref{restricted-equation-modification}) から $B$ に対する圏
(\ref{restricted-equation-modification}) への関手を与える。

\begin{lemma}
\label{restricted-lemma-Noetherian-local-ring}
$A \to B$ を，あるイデアル $I \subset A$ に関する $I$-進完備化上で
同型を誘導する Noether 環の準同型とする
（例えば $B$ が $A$ の $I$-進完備化である場合）。
% P10-E0759
このとき基底変換は，$(A, I)$ に対する圏
(\ref{restricted-equation-modification}) と $(B, IB)$ に対する圏
(\ref{restricted-equation-modification}) の間の圏同値を定める。
\end{lemma}

\begin{proof}
$X = \Spec(A)$, $T = V(I)$ と置く。
$X_1 = \Spec(B)$, $T_1 = V(IB)$ と置く。
Theorem \ref{restricted-theorem-dilatations-general}（実際には Lemma
\ref{restricted-lemma-dilatations-affine} で扱ったアフィンの場合だけでよい）により，
$X$ と $T$ に対する圏 (\ref{restricted-equation-modification}) は，
局所 Noether 形式代数空間の rig-\'etale 射
$W \to X_{/T}$ の圏と同値である。同様に，$X_1$ と $T_1$ に対する圏
(\ref{restricted-equation-modification}) は，局所 Noether 形式代数空間の rig-\'etale 射
% P10-E0760
$W_1 \to X_{1, /T_1}$ の圏と同値である。
$X_{/T} = \text{Spf}(A^\wedge)$ および
$X_{1, /T_1} = \text{Spf}(B^\wedge)$ なので（Formal Spaces, Lemma
\ref{formal-spaces-lemma-affine-formal-completion-types}），
$A^\wedge \to B^\wedge$ が同型であるという仮定により，これらの圏は同値である。
この圏同値が基底変換により与えられることの確認は省略する。
\end{proof}

\begin{lemma}
\label{restricted-lemma-Noetherian-local-ring-properties}
記法と仮定を Lemma \ref{restricted-lemma-Noetherian-local-ring} の通りとする。
$f : X' \to \Spec(A)$ がこの圏同値を介して
$g : Y' \to \Spec(B)$ に対応するとする。このとき $f$ が準コンパクト，
準分離的，分離的，proper，有限，
% P10-E0761
さらにここに追加する，のいずれかであるための必要十分条件は，$g$ がそうであることである。
\end{lemma}

\begin{proof}
% P10-E0762
準コンパクト，準分離的，分離的，proper という主張については，Lemmas
\ref{restricted-lemma-output-quasi-compact}
\ref{restricted-lemma-output-quasi-separated},
\ref{restricted-lemma-output-separated},
\ref{restricted-lemma-output-quasi-separated}, および
\ref{restricted-lemma-output-proper} を用いて対応する性質を形式完備化の性質へ翻訳し，
Lemma \ref{restricted-lemma-Noetherian-local-ring} の証明の議論を用いれば導ける。
しかし次のような fpqc 降下を用いる直接的な議論がある。
まず $A^\wedge \to B^\wedge$ は同型なので，
$A \to A^\wedge$ と $B \to B^\wedge$ について補題を証明することへ還元できる。
次に，先と同じ $U = \Spec(A) \setminus V(I)$ に対し，
$\{U \to \Spec(A), \Spec(A^\wedge) \to \Spec(A)\}$ は fpqc 被覆である。
圏 (\ref{restricted-equation-modification}) の定義により，$f$ の
$U \to \Spec(A)$ による基底変換は $\text{id}_U$ である。
$P$ を，基底上 fpqc 局所的である代数空間の射の性質とする
（Descent on Spaces, Definition
\ref{spaces-descent-definition-property-morphisms-local}）。また $P$ は恒等射について
成り立つとする。このとき $P$ が $f$ について成り立つための必要十分条件は，
$P$ が $g$ について成り立つことである。Descent on Spaces, Lemmas
\ref{spaces-descent-lemma-descending-property-quasi-compact},
\ref{spaces-descent-lemma-descending-property-quasi-separated},
\ref{spaces-descent-lemma-descending-property-separated},
\ref{spaces-descent-lemma-descending-property-proper}, および
\ref{spaces-descent-lemma-descending-property-finite} により，これは
$P$ が準コンパクト，準分離的，分離的，proper，有限である場合に適用できる。
\end{proof}

\begin{lemma}
\label{restricted-lemma-equivalence-to-completion}
$A \to B$ を局所 Noether 環の局所準同型であって，次を満たすものとする。
\begin{enumerate}
\item $A \to B$ は平坦である，
\item $\mathfrak m_B = \mathfrak m_A B$ である，
\item $\kappa(\mathfrak m_A) = \kappa(\mathfrak m_B)$ である。
\end{enumerate}
このとき $(A, \mathfrak m_A)$ に対する圏
(\ref{restricted-equation-modification}) から $(B, \mathfrak m_B)$ に対する圏
(\ref{restricted-equation-modification}) への基底変換関手は圏同値である。
\end{lemma}

\begin{proof}
これらの条件は $A \to B$ が完備化上の同型を誘導することを意味する。
More on Algebra, Lemma \ref{more-algebra-lemma-flat-unramified} を参照せよ。
従ってこの補題は Lemma \ref{restricted-lemma-Noetherian-local-ring} の特別な場合である。
\end{proof}

\begin{lemma}
\label{restricted-lemma-dominate-by-admissible-blowup}
$(A, \mathfrak m, \kappa)$ を Noether 局所環とする。
% P10-E0763
$f : X \to S$ を (\ref{restricted-equation-modification}) の対象であって，
$f$ が proper であるものとする。このとき $X$ を支配する
$U$-許容 blowup $S' \to S$ が存在する。
\end{lemma}

\begin{proof}
More on Morphisms of Spaces, Lemma
\ref{spaces-more-morphisms-lemma-dominate-modification-by-blowup} の特別な場合である。
\end{proof}
